% !TeX program = xelatex
% 数学 LaTeX 排版 Skill v4.0.0；版式保持 v3 金标准. XeLaTeX 编译至少两次. 不分发字体.
% WZ-LOCKED-CLASS-BEGIN
\begin{filecontents*}[overwrite]{elegantbook-original-adapter.cls}
\NeedsTeXFormat{LaTeX2e}
\ProvidesClass{elegantbook-original-adapter}[2026/09/23 v3.0 fixed mathematical book environments]
\LoadClass[10pt,letterpaper,oneside]{book}
\RequirePackage[UTF8,scheme=plain,fontset=fandol]{ctex}
\RequirePackage{fontspec}
\setmainfont{cmunrm.otf}[BoldFont=cmunbx.otf,ItalicFont=cmunti.otf,BoldItalicFont=cmunbi.otf]
\setCJKfamilyfont{sourceheader}[ItalicFont=FandolKai-Regular.otf]{FandolKai-Regular.otf}
\RequirePackage{amsmath,amssymb,mathrsfs}
\RequirePackage{amsthm,etoolbox}
\RequirePackage{xcolor,geometry,setspace,fancyhdr,enumitem,needspace,xurl}
\definecolor{structurecolor}{HTML}{880E4F}
\definecolor{main}{HTML}{9C27B0}
\geometry{paperwidth=612bp,paperheight=792bp,left=20mm,right=20mm,top=95.04bp,bottom=20mm,headheight=12pt,headsep=25pt,footskip=30pt}
\setstretch{1.6}
\setlength{\parindent}{2em}
\setlength{\parskip}{0pt}
\setlist{nosep,leftmargin=2em}
\AtBeginDocument{\setlength{\abovedisplayskip}{4pt plus 1pt minus 1pt}\setlength{\belowdisplayskip}{4pt plus 1pt minus 1pt}\setlength{\abovedisplayshortskip}{2pt}\setlength{\belowdisplayshortskip}{3pt}}
\fancyhf{}
\fancyhead[L]{{\itshape\CJKfamily{sourceheader}\rightmark}}
\fancyhead[R]{{\itshape\CJKfamily{sourceheader}\leftmark}}
\fancyfoot[C]{\thepage}
\renewcommand{\headrulewidth}{0.4pt}
\renewcommand{\headrule}{{\color{structurecolor}\hrule height\headrulewidth width\headwidth\vskip-\headrulewidth}}
\pagestyle{fancy}
\fancypagestyle{cover}{\fancyhf{}\fancyfoot[C]{\thepage}\renewcommand{\headrulewidth}{0pt}\renewcommand{\headrule}{}}
\RequirePackage{hyperref}
\hypersetup{unicode=true,colorlinks=true,linkcolor=structurecolor,urlcolor=structurecolor,citecolor=structurecolor,linktoc=section,bookmarksnumbered=true,bookmarksopen=true,pdfborder={0 0 0},pdfstartview=FitH}
\setcounter{tocdepth}{1}
\renewcommand*\l@section{\@dottedtocline{1}{1.5em}{2.4em}}
\renewcommand{\thesection}{\thechapter.\arabic{section}}
\newcommand{\originalchapter}[1]{\par\addvspace{14pt}\Needspace{5\baselineskip}\refstepcounter{chapter}\setcounter{section}{0}\addcontentsline{toc}{chapter}{\protect\numberline{\thechapter}#1}\markboth{\thechapter\quad #1}{}\begingroup\centering\color{structurecolor}\bfseries\fontsize{14.4}{20}\selectfont\thechapter\quad #1\par\endgroup\nobreak\vspace{8pt}}
\newcommand{\originalsection}[1]{\par\addvspace{9pt}\Needspace{4\baselineskip}\refstepcounter{section}\addcontentsline{toc}{section}{\protect\hspace*{1.5em}\protect\numberline{\protect\textcolor{structurecolor}{\thesection}}#1}\markright{\thesection\quad #1}\noindent{\fontsize{12}{17}\selectfont\color{structurecolor}\thesection\quad}{\bfseries\fontsize{12}{17}\selectfont #1}\par\nobreak\vspace{4pt}}

% Semantic environments are registered once here, not re-designed in manuscripts.
\newtheoremstyle{wzstatement}{6pt}{5pt}
  {\normalfont\kaishu}{0pt}{\normalfont\bfseries\color{structurecolor}}
  {}{.7em}{\thmname{#1}\thmnumber{ #2}\thmnote{\normalfont\ (#3)}}
\newtheoremstyle{wzdefinition}{6pt}{5pt}
  {\normalfont\songti}{0pt}{\normalfont\bfseries\color{structurecolor}}
  {}{.7em}{\thmname{#1}\thmnumber{ #2}\thmnote{\normalfont\ (#3)}}
\newtheoremstyle{wzproblem}{7pt}{5pt}
  {\normalfont\songti}{2em}{\normalfont\bfseries\color{main}}
  {}{.7em}{\thmname{#1}\thmnumber{ #2}}
\newtheoremstyle{wzremark}{4pt}{4pt}
  {\normalfont\songti}{0pt}{\normalfont\bfseries}
  {}{.7em}{\thmname{#1}}
\theoremstyle{wzstatement}
\newtheorem{theorem}{定理}[chapter]
\newtheorem{lemma}[theorem]{引理}
\newtheorem{proposition}[theorem]{命题}
\newtheorem{corollary}[theorem]{推论}
\theoremstyle{wzdefinition}
\newtheorem{definition}[theorem]{定义}
\theoremstyle{wzproblem}
\newtheorem{example}{例题}[chapter]
\newtheorem{exercise}{习题}[chapter]
\theoremstyle{wzremark}
\newtheorem*{remark}{注}
\renewcommand{\proofname}{证明}
% Retain the native amsthm QED stack; only adjust the fixed typography.
\expandafter\patchcmd\csname\string\proof\endcsname{\normalfont \topsep6\p@\@plus6\p@\relax}
  {\normalfont\songti \topsep5\p@\@plus1\p@\relax}
  {}{\ClassError{elegantbook-original-adapter}{Cannot patch proof spacing}{Check the amsthm version.}}
\expandafter\patchcmd\csname\string\proof\endcsname{\itshape}{\normalfont\bfseries}
  {}{\ClassError{elegantbook-original-adapter}{Cannot patch proof heading}{Check the amsthm version.}}
\expandafter\patchcmd\csname\string\proof\endcsname{\ignorespaces}
  {\setlength{\parindent}{2em}\ignorespaces}
  {}{\ClassError{elegantbook-original-adapter}{Cannot patch proof paragraphs}{Check the amsthm version.}}
\AtBeginEnvironment{proof}{\par\Needspace{3\baselineskip}}
\renewcommand{\qedsymbol}{\ensuremath{\square}}
\newenvironment{solution}{\begin{proof}[解答]}{\end{proof}}
\newcommand{\wzcheckchapter}{\ifnum\value{chapter}<1
  \ClassError{elegantbook-original-adapter}{Numbered mathematics before chapter 1}{Move it after the first originalchapter.}\fi}

\newcommand{\wzregisterguard}[1]{\AtBeginEnvironment{#1}{\wzcheckchapter\par\Needspace{3\baselineskip}}}
\forcsvlist{\wzregisterguard}{theorem,lemma,proposition,corollary,definition,example,exercise}
% Front matter and bibliography are not numbered mathematical chapters.
\newcommand{\originalpreface}{\par\addvspace{10pt}\Needspace{5\baselineskip}
  \noindent{\bfseries\fontsize{12}{17}\selectfont 前言}\par\nobreak\vspace{4pt}}
\newcommand{\originalcontents}{\par\addvspace{14pt}
  \begin{center}{\color{structurecolor}\bfseries\fontsize{14.4}{20}\selectfont 目录}\end{center}
  \vspace{-10pt}\@starttoc{toc}}
\renewcommand{\bibname}{参考文献}
\renewenvironment{thebibliography}[1]{
  \par\addvspace{12pt}\Needspace{3\baselineskip}\phantomsection
  \addcontentsline{toc}{chapter}{\bibname}
  \markboth{\bibname}{}
  \noindent{\color{structurecolor}\bfseries\fontsize{14.4}{20}\selectfont\bibname}\par\nobreak\vspace{6pt}
  \list{\@biblabel{\@arabic\c@enumiv}}
    {\settowidth\labelwidth{\@biblabel{#1}}\leftmargin\labelwidth
     \advance\leftmargin\labelsep\usecounter{enumiv}
     \let\p@enumiv\@empty\renewcommand\theenumiv{\@arabic\c@enumiv}
     \setlength{\itemsep}{3pt}\setlength{\parsep}{0pt}}
  \sloppy\clubpenalty4000\widowpenalty4000\sfcode`\.=1000\relax
}{\def\@noitemerr{\@latex@warning{Empty bibliography}}\endlist}

\numberwithin{equation}{chapter}
\renewcommand{\theequation}{\thechapter.\arabic{equation}}
\allowdisplaybreaks[2]
\emergencystretch=2em
\clubpenalty=6000
\widowpenalty=6000
\raggedbottom
\endinput
\end{filecontents*}
% WZ-LOCKED-CLASS-END
\documentclass{elegantbook-original-adapter}
% 本册有两位章、节号；仅调整目录号栏，不修改锁定类定义。
\makeatletter
\renewcommand*\l@section{\@dottedtocline{1}{1.5em}{3.8em}}
\makeatother

\usepackage{tikz,pgfplots,longtable,booktabs}
\usetikzlibrary{arrows.meta,calc,positioning,patterns}
\pgfplotsset{compat=1.18}
\newcommand{\dd}{\,\mathrm d}
\newcommand{\kb}{k_{\mathrm B}}
\newcommand{\avg}[1]{\langle #1\rangle}
\newcommand{\Tr}{\operatorname{Tr}}
\newcommand{\e}{\mathrm e}
\newcommand{\ii}{\mathrm i}
\newcommand{\vct}[1]{\boldsymbol{#1}}
\newcommand{\ket}[1]{\lvert #1\rangle}
\newcommand{\bra}[1]{\langle #1\rvert}
\newcommand{\source}[1]{\par\noindent\textit{来源与编写：#1}\par}
\newcommand{\psource}[1]{\par\noindent\textit{题源：#1。以下解答为编者独立补写。}\par}
\newcommand{\osource}[1]{\par\noindent\textit{题源：#1。官方答案可获取；中文解答由编者独立推导，题内说明与官方解的核对及修正。}\par}
% Keep a chapter heading with its source and opening paragraph.
\pretocmd{\originalchapter}{\Needspace{10\baselineskip}}{}{}
\newcommand{\volpart}[1]{\par\addvspace{16pt}\Needspace{5\baselineskip}\phantomsection\addcontentsline{toc}{part}{#1}\begin{center}{\color{structurecolor}\bfseries\fontsize{16}{22}\selectfont #1}\end{center}}

\hypersetup{pdftitle={高级统计力学与统计场论：MIT OCW 中文重编},pdfauthor={中文重编项目}}
\begin{document}
\frontmatter\thispagestyle{cover}
\begin{center}
{\color{structurecolor}\bfseries\fontsize{17}{23}\selectfont 高级统计力学与统计场论}\par
\vspace{5pt}{\bfseries 粒子动力学、临界现象与重整化群}\par
\vspace{5pt}根据 MIT OpenCourseWare 公开课程中文重编\par
2026年10月12日
\end{center}
\originalpreface
微观动力学保留可逆的粒子轨道，宏观气体却表现黏性、热传导与熵增；临界体系的相关长度又可远大于粒子尺度。高级统计力学的任务是说明这些描述怎样相容，并判断近似在哪个尺度有效。本册读者应已熟悉本科热力学、微正则/正则分布、理想量子气体与基本量子力学。仓库既有《热力学与统计物理》提供这些先修及8.333第15--22讲的部分补充；本卷用必要的符号连接它，重点展开动力学、关联与统计场论，不重复本科专题的全部例题。

课程主线为Mehran Kardar的\href{https://ocw.mit.edu/courses/8-333-statistical-mechanics-i-statistical-mechanics-of-particles-fall-2013/}{8.333，Fall 2013}及\href{https://ocw.mit.edu/courses/8-334-statistical-mechanics-ii-statistical-physics-of-fields-spring-2014/}{8.334，Spring 2014}。两站各公开26份课历对应讲义，官方都说明讲义来自更早版本；若文件未给确切原稿年份，本稿不猜造年份。课历的2013、2014日期、早期讲义版本与往年复习题年代分开记录。每门课有6套公开作业。8.333考试页列3套阶段复习与一套期末复习，8.334列3套阶段复习，都有公开官方解；8.333前三包标Fall2007，期末包含1998、1999、2000、2003--2006历史卷；8.334三包标Spring2014复习题，不能把复习题包称为该学期实际测试卷。8.333成绩为作业30\%、三测验45\%、期末25\%；8.334为作业30\%、三测验45\%、最终项目25\%。没有取得2013、2014当期实际考试卷及学生项目，均不补造。

第一编连接粒子、系综与流体，第二编系统建立场、相变、尺度与RG，第三编收录实际公开问题及中文解答。无公开答案的作业标为编者独立补写；复习包注明可获取官方解；中文解由编者独立推导，并在题内单列官方对照和确定可核验的修正。旧卷已写的8.333作业5仍核对原题并给固定前置定位，重复问题采用明确回指，而不作为新增解答计数。视频各26讲保留官方链接，未宣称逐分钟转录全部视频。

来源署名Mehran Kardar、MIT OpenCourseWare及文件所列作者。本项目作中文翻译、重编、推导补充与自绘矢量图，基于OCW部分采用\href{https://creativecommons.org/licenses/by-nc-sa/4.0/}{CC BY-NC-SA 4.0}，遵守署名、非商业及相同方式共享，并保留逐件例外；本项目不代表MIT认可。商业参考书只留出版信息与原题号，不打包整书或以答案反造书本题面。课程公开目录、实际文件与完成范围分别记录在学科来源清单与卷末。

\clearpage\originalcontents\clearpage\mainmatter
\volpart{第一编　粒子动力学、系综与量子气体（8.333）}
\originalchapter{热力学稳定性与概率结构}
\label{ch:thermo}
\source{8.333早期讲义对应L1--L6；本科定律仅作必要连接，重点为稳定性、累积量与大数尺度。}
\originalsection{势函数与稳定性}
平衡态熵$S(U,V,N)$满足$\dd S=T^{-1}\dd U+(p/T)\dd V-(\mu/T)\dd N$。两个可交换能量的子系统在总$U$固定时使$S_1(U_1)+S_2(U-U_1)$极大，一阶条件$T_1=T_2$，二阶条件要求相应熵曲率非正。用$\partial_U^2S=-1/(T^2C_V)$，局部稳定且$C_V$有限时$C_V>0$。稳定性来自允许的涨落方向；若约束改变，要在实际允许子空间判断Hessian，不是随便对某个变量求二阶导数。

Legendre变换把不便固定的广延量换成热库的强度量。$F=U-TS$有$\dd F=-S\dd T-p\dd V+\mu\dd N$；$G=F+pV$有$\dd G=-S\dd T+V\dd p+\mu\dd N$。混合偏导相等给Maxwell关系，例如$\partial_VS|_T=\partial_Tp|_V$。它们要求处于可微的同一热力学相，跨一阶跃迁的尖点不能无解释地交换偏导。

设$\kappa_T=-V^{-1}\partial_pV|_T>0$、$\alpha=V^{-1}\partial_TV|_p$。以$\dd S=(C_V/T)\dd T+(\partial_Tp)_V\dd V$，在定压路径$\dd V=\alpha V\dd T$，再由链式关系$(\partial_Tp)_V=\alpha/\kappa_T$，得
\begin{equation}
C_p-C_V=\frac{TV\alpha^2}{\kappa_T}\ge0.
\end{equation}
这不仅给不等式，还表明两种热容差测量热膨胀与可压缩性的耦合。相稳定的局部条件不能替代全局自由能最低条件；亚稳态可有正热容、正压缩率却不是全局平衡。

第三定律给完美晶体的非简并平衡基态$S\to0$，存在有限简并时可剩$\kb\log g_0$。无序冻结的残余熵和未达平衡的玻璃状态需要另作物理判断，不能靠选择熵零点去掉实际状态计数。
\originalsection{联合分布、累积量与相关}
随机变量$X$的特征函数$\varphi(t)=\avg{\e^{\ii tX}}$总存在，矩母函数$M(t)=\avg{\e^{tX}}$可能只在有限区间存在。定义$K(t)=\log M(t)$，累积量$\kappa_n=K^{(n)}(0)$。前几项$\kappa_1=\avg X$、$\kappa_2=\avg{X^2}-\avg X^2$、$\kappa_3=\avg{(X-\avg X)^3}$、$\kappa_4=\avg{(X-\avg X)^4}-3\kappa_2^2$。累积量把不可约相关与由低阶矩乘积产生的部分分开。

独立和$S_N=\sum_iX_i$有$M_{S_N}=\prod M_i$，所以$K_{S_N}=\sum K_i$，所有累积量相加。对联合变量，$K(t_1,\ldots,t_n)$的混合导数给连接相关。例如$\avg{XYZ}_c=\avg{XYZ}-\avg X\avg{YZ}-\avg Y\avg{XZ}-\avg Z\avg{XY}+2\avg X\avg Y\avg Z$。独立两组间的混合累积量为0；反过来若只知道二阶相关为0，不能断言独立，除非已知联合高斯等附加条件。

对Gaussian变量$K=\mu t+\sigma^2t^2/2$，高阶累积量全0。正则配分函数的参数导数也生成累积量：$\partial_\beta\log Z=-\avg H$、$\partial_\beta^2\log Z=\avg{H^2}_c$，所以$C_V=\kb\beta^2\operatorname{Var}H$。统计涨落不仅“很小”，还定量控制响应。
\originalsection{中心极限定理与大偏差}
对独立同分布且有限方差的变量，$Y_N=(S_N-N\mu)/(\sigma\sqrt N)$的特征累积量展开为$\log\varphi_{Y_N}(t)=-t^2/2+O(N^{-1/2})$，各$n>2$阶缩为$N^{1-n/2}$（若所需高阶矩存在）。更一般仅有限二阶矩也可用小$t$的$o(t^2)$余项建立Gaussian极限。重尾无有限方差时此结论失效，稳定分布可替代Gaussian。

中心极限定理只描述均值附近$O(\sqrt N)$涨落。对$S_N/N\simeq x$、$x-\mu=O(1)$，若矩母函数在0附近存在，Laplace鞍点给$P\sim\e^{-NI(x)}$，$I(x)=\sup_t[tx-K(t)]$。例如Bernoulli成功率$p$，$K=\log(1-p+p\e^t)$，最优$\e^t=x(1-p)/[p(1-x)]$，故
\begin{equation}
I(x)=x\log\frac xp+(1-x)\log\frac{1-x}{1-p}.
\end{equation}
近$x=p$展开为$(x-p)^2/[2p(1-p)]$，接回中心极限；靠近0或1的罕见事件不能直接用同一二次抛物线估计。
\originalsection{信息熵与估计}
离散熵$H=-\sum p_i\log p_i$是无量纲信息量，物理熵为$\kb H$加与相空间基准有关的常数。连续微分熵依变量的单位与坐标变换，不能单独看它的绝对值评价热力学熵。相对熵$D(p\Vert q)=\sum p_i\log(p_i/q_i)\ge0$由$\log x\le x-1$证明，等号仅$p=q$。

最大熵在归一与$\avg E=U$约束下变分，$p_i\propto\e^{-\beta E_i}$，$\beta$是能量约束的乘子；其物理温度意义还需与热库平衡条件相接。Fisher信息$I(\theta)=\avg{(\partial_\theta\log p)^2}$量化参数变化的可分辨度；在可微、积分与导数可交换、支撑不随参数突变等正则条件下，无偏估计的方差$\ge1/[NI(\theta)]$。它不是不加条件的所有统计估计精度定理。

\originalchapter{Liouville方程与BBGKY层级}
\label{ch:bbgky}
\source{8.333 L7--L8。}
\originalsection{相空间密度与可逆演化}
对$N$个经典粒子，$\Gamma=(r_1,p_1,\ldots,r_N,p_N)$是$6N$维相空间点，微观态是一个点而统计态是归一概率密度$\rho_N(\Gamma,t)$，$\int\rho_N\dd\Gamma=1$。Hamilton方程$\dot r_i=\partial_{p_i}H,\dot p_i=-\partial_{r_i}H$的相空间散度为$\sum_i(\partial_{r_i}\partial_{p_i}H-\partial_{p_i}\partial_{r_i}H)=0$。连续性给
\begin{equation}
\partial_t\rho_N+\{\rho_N,H\}=0,\qquad\{A,B\}=\sum_i(\nabla_{r_i}A\cdot\nabla_{p_i}B-\nabla_{p_i}A\cdot\nabla_{r_i}B).
\end{equation}
沿Hamilton轨道$\dd\rho_N/\dd t=0$，并保持相空间体积，这就是Liouville定理。它不保证体系遍历，也不保证任意初态最终成为正则分布。细粒度Gibbs熵$S_G=-\kb\int\rho_N\log\rho_N\dd\Gamma$随可逆演化守恒；宏观熵增必须解释丢弃信息的方式。

若$H$不显含时，任意$\rho_N=F(H)$均平稳，选择何种$F$依准备与宏观约束。对自旋、Berry波包等非正则坐标，应使用相应辛测度；直接用$\dd r\dd p$未必正确。
\originalsection{约化分布的定义}
为了研究少数粒子，定义带粒子数因子的$s$粒子分布
\begin{equation}
f_s(1,\ldots,s)=\frac{N!}{(N-s)!}\int\rho_N(1,\ldots,N)\dd(s+1)\cdots\dd N,
\end{equation}
其中$\dd i=\dd^dr_i\dd^dp_i$。于是$\int f_1\dd1=N$，$\int f_2\dd2=(N-1)f_1$。另一些文献使用归一至1的约化分布，层级方程前的$N-s$因子不同，应先看定义再比较式子。

设$H=\sum_ip_i^2/(2m)+\sum_iV_{\rm ext}(r_i)+\sum_{i<j}v(r_i-r_j)$，力$F_{ij}=-\nabla_{r_i}v(r_i-r_j)$。把Liouville方程对未保留变量积分，表面项在周期边界或足够快衰减下消失，得到
\begin{align}
\bigg[\partial_t+\sum_{i=1}^s\frac{p_i}{m}\cdot\nabla_{r_i}+\sum_{i=1}^sF_i^{\rm ext}\cdot\nabla_{p_i}+\sum_{i<j\le s}F_{ij}\cdot(\nabla_{p_i}-\nabla_{p_j})\bigg]f_s
=-\sum_{i=1}^s\int F_{i,s+1}\cdot\nabla_{p_i}f_{s+1}\dd(s+1).
\end{align}
左侧只含保留粒子的内作用，右侧通过$f_{s+1}$联系环境。这是BBGKY层级：一粒子分布方程不封闭，因为碰撞概率依赖两粒子相关。
\originalsection{闭合近似及其限制}
若长程弱作用且每粒子同时受许多粒子影响，可近似$f_2(1,2)\simeq f_1(1)f_1(2)$（有限$N$归一差另计），右侧成为平均场力，得到Vlasov方程。它仍为无碰撞、可逆的相空间输运，不能自动产生Boltzmann碰撞熵。

稀薄短程气体则两体碰撞重要。把碰撞前两粒子视为近独立，$f_2^{\rm in}\simeq f_1f_1$，但碰撞后已经关联。这个“分子混沌”假设有时间方向：若把所有速度反转后的精确相关态作为初态，不能再假定同样的入射独立性。低密度、碰撞时间短于平均自由时间、空间变化尺度大于作用范围是闭合的物理条件。三体同时碰撞和致密相关都可使其失效。

关联被层级输运到更高$f_s$，只观察$f_1$等价于忽略这些信息。细粒度熵守恒与一粒子粗粒度熵增因而没有逻辑矛盾；矛盾只会出现在声称无需任何闭合或粗粒化也能从时间可逆方程严格推出单调宏观熵时。

\originalchapter{Boltzmann碰撞与H定理}
\label{ch:boltzmann}
\source{8.333 L8--L10。}
\originalsection{两体碰撞的增益与损失}
令$f(r,p,t)$归一为粒子数。稀薄经典气的局域二体碰撞项为
\begin{equation}
C[f](p_1)=\int\dd p_2\dd\Omega\,|v_1-v_2|\frac{\dd\sigma}{\dd\Omega}\,[f(p_1')f(p_2')-f(p_1)f(p_2)].
\end{equation}
$\sigma$是两体微分截面，$p_1',p_2'$与$p_1,p_2$由动量、能量守恒及散射角联系。第一项为碰撞入当前动量的增益，第二项为离开的损失。微观可逆性使正反散射的权重及相空间测度相同。Boltzmann方程为$\partial_tf+v\cdot\nabla_rf+F\cdot\nabla_pf=C[f]$。

量子气在末态还需乘$1\pm f$（占据数按单态归一），分别表示玻色增强和Pauli阻塞。不能直接把经典$\log f$的H函数拿来声称量子H定理；相应量子熵要包含$(1\pm f)\log(1\pm f)$。
\originalsection{碰撞不变量}
对任意函数$\psi(p)$，交换粒子1、2并将正反过程成对，
\begin{equation}
\int\psi_1C[f]\dd p_1=\frac14\int\dd\mathcal C\,(f_1'f_2'-f_1f_2)(\psi_1+\psi_2-\psi_1'-\psi_2'),
\end{equation}
$\dd\mathcal C$含非负散射测度。对$\psi=1,p_i,p^2/(2m)$，最后括号为0，碰撞分别保持粒子数、总动量和动能。这给流体守恒方程的微观来源。一般外场在输运左侧改变动量或能量，不是碰撞守恒失败。
\originalsection{H定理与平衡态}
定义$H[f]=\int f\log f\dd r\dd p$。无净边界熵流时输运项对$H$积分为0。利用粒子数守恒去掉$+1$导数项，碰撞贡献化为
\begin{equation}
\frac{\dd H}{\dd t}=-\frac14\int\dd r\dd\mathcal C\,(f_1'f_2'-f_1f_2)\log\frac{f_1'f_2'}{f_1f_2}\le0.
\end{equation}
因为对正$x,y$有$(x-y)\log(x/y)\ge0$，每个碰撞通道非正。物理粗粒度熵为$-\kb H$加计数常数，所以增加。该不等式是在Boltzmann闭合之下的证明，不是对任意精确$\rho_N$的声明。

平衡要求每条可达碰撞$f_1f_2=f_1'f_2'$，故$\log f$为碰撞不变量。对能连接整个动量空间的普通弹性散射，$\log f=a+\vct b\cdot p-cp^2$，完成平方得到
\begin{equation}
f_{\mathrm{le}}(r,p)=n(r)[2\pi m\kb T(r)]^{-d/2}\exp[-|p-mu(r)|^2/(2m\kb T(r))].
\end{equation}
这是局域Maxwell分布，碰撞项局域为0，但若$n,T,u$随位置变化，输运项一般不为0；全局平衡还需满足外势与边界条件。若碰撞有额外守恒量或不能混合所有动量，其零空间可能更大。
\originalsection{可逆性问题的具体边界}
Loschmidt反转将所有粒子速度倒转，同时将出射碰撞相关转成未来的入射相关，此特殊准备不满足分子混沌。Poincaré复现适于有限、封闭、有限相空间测度体系，指出无穷时间永久单调的精确熵增不能普遍成立。两者要求辨别热力学极限、可观测时间与统计初态，不能用“粒子多”一句话把数学反例抹去。

Boltzmann方程因忽略细相关而获得有效不可逆性。当观测尺度远大于碰撞时间，闭合可准确预测输运，即使全微观态的相关信息仍在增长。关于该近似的严格适用时间存在更深的数学问题，本稿未证明任意长时间、任意密度的Boltzmann极限。

\originalchapter{流体方程与输运系数}
\label{ch:hydro}
\source{8.333 L10--L11：守恒律、零阶与一阶流体。}
\originalsection{由动量矩得到守恒方程}
定义$n=\int f\dd p$、$nu=\int(p/m)f\dd p$，局部相对速度$c=p/m-u$。将Boltzmann方程依次乘$1,p_i,p^2/(2m)$并积分，用碰撞不变量，得到
\begin{align}
&\partial_tn+\nabla\cdot(nu)=0,\\
&mn(\partial_t+u\cdot\nabla)u_i=-\partial_jP_{ij}+nF_i,\\
&\partial_te+\nabla\cdot[(e+P)u+q+\pi\cdot u]=nu\cdot F.
\end{align}
这里$P_{ij}=m\int c_ic_jf\dd p$，分成标量压强$P$与偏应力$\pi_{ij}$；$q_i=\frac m2\int c^2c_if\dd p$是相对局域流速的热流，$e$含流动动能和内能。不能把实验总能流全部称为“热流”，对流项应分开。

局域Maxwell分布给$P_{ij}=n\kb T\delta_{ij}$、$q=0$，能量闭合为理想Euler方程，$D_tT=-(2/d)T\nabla\cdot u$。局域平衡只给无耗散零阶方程，黏性和热导来自$f$偏离此分布的一阶梯度。
\originalsection{线性化碰撞算符与零模}
写$f=f_{\mathrm{le}}(1+\phi)$，线性化$C[f]=-f_{\mathrm{le}}\mathcal L\phi$。加权内积$(a,b)=\int f_{\mathrm{le}}ab\dd p$下，微观可逆性使$\mathcal L$半正，其二次型为碰撞中$[\phi_1+\phi_2-\phi_1'-\phi_2']^2$的非负积分。$1,p_i,p^2$是零模，对应守恒量。求$\phi$时需选择与这些零模正交，保证$n,u,T$由$f_{\mathrm{le}}$本身定义；否则把同一个局域平衡改变算两次。

一阶输运方程为$\mathcal L\phi=-(\partial_t+v\cdot\nabla+F\cdot\nabla_p)\log f_{\mathrm{le}}$，右侧用零阶Euler关系消掉时间导数，剩温度梯度与速度剪切源。对无内自由度理想气，标量膨胀源在该近似下与守恒量结合而不产生独立体黏性；有内弛豫的分子气体则可有体黏性。
\originalsection{弛豫时间中的显式计算}
用满足守恒匹配的简化BGK碰撞$C=-(f-f_{\mathrm{le}})/\tau$，则三维一阶偏离可写
\begin{equation}
\phi=-\tau\left[\frac{m}{\kb T}c_ic_j\mathcal E_{ij}+\frac{c\cdot\nabla T}{T}\left(\frac{mc^2}{2\kb T}-\frac52\right)\right],\quad\mathcal E_{ij}=\frac{\partial_iu_j+\partial_ju_i}2-\frac13\delta_{ij}\nabla\cdot u.
\end{equation}
$5/2$使热流源与动量零模正交，而非随意选常数。所需的Gaussian速度矩为
\begin{equation}
\avg{c_ic_j}=\frac{\kb T}{m}\delta_{ij},\qquad
\avg{c_ic_jc_lc_m}=\left(\frac{\kb T}{m}\right)^2
(\delta_{ij}\delta_{lm}+\delta_{il}\delta_{jm}+\delta_{im}\delta_{jl}).
\end{equation}
代入应力和热流定义可得
\begin{equation}
\pi_{ij}=-2\eta\mathcal E_{ij},\qquad\eta=n\kb T\tau,\qquad q=-\kappa\nabla T,\qquad\kappa=\frac52\frac{n\kb^2T\tau}{m}.
\end{equation}
热导需六阶Gaussian矩；展开$c^2c_ic_j[mc^2/(2\kb T)-5/2]$并用球对称径向矩，可核对系数$5/2$。BGK预测Prandtl数1，而真实单原子稀薄气接近$2/3$，所以一个共同$\tau$只提供结构与量级，不是精确碰撞截面解。

动量扩散$\eta/(mn)$与热扩散$\kappa/(nc_p)$通常近$\bar v\ell$，$\ell=\bar v\tau$。稀薄气碰撞率$\tau^{-1}\sim n\sigma\bar v$，$\eta\sim m\bar v/\sigma$近乎与$n$无关：载流粒子更多但平均自由程更短，这解释一个看似反常的密度取消。
\originalsection{熵产生与声模}
局域熵平衡可整理为$\partial_ts+\nabla\cdot(su+q/T)=\kappa(\nabla T)^2/T^2+2\eta\mathcal E_{ij}\mathcal E_{ij}/T+\zeta(\nabla\cdot u)^2/T\ge0$。正黏性与热导源于线性化碰撞算符半正，任意设负系数会违反近平衡稳定性。

在线性均匀背景，横向速度模$\omega=-\ii\eta q^2/(mn)$、热扩散模$\omega=-\ii D_Tq^2$，纵向密度与温度耦合产生$\omega=\pm c_sq-\ii\Gamma_sq^2$，$c_s^2=(\partial p/\partial\rho)_s$。它们是流体尺度的集体模式，成立条件$q\ell\ll1,\omega\tau\ll1$；在无碰撞窗口需回到动力学方程，不能把同一声速外推任意$q$。

\originalchapter{系综、密度矩阵与量子计数}
\label{ch:ensemble}
\source{8.333 L12--L14、L19--L23；本科气体公式仅保留高级推导所需连接。}
\originalsection{从热库到正则分布}
封闭总体系在能量壳等概率，子系统能量$E$的权重是热库态数$\Omega_B(U-E)$。展开$\log\Omega_B(U-E)=\log\Omega_B(U)-\beta E+O(E^2/C_BT^2)$，足够大的热库给$p(E)\propto\e^{-\beta E}$。下一阶定量说明有限热库误差，若子系统可交换的能量与热库热容同量级，正则近似不再充分。

经典单原子理想气$Z_N=V^N/(N!\lambda_T^{dN})$，$\lambda_T=\sqrt{2\pi\hbar^2/(m\kb T)}$。$N!$来自相同粒子不重复计数；去掉它会使熵混合悖论与不正确的非广延项出现。对不同种粒子各自除$N_\alpha!$，混合熵仍真实存在。微正则、正则、巨正则在热力学极限的局域观测可等价，需要稳定短程作用与适当可微大偏差结构；长程引力等系统不能无条件套这项等价。
\originalsection{能量与粒子数的涨落矩阵}
巨正则$\Xi=\Tr\exp[-\beta(H-\mu N)]$，以$u=-\beta,v=\beta\mu$作独立参数，则$\partial_u\log\Xi=\avg H$、$\partial_v\log\Xi=\avg N$，Hessian为$(H,N)$协方差矩阵，半正。固定$T$时$\partial_\mu\avg N=\beta\operatorname{Var}N$，结合$n^2\kappa_T=\partial n/\partial\mu$得$\operatorname{Var}N=V\kb T n^2\kappa_T$。交叉协方差解释热膨胀及不同响应之间的联系，不能一边固定$N$一边保留巨正则$N$涨落公式。
\originalsection{量子熵与变分原理}
量子统计态是正算符$\rho$、$\Tr\rho=1$，熵$S=-\kb\Tr\rho\log\rho$。canonical态$\rho_\beta=\e^{-\beta H}/Z$同时包含经典概率与本征态相干结构的可能性；热平衡在能量基对角不意味着所有非平衡态也对角。

相对熵$D(\rho\Vert\sigma)=\Tr\rho(\log\rho-\log\sigma)\ge0$可在有限维直接证明。以$\rho$本征基$\ket i$及$\sigma$本征值$s_j$写第二项，$\avg{i|\log\sigma|i}=\sum_j|\avg{i|j}|^2\log s_j\le\log\avg{i|\sigma|i}$由凹性得到；于是$D\ge\sum_ir_i\log[r_i/q_i]\ge0$，$q_i=\avg{i|\sigma|i}$。支撑不包含时相对熵为$+\infty$。

取$\sigma=\rho_\beta$，得$\Tr\rho H-TS(\rho)\ge F$，所以平衡态极小自由能泛函。再取参考$H_0$的热态$\rho_0$，得到Bogoliubov变分不等式
\begin{equation}
F\le F_0+\avg{H-H_0}_0.
\end{equation}
它对量子非对易$H,H_0$也成立；不需要错误地把$\e^{-\beta(H_0+V)}$拆成两个可交换指数。无限维需相应配分函数为trace class及期望有限。
\originalsection{分子内部自由度与量子极限}
多原子气体的平移、转动、振动在弱耦合近似下配分函数相乘。刚性线性转子$E_l=\hbar^2l(l+1)/(2I)$、简并$2l+1$，高温转动配分近$2I\kb T/\hbar^2$，低温转动态冻结；同核分子的核自旋交换可限制偶奇$l$，不能把异核结果直接套去。谐振模$Z=\e^{-\beta\hbar\omega/2}/(1-\e^{-\beta\hbar\omega})$，高温恢复每模$\kb$热容，低温指数冻结。

固体声子与黑体辐射的低温态密度积分详见既有卷；这里强调计数约束：光子与平衡声子数不守恒，$\mu=0$；原子玻色气$N$守恒，$\mu$需由数方程求。把“玻色子”当作自动$\mu=0$会使BEC数方程失去意义。
\originalsection{相同粒子的排列与交换}
两个相同粒子的Hilbert空间先取张量积，再投到对称或反对称子空间。$N$粒子的投影与配分函数为
\begin{equation}
P_\pm=\frac1{N!}\sum_{\pi\in S_N}(\pm1)^\pi U_\pi,
\qquad Z_N^\pm=\Tr\!\left[P_\pm\e^{-\beta\sum_i h_i}\right].
\end{equation}
排列分解为长度$l$循环，每个循环贡献单体迹$Z_1(l\beta)$；循环之间的组合计数给
\begin{equation}
\log\Xi_\pm=\sum_{l=1}^\infty\frac{(\pm1)^{l-1}z^l}{l}Z_1(l\beta).
\end{equation}
对自由气$Z_1(l\beta)=V/(\lambda_T^dl^{d/2})$，交换的量子virial修正由$l=2$开始。这里不需诉诸额外相互作用就得到玻色趋向聚集、费米趋向排斥的有效第二virial符号。高温$z\ll1$或$n\lambda_T^d\ll1$只保留$l=1$，恢复经典不可分辨计数。

\originalchapter{连接图、virial与液气相变}
\label{ch:clusters}
\source{8.333 L15--L18及L19的课历回顾；旧卷相关公式在此用于连接图与临界问题。}
\originalsection{Mayer函数与连接展开}
两体势$v_{ij}$使$\e^{-\beta\sum v_{ij}}=\prod_{i<j}(1+f_{ij})$，$f_{ij}=\e^{-\beta v_{ij}}-1$。每项可画为粒子顶点与所选$f$边的图。构型积分中不相连的图分解成乘积，而$\log\Xi$只保留连接图，这是累积量与连接相关的同一组合结构。

写$\beta p=\sum_{l\ge1}b_lz^l$、$n=\sum l b_lz^l$，$z$无量纲。平移不变、短程势下
\begin{align}
b_1&=\lambda_T^{-d},\quad b_2=\frac1{2\lambda_T^{2d}}\int f(r)\dd^dr,\\
b_3&=\frac1{6\lambda_T^{3d}}\iint[ f_{12}f_{13}+f_{12}f_{23}+f_{13}f_{23}+f_{12}f_{13}f_{23}]\dd^dr_2\dd^dr_3,
\end{align}
固定$r_1=0$后总体积已除去。三个树图与一个三角图皆为连接，不能只留三角便称为$b_3$。
\begin{center}\begin{tikzpicture}[scale=.8]
\foreach \x in {0,3.4,6.8}{\fill(\x,0)circle(2pt);\fill(\x+1,1)circle(2pt);\fill(\x+2,0)circle(2pt);\draw(\x,0)--(\x+1,1)--(\x+2,0);}
\draw(6.8,0)--(8.8,0);\node[below]at(1,-.3){连接树};\node[below]at(4.4,-.3){另一树（标签不同）};\node[below]at(7.8,-.3){不可约三角};
\end{tikzpicture}\end{center}
图中的第二树仅示形状；实际标签置换在代数中给三份，不能把形状数当带标签图数。
\originalsection{消去逸度与virial系数}
由数密度级数反解$z=n/b_1-2b_2n^2/b_1^3+[8b_2^2/b_1^5-3b_3/b_1^4]n^3+\cdots$，回代压力得到
\begin{equation}
\beta p=n+B_2n^2+B_3n^3+\cdots,\quad B_2=-\frac{b_2}{b_1^2},\quad B_3=\frac{4b_2^2}{b_1^4}-\frac{2b_3}{b_1^3}.
\end{equation}
因此$B_2=-\frac12\int f$，$B_3=-\frac13\iint f_{12}f_{23}f_{31}$，树图在由逸度换成密度时取消。这个取消解释为何virial展开只含相应二连接图，而不是随意删掉的经验规则。

三维硬球直径$\sigma$有$f=-1$于$r<\sigma$、其余0，$B_2=2\pi\sigma^3/3$，$B_3=5\pi^2\sigma^6/18$。后者为三个两两距离小于$\sigma$的几何体积，需做重叠球积分；可取$r_{12}=r$，两个半径$\sigma$球交体积$\pi(4\sigma+r)(2\sigma-r)^2/12$，再$\int_0^\sigma4\pi r^2\dd r$并除3，得到所给值。这样保留第三virial的可复算主计算。
\originalsection{van der Waals与共存}
排斥硬芯给$b=2\pi\sigma^3/3$，弱吸引外尾把$f\simeq-\beta v$，$B_2=b-a/(\kb T)$，$a=-2\pi\int_\sigma^\infty r^2v(r)\dd r>0$。把排斥近似为有效自由体积、吸引近似为平均密度作用，得到$p=n\kb T/(1-bn)-an^2$。这个重组比仅截断$B_2$更能表现密度非线性，却不是精确簇展开到所有阶。

临界点满足$\partial_np=\partial_n^2p=0$，求得$n_c=1/(3b)$、$\kb T_c=8a/(27b)$、$p_c=a/(27b^2)$。低于$T_c$曲线有$\partial_np<0$的不稳定区。平衡是自由能的凸包，两相须同$T,p,\mu$；对摩尔体积作Maxwell等面积构造正是相同化学势条件的积分表达，而非把负压缩率支保留下来。
\originalsection{平均场指数与临界涨落的入口}
临界附近取密度差$\phi=n-n_c$，消去平移后自由能近$F/V=f_0+r\phi^2/2+u\phi^4/4-h\phi$，$u>0$。极小条件$r\phi+u\phi^3=h$给$\phi\propto(-r)^{1/2}$、$\chi\propto r^{-1}$、临界$\phi\propto h^{1/3}$，即$\beta=1/2,\gamma=1,\delta=3$，热容有有限跳跃$\alpha=0$。这些指数忽略空间涨落；第二编将通过Gaussian关联、Ginzburg准则与RG说明何时可用、何时需要新的固定点。一般临界行为不能仅由局部势能的一阶极小化决定。

\originalchapter{简并量子气体与超流}
\label{ch:qfluid}
\source{8.333 L23--L26；理想气体先修见既有卷，重点为低能激发与响应。}
\originalsection{Fermi面与低温热激发}
单态占据$\bar n_k=[\e^{\beta(\epsilon_k-\mu)}+1]^{-1}$由$\Xi_k=1+z\e^{-\beta\epsilon_k}$直接得。零温费米面是占据突变的动量边界。低温只有其附近宽$T$的能量窗口可热激发，固定粒子数的Sommerfeld展开给$C/V=\pi^2D(E_F)\kb^2T/3$。一般晶格或各向异性色散，应以实际费米面$\int\dd S/[(2\pi)^d|\nabla\epsilon|]$计算态密度，不把自由气的球面公式当通用式。

相互作用Fermi液体的准粒子寿命率近$E^2+T^2$，其速度与热态密度可重整化；磁化率、压缩率还含Landau参数，详细构造见新凝聚态卷。统计力学中的“低温像自由准粒子”不意味着原相互作用消失，也不意味着所有观测都只需把$m$换成$m^*$。
\originalsection{Bose凝聚与超流的区别}
玻色占据$\bar n_k=[\e^{\beta(\epsilon_k-\mu)}-1]^{-1}$要求$\mu\le\epsilon_{\min}$。在三维抛物线气体，激发态能容纳的密度上限$n_{\mathrm{ex}}=\zeta(3/2)/\lambda_T^3$，超出的粒子进入最低态。均匀$d\le2$理想气有限温的低能积分发散，因此没有同样的有限温凝聚。有限势阱、相互作用与二维BKT需分别讨论。

凝聚是单体密度矩阵的一个本征值$O(N)$，超流是对相位扭曲、旋转或流动的刚度和耗散性质。理想Bose气可凝聚，却在$\epsilon(p)=p^2/2m$下$\min\epsilon/p=0$，无正Landau临界速度。弱排斥作用使Bogoliubov谱$E_p=\sqrt{(p^2/2m)(p^2/2m+2gn)}$低能线性，才给非零声速尺度。
\originalsection{Landau准则、声子与roton}
在速度$u$的流体中，创建一激发同时减少流体动量$p$。Galilean能量变化为$\Delta E=\epsilon(p)-u\cdot p$，若所有激发皆正则弱扰动不能自发激发它们，所以$u_c=\min_p\epsilon(p)/|p|$。$^4$He的声子线性段给$c$，roton谷$\epsilon(p)\simeq\Delta_r+(p-p_0)^2/(2m_r)$可给更小约$\Delta_r/p_0$。涡旋成核、边界与非线性过程可使实际临界速度更小，Landau准则不是任意实验的充分无耗散保证。

用密度激发试态$\ket q=\rho_q\ket0$，归一为$N S(q)$，双对易和规则给激发变分能
\begin{equation}
E_F(q)=\frac{\hbar^2q^2}{2mS(q)}.
\end{equation}
它是Feynman密度试态的能量上界而非所有激发的精确色散。静态结构因子峰表示密度排列的强相关，对应上界变低，解释roton与液体短程结构的联系。
\originalsection{正常密度与两流体}
激发在小相对速度$w$下分布$n_B(\epsilon_p-p\cdot w)$，其动量密度线性为$\rho_n w$，所以
\begin{equation}
\rho_n=\frac1d\int\frac{\dd^dp}{(2\pi\hbar)^d}p^2[-\partial_{\epsilon_p}n_B].
\end{equation}
三维声子$\epsilon=cp$给$\rho_n=2\pi^2(\kb T)^4/(45\hbar^3c^5)$，积分中$\int x^4\e^x/(\e^x-1)^2=4\pi^4/15$。roton附近按Gaussian径向能量积分贡献$\propto T^{-1/2}\e^{-\Delta_r/\kb T}$，数值前因子含$p_0,m_r$。超流密度$\rho_s=\rho-\rho_n$。

两流体总质量流$\rho_su_s+\rho_nu_n$，熵主要由正常部分携带。线性化密度与熵守恒可有第一声（主要密度波）与第二声（主要熵温度波），低热膨胀近似$u_2^2=(\rho_s/\rho_n)Ts^2/c_p$，此处$s,c_p$按单位质量。式中若改成每体积熵须同步换单位，否则量纲会错。两流体是低能局域平衡描述，不从“有BEC”一词自动推出。

\volpart{第二编　统计场论、重整化与集体物态（8.334）}
\originalchapter{弹性、序参量与对称破缺}
\label{ch:landau}
\source{8.334早期讲义对应L1--L3；符号与量纲在此统一。}
\originalsection{从粒子到连续场}
晶体平衡位置$R_i$附近的位移$u_i$在长波可插值为$u(x)$。均匀平移$u\to u+a$不应改变能量，因此最低项含梯度。刚体转动也不产生弹性能，真正的应变为$u_{ij}=(\partial_iu_j+\partial_ju_i)/2$而不是所有位移梯度。各向同性小应变自由能
\begin{equation}
F_{\rm el}=\int\dd^dx\,[\lambda(\nabla\cdot u)^2/2+\mu u_{ij}u_{ij}],
\end{equation}
$\mu>0$且体积模量满足正条件才稳定。Fourier分解成纵横投影$P_L=\hat q\hat q^T,P_T=I-P_L$，静态关联为$\avg{u_i(q)u_j(-q)}=\kb T[P_L/(\lambda+2\mu)+P_T/\mu]_{ij}/q^2$。零波矢整体平移须固定或单独积分，否则Gaussian配分出现零模无限体积。

关联$\int\dd^dq/q^2$在$d\le2$红外发散，长波位移可使二维晶体的通常平移长程序变为代数序，但局部应变含额外$q$因子，仍可有限。这说明破缺哪种对称与观察哪个场同等重要，不能把大位移涨落一概称晶体马上熔化。
\originalsection{Landau--Ginzburg泛函}
序参量$\phi(x)$代表局域平均密度差、磁化或其他慢变量。短程相互作用、缓慢空间变化、相关对称约束允许
\begin{equation}
\mathcal H[\phi]\equiv\frac F{\kb T}=\int\dd^dx\left[\frac K2(\nabla\phi)^2+\frac r2\phi^2+\frac u4\phi^4-h\phi\right].
\end{equation}
这里$\mathcal H$无量纲，$K,r,u,h$已吸收温度与场的单位。场积分必须有微观截止$a$，不是未定义的无穷多个独立点积分。对Ising$\phi\to-\phi$，$h=0$禁奇项；对无这种对称的液气临界点可移场并调参，使最低临界形状同标量模型。$u>0$稳定四阶截断；若$u<0$需六阶正项，可能一阶转变或三临界。

均匀鞍点$r\phi+u\phi^3=h$给平均场指数。空间Euler方程$-K\nabla^2\phi+r\phi+u\phi^3=h$也给非均匀构造。例如$r<0,h=0$的平面畴壁$\phi(x)=\phi_0\tanh[x/(\sqrt2\xi)]$，$\phi_0=\sqrt{-r/u},\xi=\sqrt{K/|r|}$，代回可核对二阶导与局域势平衡。界面张力$\sigma=\kb T\int K(\partial_x\phi)^2\dd x=2\sqrt2\kb T\sqrt K|r|^{3/2}/(3u)$，不是单纯两极小值的自由能差。
\originalsection{连续对称与Goldstone模}
$n$分量$O(n)$模型将$\phi^2$换成$\vct\phi^2$。$r<0$的均匀极小在半径$\phi_0$球面上，选$n$方向后写纵$\sigma$与$n-1$横$\pi$。二次能量为$\frac12K(\nabla\sigma)^2+|r|\sigma^2+\frac12K(\nabla\pi)^2$。纵模有质量平方$2|r|/K$，横模没有局域恢复项，因为均匀转向不改能量。这是Goldstone模的具体来源。

热横涨落$\avg{\pi^2}\propto(n-1)T\int\dd^dq/(Kq^2)$，$d\le2$会在无限尺度增长。对有限温、短程、连续内部对称的低维体系，这提示通常长程序被破坏；严谨Mermin--Wagner定理还要求其标准相互作用与对称假设。离散Ising无连续转向Goldstone模，二维可以有有限温序。外场、各向异性和长程力可给质量或改变红外幂，不能套同一禁止结论。
\originalsection{临界模式与对称允许项}
平移、旋转、内部对称以及守恒量决定允许的低能项；忽略一个对称允许且相关的项会得到错误普适类。例如有各向异性$\sum_i\phi_i^4$时并非严格$O(n)$，而有长程势$|q|^\sigma$时梯度核不再$q^2$。Landau泛函的作用是把微观细节组织为耦合常数，下一步RG才决定哪些细节在长距离留下。仅写出四次多项式并不证明它是某材料的完整有效理论。

\originalchapter{Gaussian关联与Ginzburg准则}
\label{ch:gaussian}
\source{8.334 L4--L5及L8的Gaussian直接解。}
\originalsection{Gaussian积分与生成泛函}
有限维正定矩阵$A$有$Z[J]=\int\dd^nx\exp[-x^TAx/2+J^Tx]=(2\pi)^{n/2}(\det A)^{-1/2}\exp[J^TA^{-1}J/2]$。证明是完成平方$x\to x-A^{-1}J$再以正交矩阵对角化；若有零本征值或负本征值，积分并不存在为同一有限数。连续场在截断动量或有限晶格上使用同一公式，最后再取受控极限。

对$\mathcal H_0=\frac12\int[K(\nabla\phi)^2+r\phi^2]$，$G(q)=\avg{\phi_q\phi_{-q}}=1/(Kq^2+r)$、$\xi=\sqrt{K/r}$。配分自由能密度的Gaussian涨落项$\frac12\int\log(Kq^2+r)\dd^dq/(2\pi)^d$包含截止依赖的常数和解析背景，物理临界奇异部分需与之区分。

Wick定理由$Z[J]$的指数二次形式直接生成：奇数零均值矩为0，偶数矩为所有配对的$G$乘积。四点$\avg{1234}=G_{12}G_{34}+G_{13}G_{24}+G_{14}G_{23}$；连通四点在纯Gaussian模型为0，加入相互作用后才出现。
\originalsection{实空间相关与散射}
三维Fourier反演$G(r)=\e^{-r/\xi}/(4\pi Kr)$，可通过$(-K\nabla^2+r)G=\delta$的径向Green方程得到。临界$r\to0$，一般$d>2$的长距离$G(x)\propto1/|x|^{d-2}$。因此Gaussian指数$\eta=0,\nu=1/2$，静态磁化率$\chi=G(q=0)=1/r$。这里$r$同时表示质量参数和有时实空间距离容易混淆，正文距离用$|x|$或明确下标。

实验小角散射测结构因子$S(q)$近$1/(q^2+\xi^{-2})$的Lorentz峰，临界相关长度增长使峰变窄。真实临界$S(q)\sim q^{-2+\eta}$，非零$\eta$表示长距离场尺度已被相互作用修正。有限分辨率、样品尺寸与非临界背景须先校正，再把峰宽解释为$\xi^{-1}$。
\originalsection{涨落修正的尺度}
有序相鞍点$\phi_0^2=|r|/u$。一个相关体积内粗粒化的涨落方差主要来自$q\lesssim\xi^{-1}$，$\avg{(\delta\phi)^2}_\xi\sim\int_0^{1/\xi}\dd^dq/[Kq^2+O(|r|)]\sim K^{-d/2}|r|^{(d-2)/2}$。将它与$\phi_0^2$比较，
\begin{equation}
\frac{\avg{(\delta\phi)^2}_\xi}{\phi_0^2}\sim uK^{-d/2}|r|^{(d-4)/2}\ll1
\end{equation}
是平均场自洽条件。$d>4$临界附近比值下降，$d<4$增长，$d=4$出现边际对数。这给上临界维数4。它不是简单比较某点的全紫外方差；后者通常含大的非临界截止项。

一回路$tadpole$修正$\delta r\propto u\int\dd^dq/(Kq^2+r)$可移临界温度，导数$\partial_r\delta r\propto-u\int\dd^dq/(Kq^2+r)^2$在$d\le4$有红外奇异。这与Ginzburg比较是同一失控尺度的两种表现。质量偏移依微观截止，普适临界指数由后续尺度流的线性化得到。
\originalsection{下与上临界维数}
下临界维数问涨落或缺陷是否连有序相都破坏，上临界维数问临界指数何时由Gaussian/平均场决定。连续$O(n>1)$的热自旋波下临界2，标量Ising下临界1，普通短程$\phi^4$上临界4。二维XY在边界维数仍可有代数序与BKT，不能把“无通常长程序”解释成所有相关都短程；二维Ising则在其下临界之上、上临界之下，有非Gaussian有限温临界点。

\originalchapter{尺度假设、有限尺寸与重整化}
\label{ch:scaling}
\source{8.334 L6--L8。}
\originalsection{自由能的齐次结构}
以$t=(T-T_c)/T_c$为温度尺度场、$h$为外场，假设奇异自由能在稳定临界固定点附近满足$f_s(t,h)=b^{-d}f_s(tb^{y_t},hb^{y_h})$。这个齐次形式来自粗粒化后每个相关块贡献有限自由能及尺度场的线性流，是需要临界尺度区的假设，不是任意函数恒等式。

选$b=|t|^{-1/y_t}$，得$f_s=|t|^{d/y_t}\Phi_\pm(h/|t|^{y_h/y_t})$。定义$\xi\sim|t|^{-\nu}$、$M\sim(-t)^\beta$、$\chi\sim|t|^{-\gamma}$、$C_s\sim|t|^{-\alpha}$、临界$M\sim h^{1/\delta}$，对该自由能求导得到
\begin{equation}
\nu=\frac1{y_t},\quad\alpha=2-d\nu,\quad\beta=\nu(d-y_h),\quad\gamma=\nu(2y_h-d),\quad\delta=\frac{y_h}{d-y_h}.
\end{equation}
临界关联$G(x)\sim|x|^{-(d-2+\eta)}$给场尺度$x_\phi=(d-2+\eta)/2$，源项$\int h\phi$无量纲要求$y_h=d-x_\phi$，从而$\gamma=\nu(2-\eta)$。这些推导解释指数关系为何不是独立猜测。
\originalsection{假设边界与危险无关变量}
$\alpha=2-d\nu$称超尺度关系，要求只有$\xi$控制奇异自由能体积尺度。在$d>4$的$\phi^4$，$u$在RG中流向0，却有序相$\phi_0\sim\sqrt{|r|/u}$，不能在求振幅前置$u=0$。这种危险无关变量使通常超尺度失效，平均场$\nu=1/2,\alpha=0$不会在任意$d$满足$2-d\nu$。边际维数还可有对数；BKT的$\xi$为本质奇异，也不能硬套普通温度幂指数。
\originalsection{RG变换与相关、无关方向}
先积分掉尺度$a$到$ba$之间的短波，再把坐标缩回原截止，最后重标场。所有对称允许的耦合组成向量$g$，变换$g'=R_b(g)$。固定点$R_b(g_*)=g_*$，在其附近$\delta g_i'=\sum_jM_{ij}\delta g_j$，对角尺度方向的本征值$b^{y_i}$。$y_i>0$相关、$<0$无关、$=0$边际，边际需看非线性项。对相关方向的名字应与所选基相关，原始温度可混入其他参数。

多个微观模型流到同一固定点，就有相同长距离指数与经尺度单位归一的普适函数；“所有材料都一样”并非普适性的含义。对称、维数、相互作用范围、守恒与拓扑等改变可使流到不同固定点。

Gaussian模型在$x=bx'$、$\phi(x)=b^{-(d-2)/2}\phi'(x')$下保持梯度项，质量$r'=b^2r$，四次耦合$u'=b^{4-d}u$。这直接给$y_r=2,y_u=4-d$与上临界维数。后续相互作用RG将这些工程维数修正。
\originalsection{有限尺寸标度}
当样品边长$L$限制相关长度，临界量依赖$tL^{1/\nu}$。例如$M=L^{-\beta/\nu}\mathcal M(tL^{1/\nu},hL^{y_h})$、$\chi=L^{\gamma/\nu}\mathcal X(tL^{1/\nu})$。Binder比$U_L=1-\avg{M^4}/[3\avg{M^2}^2]$无量纲，在相同几何边界下临界曲线近交于普适值。有限尺寸修正$L^{-\omega}$、外形及边界都会移动交点，不能用两条小样品曲线交点声称精确$T_c$。

有限尺寸谱或模拟数据做塌缩时，应先固定可检验的指数、误差与修正方案，再检查不同尺寸是否一致。任意调横纵坐标直到看起来重合不构成独立指数验证。

\originalchapter{动量壳RG与Wilson--Fisher固定点}
\label{ch:wilson}
\source{8.334 L9--L12；采用$u/4!$规范统一图的组合系数。}
\originalsection{快慢模与连接展开}
设$\mathcal H=\int[\frac12(\nabla\vct\phi)^2+\frac r2\vct\phi^2+\frac u{4!}(\vct\phi^2)^2]$，$n$分量，截止$\Lambda$。分$\phi=\phi_<+\phi_>$，慢模$|q|<\Lambda\e^{-\dd l}$，快模在薄壳。有效作用量$\mathcal H_{\rm eff}[\phi_<]=-\log\avg{\e^{-\mathcal H_{\mathrm{int}}}}_>$加慢Gaussian项，展开为$\avg{\mathcal H_{\mathrm{int}}}_>-\frac12\avg{\mathcal H_{\mathrm{int}}^2}_{>,c}+\cdots$。负二阶号来自对数，不能把各阶相同符号堆上去。

Gaussian快模收缩$\avg{\phi_{>,a}(q)\phi_{>,b}(-q)}=\delta_{ab}/(q^2+r)$。一次四腿顶角中收缩两快腿给质量tadpole；两个顶角各留两慢腿并两条快传播子相连给四次顶角修正。只保留连接图避免重复计入真空归一。
\originalsection{一回路系数的计算}
薄壳积分$I_1=K_d\Lambda^{d-2}\dd l/(1+\widetilde r)$、$I_2=K_d\Lambda^{d-4}\dd l/(1+\widetilde r)^2$，$K_d=S_{d-1}/(2\pi)^d$、$\widetilde r=r/\Lambda^2$。为算出$n$的组合系数，写$\phi=\varphi+\chi$（慢、快），四次顶角中恰有两条快腿的部分为
\begin{equation}
V_2=\frac u{24}\int[2\varphi^2\chi^2+4(\varphi\cdot\chi)^2]
=\int\chi_a B_{ab}(\varphi)\chi_b,\qquad
B_{ab}=\frac u{12}(\varphi^2\delta_{ab}+2\varphi_a\varphi_b).
\end{equation}
快场同点收缩$\langle\chi_a\chi_b\rangle=\delta_{ab}I_1$，所以$\langle V_2\rangle=(n+2)uI_1\int\varphi^2/12$；与$r\varphi^2/2$比较即$\delta r=(n+2)uI_1/6$。二阶连接收缩$\langle\chi_a(x)\chi_b(x)\chi_c(y)\chi_d(y)\rangle_c=G_>(x-y)^2(\delta_{ac}\delta_{bd}+\delta_{ad}\delta_{bc})$。两条快线积分给$I_2$，而
\begin{equation}
\operatorname{tr}B^2=\frac{u^2}{144}(n+8)(\varphi^2)^2,\qquad
-\frac12\langle V_2^2\rangle_c=-I_2\int\operatorname{tr}B^2.
\end{equation}
与$u(\varphi^2)^2/24$比较，$\delta u=-(n+8)u^2I_2/6$。这里的$n$是闭合指标环，额外的2和8来自纵向投影与两种跨顶角Wick配对；例如$n=1$得$\delta r=uI_1/2$、$\delta u=-3u^2I_2/2$。零外动量的局部顶角近似要求外部慢波矢远小于薄壳动量。

缩回截止后，以$g=K_du\Lambda^{d-4}$，$\epsilon=4-d$，得到
\begin{align}
\frac{\dd\widetilde r}{\dd l}&=2\widetilde r+\frac{n+2}{6}\frac g{1+\widetilde r}+O(g^2),\\
\frac{\dd g}{\dd l}&=\epsilon g-\frac{n+8}{6}\frac{g^2}{(1+\widetilde r)^2}+O(g^3).
\end{align}
非普适临界质量被tadpole移为负，不能直接把$r=0$当相互作用模型的精确临界温度。固定点$g_*=6\epsilon/(n+8)+O(\epsilon^2)$，$\widetilde r_*=-(n+2)\epsilon/[2(n+8)]+O(\epsilon^2)$。
\originalsection{临界指数与控制参数}
在临界面附近，温度方向本征值$y_t=2-(n+2)\epsilon/(n+8)+O(\epsilon^2)$，所以$\nu=1/2+(n+2)\epsilon/[4(n+8)]+O(\epsilon^2)$。一回路tadpole与外动量无关，故这一阶$\eta=0$。由前章指数关系，
\begin{equation}
\gamma=1+\frac{n+2}{2(n+8)}\epsilon+O(\epsilon^2),\qquad\beta=\frac12-\frac3{2(n+8)}\epsilon+O(\epsilon^2),\qquad\delta=3+\epsilon+O(\epsilon^2).
\end{equation}
这是一种在$\epsilon\ll1$且固定点$g_*=O(\epsilon)$下可控的非Gaussian展开。将$\epsilon=1$代入可估三维，却有未计高阶误差；不应把一回路三维数当精确值。$n=1$给$\nu\simeq7/12$，与精确二维$\nu=1$的远维数结果更不能仅以一阶硬外推。

两回路首次产生外动量依赖的sunset自能。课程给出的标准$O(n)$结果$\eta=(n+2)\epsilon^2/[2(n+8)^2]+O(\epsilon^3)$需同阶波函数与耦合重整化。此处引用其结果定位两回路的作用，主稿完整算到一回路，未把一回路图称作该二阶系数的证明；相应作业中若要求具体二阶壳图，按其模型另算。
\originalsection{其他相互作用与边际维数}
算符$\phi^{2p}$的Gaussian耦合尺度$y_{2p}=d-p(d-2)$，在$d=4$、$p>2$为负；额外梯度也降低工程相关性。因此近四维的普通四次固定点只需有限个相关调参，而不是每个微观项都同等控制临界。但当四次项也被调到0，六次项可变为三临界主项，上临界维数改为3，不能仍删它。

$d=4$时$g'=-(n+8)g^2/6$，$g(l)=g_0/[1+(n+8)g_0l/6]$缓慢趋0。把这个$1/\log$流放入质量与场响应会产生平均场幂乘对数修正，这正是边际项需要非线性流的具体例子。

\originalchapter{实空间RG与Monte Carlo}
\label{ch:realrg}
\source{8.334 L13--L14与课程大纲中的格点模拟。}
\originalsection{一维Ising精确消元}
一维零场$\mathcal H=-K\sum_i s_is_{i+1}$，$s_i=\pm1$。积分每个偶站使$\sum_{s_2}\e^{K(s_1s_2+s_2s_3)}=2\cosh[K(s_1+s_3)]=A\e^{K's_1s_3}$。比较$s_1=s_3$与$s_1=-s_3$，得$K'=\frac12\log\cosh2K$、$A=2\sqrt{\cosh2K}$。有限$K>0$有$K'<K$，流向高温无序点；$K=\infty$为零温临界端点。

精确相关$\avg{s_0s_r}=(\tanh K)^r$给$\xi^{-1}=-\log\tanh K$。低温定义热缺陷逸度$t=\e^{-2K}$，$\xi\simeq1/(2t)$，消元$t'\simeq2t$给$y_t=1$。此$t$不是线性$T-T_c$，所以不能写一维$\xi\propto T^{-1}$；它随$\e^{2J/\kb T}$增长。
\originalsection{块自旋与累积近似}
高维选择块变量$S_B$并保持原自旋与块变量的条件权重，$\e^{-\mathcal H'[S]}=\sum_sP(S|s)\e^{-\mathcal H[s]}$。精确变换会产生更多邻距、多体和场耦合，一般不是只剩一个$K'$。

Niemeijer--van Leeuwen做法将块内$H_0$精确处理、块间$V$按连接累积展开：$\mathcal H'=-\log Z_0+\avg V_0-\avg{V^2}_{0,c}/2+\cdots$。块内条件分布由所选块自旋定义决定，$\avg s_i$与$\avg{s_is_j}_c$构成重整化系数。截断阶数与舍弃生成耦合必须明示，它不是高维Ising的精确消元。

Migdal--Kadanoff近似用移键与消元改变连通结构。先串联消元再并联的一种$b$倍次序给
\begin{equation}
K'=b^{d-1}\operatorname{atanh}[(\tanh K)^b].
\end{equation}
$d=1$时是精确串联，$d>1$改变图的连通结构。$d=2,b=2$有$K'=\log\cosh2K$，临界不动点约0.609，而方格精确值为0.4407。线性化$\nu=\log b/\log R'(K_c)$可检查估计精度。本卷官方复习题采用另一种先移键后消元次序，数值另列；两种都属近似，不能混用。近似仍保留两个流域与临界相关方向，物理结构可比数值更可靠。
\originalsection{详细平衡与采样误差}
对构型$s$取目标$\pi(s)=\e^{-\mathcal H(s)}/Z$。对称提议翻一自旋，接受$P_{\mathrm{acc}}=\min(1,\e^{-\Delta\mathcal H})$。两态概率比恰为$\e^{-\Delta\mathcal H}$，故详细平衡$\pi(s)P(s\to s')=\pi(s')P(s'\to s)$。加可达性与非周期性保证有限状态Markov链趋目标分布；详细平衡本身不保证混合快。

能量或磁化的连续采样相关。定义积分自相关时间
\begin{equation}
\tau_{\mathrm{int}}=\frac12+\sum_{t>0}\frac{C(t)}{C(0)},\qquad
\operatorname{Var}(\overline A)\simeq\frac{2\tau_{\mathrm{int}}}{N_{\mathrm{samples}}}\operatorname{Var}(A).
\end{equation}
热化前样本应舍去；跨临界慢化时必须用更长轨迹、块平均或簇算法，不能把每一步当独立样本制造过小误差。
\originalsection{临界模拟的可复核对象}
周期$L\times L$方格的能量$E=-J\sum_{\avg{ij}}s_is_j$，翻自旋的$\Delta E=2Js_i\sum_{j\sim i}s_j$。测$C=\beta^2(\avg{E^2}-\avg E^2)/L^2$、$\chi=\beta L^2(\avg{m^2}-\avg m^2)$和Binder比，用多个$L$检查标度。有限零场系统对称性常使$\avg m=0$，报告$|m|$作为序参量时磁化率定义仍要小心，不能未经说明把$\avg m^2-\avg{|m|}^2$当严格线性响应。

本卷数值附档以固定种子执行短规模模拟并对照一维精确结果、Gaussian积分与RG固定点。它验证程序与计算步骤，不以少量样品替代热力学极限证明。作业需要的具体图、误差与有限尺度范围另按题登记。

\originalchapter{级数、闭环与对偶}
\label{ch:duality}
\source{8.334 L15--L17。}
\originalsection{高温展开与低温畴壁}
Ising键$\e^{Ks_is_j}=\cosh K(1+s_is_j\tanh K)$，逐键相乘后对所有自旋求和，任何顶点若有奇数入边则该自旋求和为0。故$Z=2^N(\cosh K)^{N_b}\sum_{\rm even\ graphs}(\tanh K)^{|G|}$，图在每顶点偶度，在二维可分成闭环但交叉与重叠的计数须保留。

低温从全同自旋出发，翻一个畴的能量只随边界长度：每条不满足键代价$2J$。于是低温图权重$\e^{-2K\ell}$，图画在对偶格。Peierls估计在二维把穿过某邻域的长度$\ell$闭环数上界为常数乘$3^\ell$，若$3\e^{-2K}<1$，长畴壁总概率小，可证明足够低温有自发磁化。这个粗界证明存在有序区，不给精确$T_c$。
\originalsection{一维传递矩阵}
含场$h$时$T_{ss'}=\exp[Kss'+h(s+s')/2]$，$Z=\Tr T^N$。本征值$\lambda_\pm=\e^K\cosh h\pm\sqrt{\e^{2K}\sinh^2h+\e^{-2K}}$，无限链$f=-\kb T\log\lambda_+$。有限$K$最大根光滑且唯一，因此无有限温非解析相变。零场$\lambda_-/\lambda_+=\tanh K$，再次得到相关长度。

高温闭环或低温缺陷在一维都可完整求和，适合核对展开系数和有限尺寸实现。把一维无相变推到二维的错误在于域壁熵随长度的增长和环结构已变化。
\originalsection{二维自对偶与三维规范对偶}
方格的低温畴壁与对偶格高温偶图一一对应，令$\e^{-2K^*}=\tanh K$。扣去显式$2^N,\cosh K$等解析因子，$Z(K)$映到$Z(K^*)$。自对偶条件$\sinh2K_c=1$给$K_c=\frac12\log(1+\sqrt2)$，但只凭自对偶还须知道相变位置与唯一性；下章精确自由能核对此唯一临界奇异。

三维Ising低温图是闭合的二维畴壁面，而非二维闭线。$\mathbb Z_2$格点规范模型边变量$u_l=\pm1$、作用$-K_g\sum_p\prod_{l\in\partial p}u_l$的高温展开也给闭面，因为对每条边求和要求相邻选面数偶。因此$\e^{-2K_{\mathrm{Ising}}}=\tanh K_g$联系三维Ising与对偶规范模型。不能把三维模型错误称为与同一三维Ising自对偶。
\originalsection{随机游走生成与phantom环}
Gaussian格点场的二次矩阵写$A=I-w\mathsf J$，$\mathsf J$为邻接矩阵，则$\log Z=-\frac12\log\det A+\mathrm{const}=\frac12\sum_{l\ge1}w^l\Tr\mathsf J^l/l+\mathrm{const}$。$\Tr\mathsf J^l$计起点相同的闭游走，$1/l$除去起点重复，$1/2$与无向或实场计数相联。这些游走可自交、重复走边，故称phantom环。

它可总结高温图的Gaussian近似，解释高维临界的均场结构，但不是把每个Ising偶图精确计一次。二维精确求和还需转角相位，使交叉路径的过计数抵消（Kac--Ward/Pfaffian形式）。图的权重与符号结构不能只用“所有闭环都有同样权”替代。

\originalchapter{方格Ising的精确自由能}
\label{ch:ising2}
\source{8.334 L18--L19给Kac--Ward精确自由能并说明费米传递解释；以下Jordan--Wigner对角化为编者补足的等价计算链。}
\originalsection{行传递与对偶耦合}
考虑无限二维方格、零场、各向同性最近邻$J>0$。一行$N$自旋间的垂直键矩阵每站为$\left(\begin{smallmatrix}\e^K&\e^{-K}\\\e^{-K}&\e^K\end{smallmatrix}\right)=\sqrt{2\sinh2K}\,\e^{K^*\sigma^x}$，$\tanh K^*=\e^{-2K}$。水平键对角，因此对称行传递矩阵可写
\begin{equation}
V=(2\sinh2K)^{N/2}\e^{(K/2)\sum\sigma_j^z\sigma_{j+1}^z}\e^{K^*\sum\sigma_j^x}\e^{(K/2)\sum\sigma_j^z\sigma_{j+1}^z}.
\end{equation}
每个因子由显式矩阵恒等式得到，$Z=\Tr V^{N_y}$，大$N_y$时最大本征值控制自由能。
\originalsection{Majorana化与每模对角化}
为了不把“化为自由费米子”当黑箱，取一行的串算符$S_j=\prod_{\ell<j}\sigma_\ell^x$，定义两个Hermitian Majorana算符
\begin{equation}
a_j=S_j\sigma_j^z,\qquad b_j=S_j\sigma_j^y,
\qquad\{a_j,a_k\}=\{b_j,b_k\}=2\delta_{jk},\quad\{a_j,b_k\}=0.
\end{equation}
沿串移动一个站点会得到一次反对易；同站用Pauli乘法直接核对$\sigma_j^x=\ii a_jb_j$及$\sigma_j^z\sigma_{j+1}^z=\ii b_ja_{j+1}$。故上节$V=PQP$的两种指数均为Majorana双线性，其中
\begin{equation}
P=\exp\!\left(\frac{\ii K}{2}\sum_jb_ja_{j+1}\right),\qquad
Q=\exp\!\left(\ii K^*\sum_ja_jb_j\right).
\end{equation}
周期自旋行的末键带总宇称$\mathcal P=\prod_j\sigma_j^x$，在每个固定宇称扇区选对应的周期或反周期Majorana边界。以下$q$为该扇区允许的动量；两个动量网格相差半格，体极限相同。

取$a_q=N^{-1/2}\sum_j\e^{-\ii qj}a_j$，$b_q$同。由$[\ii\gamma_1\gamma_2,\gamma_1]=-2\ii\gamma_2$与另一条对易式，可逐项计算$P,Q$共轭作用于列向量$(a_q,b_q)^{\mathsf T}$。记$c=\cosh K,s=\sinh K$以及$c_*=\cosh2K^*,s_*=\sinh2K^*$，两种作用矩阵分别为
\begin{equation}
H(q)=\begin{pmatrix}c&\ii s\e^{-\ii q}\\-\ii s\e^{\ii q}&c\end{pmatrix},
\qquad
Q_*(q)=\begin{pmatrix}c_*&-\ii s_*\\\ii s_*&c_*\end{pmatrix}.
\end{equation}
例如$Pa_jP^{-1}=ca_j+\ii s b_{j-1}$，$Pb_jP^{-1}=cb_j-\ii s a_{j+1}$；Fourier变换正给$H(q)$中的两个相位。对称行传递的$2\times2$块为$M(q)=H(q)Q_*(q)H(q)$。两因子行列式均为1，因此$\det M=1$；利用迹的循环性，$\operatorname{tr}M=\operatorname{tr}[Q_*H^2]$，且$H^2$只需把$K$改成$2K$。乘开两项对角元，得到
\begin{equation}
\frac12\operatorname{tr}M(q)=\cosh\gamma_q
=\cosh2K\cosh2K^*-\sinh2K\sinh2K^*\cos q.
\end{equation}
因$\det M=1$，其本征值必为$\e^{\pm\gamma_q}$；对偶关系$\sinh2K\sinh2K^*=1$又把上式改写为$\cosh\gamma_q=\cosh2K\coth2K-\cos q$。每个正能复费米模的两态传递本征值相差$\e^{\gamma_q}$，最大本征态在每模取$+\gamma_q/2$，故
\begin{equation}
\log\lambda_{\max}=\frac N2\log(2\sinh2K)+\frac12\sum_q\gamma_q.
\end{equation}
有限行的宇称投影会调整允许动量和个别态，却不改上述每站体自由能；没有投影时不可把两个完整费米Fock空间重复计为原自旋空间。

因此每站精确自由能为
\begin{equation}
-\beta f=\frac12\log(2\sinh2K)+\frac1{2\pi}\int_0^\pi\operatorname{arcosh}[\cosh2K\coth2K-\cos q]\dd q.
\end{equation}
高温$K\to0$两个对数奇异相消，给$\log2$；低温给$2K$，即每站两条满意键的基态能$-2J$。两个边界极限是公式归一的直接检查。用$t=\tanh K$将此一维$\operatorname{arcosh}$积分还原为第\ref{ch:duality}章的双动量行列式积分，两式必须给相同自由能；例如$K=0.2,0.4,0.7$时两式数值差分别小于$3\times10^{-15}$。
\originalsection{能隙、相关长度与热容}
$\gamma_0=2|K-K^*|$，只在$K=K^*$闭合，故精确临界满足$\sinh2K_c=1$。近临界，$\gamma_q\simeq\sqrt{m^2+q^2}$，$m=2|K-K^*|\propto|T-T_c|$，因此$\xi\sim m^{-1}$、$\nu=1$。奇异自由能来自$\int_0^\Lambda\sqrt{m^2+q^2}\dd q$的$-m^2\log|m|$项，二阶温度导数给对数热容，$\alpha=0$带对数，而非有限平均场跳跃。

零场自发磁化的精确关联结果为$m(T)=[1-\sinh^{-4}(2K)]^{1/8}$于$K>K_c$，高温0，故$\beta=1/8$。这一磁化公式需要非局域自旋相关的Toeplitz/Pfaffian分析；本章上述单粒子传递谱完整求出自由能，却没有把它冒称为磁化公式的全证明。课程L19引用该标准结果，本稿保留这个证明范围。

临界自旋相关$G(r)\sim r^{-1/4}$给$\eta=1/4$，配尺度关系得$\gamma=7/4,\delta=15$。这些不同于近四维一回路外推，说明二维强涨落固定点需要真正的非Gaussian结构。含外场、非最近邻或一般三维Ising没有此自由费米传递解，不能把零场精确公式用于任意Ising模型。
\originalsection{可计算的数值核对}
随卷计算用稳定的$\operatorname{arcosh}$与数值积分对照高低温展开，并以小周期方格枚举所有$2^{L^2}$构型核对有限尺寸配分与能量；无限体公式与小$L$有有限尺寸误差，不强求相等。积分在临界端点的导数奇异需单独处理，普通粗网格微分会低估对数峰。

\originalchapter{连续自旋、涡旋与BKT转变}
\label{ch:bkt}
\source{8.334 L20--L22。}
\originalsection{非线性sigma模型}
$O(n)$低温固定幅度场$n(x)$满足$n^2=1$，作用$\mathcal H=(1/2g)\int(\nabla n)^2$，$g$与温度/刚度之比有关。参数化$n=(\pi,\sqrt{1-\pi^2})$，展开得$\mathcal H=(1/2g)\int[(\nabla\pi)^2+(\pi\cdot\nabla\pi)^2+\cdots]$，约束造成梯度相互作用，不能把$n$的$n$个分量都当独立Gaussian。

在$d=2+\epsilon$，低温动量壳积分的自旋波重整化给$\dd g/\dd l=-\epsilon g+(n-2)g^2/(2\pi)+\cdots$。$n>2,d=2$耦合随尺度增长，相关长约$a\exp[2\pi/((n-2)g)]$，自旋波弱耦合最终失控。这个一回路系数来自约束相互作用与场重标的组合，$n=2$时其抵消与可用单角度Gaussian表述一致；但是涡旋仍是非微扰自由度。
\originalsection{XY自旋波与代数序}
$n=2$写$n=(\cos\theta,\sin\theta)$，$\mathcal H=K\int(\nabla\theta)^2/2$，$K=\rho_s/(\kb T)$。在二维高斯场$\avg{[\theta(r)-\theta(0)]^2}=\log(r/a)/(\pi K)$，所以$\avg{\e^{\ii\theta(r)}\e^{-\ii\theta(0)}}=(a/r)^{\eta}$，$\eta=1/(2\pi K)$。相位差的对数增长使通常磁化为0，却有幂律关联，正是准长程序。

紧致角还允许$\oint\nabla\theta\cdot\dd l=2\pi m$涡旋。单涡$\theta=m\arg(x+\ii y)$，梯度$m/r$，能量$E_v=\pi\rho_sm^2\log(L/a)+E_c$。可放置位置的熵约$2\kb\log(L/a)$，对单位涡的简单能熵竞争给临界刚度量级$\rho_s/(\kb T)=2/\pi$。因其他涡与自旋波会重整刚度，应使用长距离的$K_R$，不能把这个估计中的裸刚度直接叫精确转变温度。
\originalsection{Coulomb气与介电屏蔽}
多个涡的能量在中性条件$\sum_im_i=0$下为$-2\pi\rho_s\sum_{i<j}m_im_j\log(r_{ij}/a)+E_c\sum_im_i^2$。小涡反涡偶极子在大涡产生的场中极化，屏蔽对数作用、降低有效$K$。赋予单位涡逸度$y=\e^{-E_c/\kb T}$，积分尺寸$a$到$a\e^{\dd l}$的偶极子并重标，标准归一给
\begin{equation}
\frac{\dd K^{-1}}{\dd l}=4\pi^3y^2+O(y^4),\qquad\frac{\dd y}{\dd l}=(2-\pi K)y+O(y^3).
\end{equation}
第二式的2是位置积分维数，$-\pi K$是涡自能随截止变动的尺度代价；第一式正号表明偶极屏蔽只降低刚度，系数依$y$的定义。保持中性与偶极角平均才能获得这个二阶修正。
\originalsection{流、普适跳跃与本质奇异}
$y=0,K>2/\pi$是一条低温固定线，涡无关，相关幂随$K_R$连续变化；$K<2/\pi$时$y$增长，自由涡破坏代数序。分界流抵达$K_R=2/\pi$，故$\rho_s(T_c^-)=2\kb T_c/\pi$，临界$\eta=1/4$。

在端点附近令$x=2-\pi K$并适当重标$Y\propto y$，最低阶$\dot x\propto Y^2,\dot Y=xY$，有双曲不变量$x^2-Y^2$。高温距分界$\tau>0$时流在临界区域停留$l_*\propto1/\sqrt\tau$，故$\xi\sim a\exp(b/\sqrt\tau)$。这是本质奇异，不能赋一个普通有限$\nu$来模拟。$b$依微观逸度，刚度跳跃与临界指数则由固定点结构决定。

有限系统即使无穷体系无磁化也可有大$\avg{|M|}$；应测关联与扭曲刚度随$L$的对数修正，再区分BKT与普通幂律临界。孤立单涡不能在严格周期环面中净存在，模拟要用中性涡对或相应拓扑扇区。

\originalchapter{二维晶体与拓扑熔化}
\label{ch:melting}
\source{8.334 L23。}
\originalsection{平移代数序与取向}
二维弹性Gaussian场的位移差方差按$\log r$增长。晶体密度的倒格分量$\rho_G(x)\propto\e^{\ii G\cdot[x-u(x)]}$，故关联$\avg{\rho_G(r)\rho_G(0)^*}\propto r^{-\eta_G}$，$\eta_G$由纵横弹性常数及$G$决定。取向角近$\omega=(\partial_xu_y-\partial_yu_x)/2$，含一个梯度，热长波角涨落在无缺陷晶体中有限。所以二维固体可有代数平移序而长期取向序，并不与位移红外发散矛盾。
\originalsection{位错与位向旋错}
沿闭路位移增量$\oint\partial_i u_j\dd x_i=b_j$为Burgers向量，非零则存在位错；它不能由单值光滑位移场消掉。对六角晶格，键取向绕路改$\pm\pi/3$为基本位向旋错。一个位错可视为相邻反号旋错偶极，两个缺陷层级的能量尺度不同，因而可能两步解束缚。

弹性作用求出的单个位错能$E_d\simeq Yb^2\log(L/a)/(8\pi)+E_c$，$Y=4\mu(\lambda+\mu)/(\lambda+2\mu)$是二维Young模量。位置熵$2\kb\log(L/a)$给长距离重整模量的解束缚条件$Ya^2/(\kb T)=16\pi$，取最短Burgers长度$a$。用裸模量代替已屏蔽模量只是一阶估计。
\originalsection{hexatic与两步图景}
位错解束缚破坏平移代数序，但仍可有六重键角序$\psi_6=\e^{6\ii\theta}$，称hexatic相。其长波角自由能$F_\theta=K_A\int(\nabla\theta)^2/2$。旋错电荷$1/6$，同一涡能熵分析给$K_A/(\kb T)=72/\pi$临界值，$\psi_6$关联指数在临界为$1/4$。再解束缚旋错后为各向同性液体。

KTHNY图景要求缺陷核心代价使连续解束缚流可达相应固定点；某些材料或相互作用可能因核心能低而出现一阶熔化、晶界机制或其他中间结构。因此“所有二维固体必经hexatic”不是该理论的结论。解释实验应同时看平移和取向相关、弹性模量、缺陷密度与有限尺寸，不能仅凭一个结构因子峰宽标成hexatic。

\originalchapter{Langevin、Fokker--Planck与临界动力学}
\label{ch:dynamics}
\source{8.334 L24--L25。}
\originalsection{Brownian粒子的平衡结构}
过阻尼粒子满足$\dot x=-\Gamma U'(x)+\eta(t)$，白噪声$\avg\eta=0$、$\avg{\eta(t)\eta(t')}=2D\delta(t-t')$。短时间增量$\Delta x=-\Gamma U'\Delta t+\sqrt{2D\Delta t}\,\zeta$，$\zeta$标准Gaussian；对任意平滑测试函数Taylor展开至$O(\Delta t)$，平均生成算符$\mathcal Lf=-\Gamma U'f'+Df''$。分部积分得到
\begin{equation}
\partial_tP=\partial_x[\Gamma U'P+D\partial_xP].
\end{equation}
概率流$J=-\Gamma U'P-D\partial_xP$。要求无流的$P_{\rm eq}\propto\e^{-\beta U}$，得到Einstein关系$D=\Gamma\kb T$。噪声与摩擦若不满足该关系，不能把稳态擅自称同温度的热平衡。

谐势$U=kx^2/2$的Ornstein--Uhlenbeck过程有$\avg{x(t)x(0)}=(\kb T/k)\e^{-\Gamma k|t|}$，响应$\chi(\omega)=\Gamma/(\Gamma k-\ii\omega)$，噪声谱$S_x=2D/[(\Gamma k)^2+\omega^2]$，直接满足经典涨落耗散$S_x=2\kb T\chi''/\omega$。这提供从随机方程到可测谱的完整例子。
\originalsection{非守恒序参量：Model A}
对场的物理自由能$F[\phi]$，取$\partial_t\phi=-\Gamma\delta F/\delta\phi+\zeta$，$\avg{\zeta(x,t)\zeta(x',t')}=2\Gamma\kb T\delta(x-x')\delta(t-t')$。泛函Fokker--Planck流同粒子情形，保证$P[\phi]\propto\e^{-\beta F}$。若Gaussian$F=\frac12\int[r\phi^2+K(\nabla\phi)^2]$，每$q$模衰减率$\Gamma(r+Kq^2)$，$\tau\sim r^{-1}\sim\xi^2$，Gaussian动态指数$z=2$。

响应$R(q,\omega)=[-\ii\omega+\Gamma(r+Kq^2)]^{-1}$，结构谱$C(q,\omega)=2\Gamma\kb T|R|^2$。频率积分给静态$\avg{|\phi_q|^2}=\kb T/(r+Kq^2)$，恢复前章Gaussian结果，这验证噪声归一及静动态一致性。
\originalsection{守恒序参量：Model B}
若$\int\phi$必须守恒，动力学应由电流$J=-\Gamma\nabla(\delta F/\delta\phi)+\eta_J$形成，$\partial_t\phi=\Gamma\nabla^2\delta F/\delta\phi-\nabla\cdot\eta_J$。噪声为电流噪声，谱额外$q^2$，不能继续用空间独立非守恒白噪声。

Gaussian衰减率$\Gamma q^2(r+Kq^2)$，临界$r=0$为$q^4$，$z=4$。相互作用情况下在标准Model B且无额外耦合的条件下守恒与静态异常维数联系为$z=4-\eta$；若序参量耦合动量、能量或其他慢量，可能属于其他动态普适类。相同静态Ising指数的两种模型可有不同弛豫时间。
\originalsection{动态RG与可观测时间}
动态尺度形式$C(q,\omega)=q^{-2+\eta-z}\mathcal C(\omega/q^z,q\xi)$，频率积分需给静态$q^{-2+\eta}$。$z$控制临界慢化，实验扫描速度若快于$\tau(\xi)$，就可能看到非平衡滞后而非平衡指数。局域Markov随机方程已忽略热库记忆，若噪声有长时间相关或惯性重要，需要扩展而不能只改$\Gamma$数值。

响应场路径积分可把噪声Gaussian概率积分，得到包含$\widetilde\phi(\partial_t\phi+\Gamma\delta F/\delta\phi)$与噪声二次项的动态作用量；其传播子是上述推迟响应，相关线为噪声乘两个响应。这个构造提供动态微扰图的起点，不能把静态传播子直接当频率依赖响应而忽略因果性。

\originalchapter{守恒约束与非平衡表面增长}
\label{ch:growth}
\source{8.334 L26：平衡的一般尺度不变性与开放系统增长。}
\originalsection{无临界调参的长程结构}
临界长相关常需调温度到特殊点，但守恒、连续对称或约束也可使长波无质量。不可压缩流体$\nabla\cdot u=0$的均衡动量能$F=\rho\int u^2/2$，允许子空间为横向，故$\avg{u_i(q)u_j(-q)}=(\kb T/\rho)(\delta_{ij}-q_iq_j/q^2)$。尽管振幅无$q$发散，投影在$q=0$方向非解析，实空间含非局域偶极型约束相关。不能仅以“谱有限”判断相关一定纯局域。

这些约束性尺度结构不同于标量液气临界的温度相关方向。开放系统的持续驱动、守恒噪声及模式耦合又可产生新的无临界调参标度，必须从其实际随机方程与对称判断，不能以平衡配分函数自动描述。
\originalsection{Edwards--Wilkinson模型}
表面高度$h(x,t)$相对平均面写$\partial_th=\nu\nabla^2h+\eta$，$\avg{\eta(x,t)\eta(x',t')}=2D\delta^d(x-x')\delta(t-t')$。减掉$q=0$平均高度增长后，每非零模为OU过程，稳态$S_h(q)=D/(\nu q^2)$。

粗糙度$W^2=L^{-d}\int\avg{[h(x)-\bar h]^2}\dd x\sim\int_{1/L}^{1/a}\dd^dq/q^2$。$d<2$给$W\sim L^\alpha$、$\alpha=(2-d)/2$，$d=2$为$W^2\sim\log L$，$d>2$红外有限。动态模时间$\sim q^{-2}$，$z=2$；未饱和时$W\sim t^\beta$，$\beta=\alpha/z=(2-d)/4$。这些指数从精确Gaussian过程来，适于检查模拟和噪声归一。
\originalsection{KPZ非线性与倾斜对称}
若沿表面法向沉积，转成垂直高度速度会出现$(\nabla h)^2$项，KPZ方程为$\partial_th=\nu\nabla^2h+\lambda(\nabla h)^2/2+\eta$。相对平面倾斜$h\to h+v\cdot x$可由坐标漂移$x\to x+\lambda vt$与常数速度补偿，给Galilean倾斜对称。若RG固定点保持该对称且$\lambda$有限，尺度要求$\alpha+z=2$，不是独立地把两个指数都设成Gaussian值。

在一维周期边界，稳态梯度分布仍为$P[h]\propto\exp[-\nu\int(\partial_xh)^2/(2D)]$。Fokker--Planck中非线性流的额外贡献经泛函分部后正比$\int(\partial_xh)^2\partial_x^2h\dd x=\frac13\int\partial_x[(\partial_xh)^3]\dd x=0$。因此$\alpha=1/2$；结合倾斜关系得$z=3/2,\beta=1/3$，不同于EW的$z=2$。这项一维稳态相消一般不在高维成立，也不证明所有噪声或边界下都同属KPZ。
\originalsection{Burgers与Hopf--Cole连接}
令斜率$u=\partial_xh$，得到随机Burgers方程$\partial_tu=\nu\partial_x^2u+\lambda u\partial_xu+\partial_x\eta$，高度不守恒而斜率总积分在周期边界守恒。另一变换$Z=\exp[\lambda h/(2\nu)]$把非线性梯度抵消，形式上$\partial_tZ=\nu\nabla^2Z+\lambda\eta Z/(2\nu)$，是乘性随机热方程。白噪声下需固定Itô/Stratonovich与正规化常数，不能把光滑链式法则直接当作无紫外问题的精确随机恒等式。

随机热方程的路径积分可解释为随机环境中的有向聚合物配分，$h\propto\log Z$是其自由能。这连接非平衡增长与无序多路径问题；下一章讨论无序平均为何不能只把随机势换成均值。

\originalchapter{淬火无序、聚合物与大纲补充}
\label{ch:random}
\source{8.334大纲列随机体系与聚合物；26份公开讲义未为这些大纲项各给独立全讲，本章为编者补写的基础连接，不冒称遗漏的课堂原稿。}
\originalsection{退火与淬火的区别}
若杂质随系统充分运动，可与热变量一起求和，退火自由能$F_a=-\kb T\log\overline Z$；若杂质在观测时间内固定，每个样本先取热自由能，淬火$F_q=-\kb T\,\overline{\log Z}$。由$\log$凹性$\overline{\log Z}\le\log\overline Z$，故$F_q\ge F_a$。平均哈密顿量、平均配分函数、平均自由能不是同一运算，不能互换。

例如独立随机能级$E$的单态$Z=\e^{-\beta E}$，淬火$F_q=\overline E$，退火$F_a=\overline E-\beta\operatorname{Var}E/2$对Gaussian成立。退火结果被罕见低能样本偏重，差异不是仅一个数值误差。
\originalsection{Harris与Imry--Ma判据}
短程独立随机局域临界温度在相关体积$\xi^d$内的平均涨落$\delta t_\xi\sim\xi^{-d/2}$。洁净体系距离临界$t\sim\xi^{-1/\nu}$。若随机涨落比$t$更小，要求$d\nu>2$；反之$d\nu<2$无序相关。边际等号需更细RG，不能据此判定必然相关或无关。若洁净超尺度成立，等价$\alpha<0$使弱随机质量无关。

随机场耦合$-\int h(x)\phi(x)$则不同。Ising翻一个尺度$L$畴的壁代价$\sim\sigma L^{d-1}$，随机场能增益$\sim hL^{d/2}$。$d<2$时大畴增益压过壁，低维长程序被破坏；$d=2$边际需更精细分析。连续自旋可通过缓慢转向降低代价为$L^{d-2}$，比较给下临界尺度4。随机键与随机场不能用同一个Harris不等式混称。
\originalsection{副本与自旋玻璃的范围}
副本恒等式形式$\overline{\log Z}=\lim_{n\to0}(\overline{Z^n}-1)/n$，对正整数$n$可先计算$n$份热变量共享同一无序。若随机场Gaussian，平均$\exp[\beta\int h\sum_a\phi_a]$生成$\frac12\beta^2\Delta\int(\sum_a\phi_a)^2$，使不同副本耦合。随机质量则生成四次副本耦合。这解释副本作用量项的来源。

从整数$n$解析继续到0并非无需证明的普遍合法步骤；副本对称鞍点也可能失稳。自旋玻璃以无序冻结局域磁矩与大量亚稳态为特征，可用Edwards--Anderson量$q_{EA}=\overline{\avg{s_i}^2}$区别均匀磁化为0的冻结态。完整有限维自旋玻璃与副本破缺的严谨解不属于这组公开讲义的可重建证明，本章给定义、平均结构与判据，不以一个$q$方程宣称完成所有模型。
\originalsection{随机链与Flory尺度}
理想柔性链有$N$段、段长$a$，独立方向使$\avg{R^2}=Na^2$，Gaussian端点分布$P(R)\propto\exp[-dR^2/(2Na^2)]$。排斥体积$ v>0$可近似自由能$F/\kb T\sim R^2/(Na^2)+vN^2/R^d$，前项是伸展熵代价，后项为单体密度平方积分。极小$R^{d+2}\sim v a^2N^3$，Flory指数$\nu_F=3/(d+2)$，三维为$3/5$。

这是量级平衡而非精确指数；它忽略链段相关，二维虽恰得$3/4$也不能据此称该方法严格。$d>4$排斥体积无关，应回理想$\nu=1/2$，而不是无条件继续$3/(d+2)$。$n\to0$的$O(n)$场论把两点高温图中一条开链保留、闭环因$n$权重消失，建立自避行走与RG的联系；近四维前章$\nu=1/2+\epsilon/16+O(\epsilon^2)$给可控修正。

吸引与排斥抵消的$\theta$点涉及三体稳定项与三临界，上临界维数3。更复杂网络需要按节点功能与连接拓扑另立熵计数，不能把所有网络都视为单条线性链。课程大纲的这些扩展保留实际覆盖与证明边界。

\volpart{第三编　8.333公开作业与历史复习题}
\originalchapter{8.333作业1：热力学}
\psource{Mehran Kardar，8.333 Fall 2013，PS1，截止2013年9月18日；与阶段复习1的同题另作对应。}
\begin{exercise}[原题1(a)--(d)]\label{p:333-1-1}
弹性丝位移$x$所需张力$J=ax-bT+cTx$，定$x$热容$C_x=A(x)T$。(a)由Maxwell关系求$S_x|_T$。(b)证明$A$不依$x$。(c)取$S(0,0)=S_0$求熵。(d)以$T,J$求定张力热容$C_J$，注明稳定区。
\end{exercise}
\begin{solution}
$\dd E=T\dd S+J\dd x$，$\dd F=-S\dd T+J\dd x$，所以$S_x=-J_T=b-cx$。而$S_T=C_x/T=A(x)$，混合偏导要求$A'(x)=\partial_T(b-cx)=0$。积分$S=S_0+AT+bx-cx^2/2$。定$J$时$x=(J+bT)/(a+cT)$，$\partial_Tx|_J=(ab-cJ)/(a+cT)^2$，故
\begin{equation}
C_J=AT+\frac{T(ab-cJ)^2}{(a+cT)^3}.
\end{equation}
局部机械稳定需$\partial_xJ=a+cT>0$，通常热稳定还需$A>0$。有限温经验方程不必适用于所有$T$，$T=0$积分常数是题给的数学延拓条件。
\end{solution}
\begin{exercise}[原题2(a)--(d)]\label{p:333-1-2}
硬芯气满足$p(V-Nb)=N\kb T$，固定$N$、$C_V$不随温度。(a)求涉及$S_V|_T$的Maxwell关系。(b)求$\dd E$并证明$E$不依$V$。(c)证明$\gamma=C_p/C_V=1+N\kb/C_V$。(d)证明绝热$p(V-Nb)^\gamma$恒定。
\end{exercise}
\begin{solution}
$S_V=p_T=N\kb/(V-Nb)$，$\dd E=C_V\dd T+[Tp_T-p]\dd V=C_V\dd T$，所以$E=C_VT+E_0(N)$。定压膨胀$V-Nb=N\kb T/p$，焓导数$C_p=C_V+N\kb$。熵$S=C_V\log T+N\kb\log(V-Nb)+S_0(N)$，绝热$T(V-Nb)^{N\kb/C_V}$恒定，乘状态方程即所求。要求$V>Nb,T>0$。
\end{solution}
\begin{exercise}[原题3(a)--(e)]\label{p:333-1-3}
零场超导与正常金属热容$C_s=V\alpha T^3,C_n=V(\beta T^3+\gamma T)$，体积变化忽略。(a)第三定律求两熵。(b)零场转变无潜热，求$T_c$。(c)$E_n(0)=E_0,E_s(0)=E_0-V\Delta$，求有限温能量。(d)相平衡求$\Delta$。(e)采用Gaussian磁功$\dd E=T\dd S+B\dd M+\mu\dd N$，$M_s=-VB/(4\pi),M_n=0$，求使超导失稳的$B_c(T)=B_0[1-T^2/T_c^2]$与$B_0$。
\end{exercise}
\begin{solution}
熵$S_s=V\alpha T^3/3,S_n=V(\beta T^3/3+\gamma T)$，无潜热要求$S_s(T_c)=S_n(T_c)$，得$T_c^2=3\gamma/(\alpha-\beta)$，需$\alpha>\beta,\gamma>0$。能量$E_s=E_0-V\Delta+V\alpha T^4/4$、$E_n=E_0+V\beta T^4/4+V\gamma T^2/2$。自由能差
\begin{equation}
F_n-F_s=V[\Delta+(\alpha-\beta)T^4/12-\gamma T^2/2].
\end{equation}
$T_c$相等给$\Delta=3\gamma^2/[4(\alpha-\beta)]=\gamma T_c^2/4$，所以差$V\Delta(1-T^2/T_c^2)^2$。磁场自然势$G=E-TS-BM$有$G_B=-M$，超导额外$VB^2/(8\pi)$、正常不变。平衡给$B_c^2=8\pi\Delta(1-T^2/T_c^2)^2$，故$B_0=\sqrt{8\pi\Delta}$，此形式只用于$T\le T_c$与题示热容模型。
\end{solution}
\begin{exercise}[原题4(a)--(e)]\label{p:333-1-4}
用光子气的无穷小Carnot循环，从相邻温度$T,T+\dd T$、压强$p,p+\dd p$与体积区间$V,V+\dd V$推黑体能量。(a)功$W$。(b)等温膨胀吸热$Q$。(c)用Carnot效率联系二者。(d)已知$p=AT^4,A=\pi^2\kb^4/[45(\hbar c)^3]$、$E(T,0)=0$求$E$。(e)求绝热线。
\end{exercise}
\begin{solution}
循环面积$W=\dd p\dd V$（取发动机方向），$Q=[(\partial_VE)_T+p]\dd V$，效率$\dd T/T$给$(\partial_VE)_T+p=T\dd p/\dd T$。代$p=AT^4$得$E=3AT^4V$。$S=(E+pV)/T=4AT^3V$，绝热$VT^3$恒定，亦$pV^{4/3}$恒定。原图的相邻曲线信息已由这些状态条件完整表述，不把示意图面积当有限循环精确公式。
\end{solution}
\begin{exercise}[原题5(a)--(b)]\label{p:333-1-5}
(a)两个初温不同的物体只通过热交换趋向共同温度，利用热流方向证明总熵增加。(b)可调体积气体置于恒$T,p$热库，热库熵变$-(\Delta E+p\Delta V)/T$，证明平衡极小$G=E+pV-TS$。
\end{exercise}
\begin{solution}
取$\delta Q>0$从较热物体1流到较冷2，均匀体的熵变化$\delta S=\delta Q(1/T_2-1/T_1)>0$，积分给正总增量；若物体内部有不可逆温差还可增加更多熵。(b)$\Delta S_{\rm tot}=\Delta S_{\rm gas}-(\Delta E+p\Delta V)/T=-\Delta G/T\ge0$，所以固定热库条件下自发过程$G$下降，稳定平衡取可达状态中最低$G$。这是组合体系熵原理，并非开放气体自身熵必增。
\end{solution}
\begin{exercise}[原题6(a)--(b)]\label{p:333-1-6}
两物体热容常数$C_1,C_2$、$T_1>T_2>0$。(a)直接接触且无外界交换，求最终温度与熵变。(b)用可逆Carnot发动机在二者间传热，求最终温度与输出功。
\end{exercise}
\begin{solution}
(a)能量守恒给$T_F=(C_1T_1+C_2T_2)/(C_1+C_2)$，$\Delta S=C_1\log(T_F/T_1)+C_2\log(T_F/T_2)\ge0$，正号由加权对数凹性。(b)总熵守恒给$T_R=T_1^{C_1/(C_1+C_2)}T_2^{C_2/(C_1+C_2)}$，输出$W=C_1T_1+C_2T_2-(C_1+C_2)T_R$。$T_R\le T_F$来自算术--几何均值；更多初始内能可作为功输出。
\end{solution}
\begin{exercise}[原题7]\label{p:333-1-7}
发动机与一组$T_{max}=T_1>\cdots>T_n=T_{min}$热源交换，吸入$Q_\alpha^+>0$、放出$Q_\beta^->0$，$W=Q^+-Q^-$。证明效率不超过仅用最高最低温的Carnot效率，说明严格小于的条件。
\end{exercise}
\begin{solution}
周期热源熵要求$\sum Q_\alpha^+/T_\alpha^+\le\sum Q_\beta^-/T_\beta^-$。左侧$\ge Q^+/T_{max}$，右侧$\le Q^-/T_{min}$，所以$\eta=1-Q^-/Q^+\le1-T_{min}/T_{max}$。非零热量在中间温度吸入或释放、或循环不可逆时不等号严格；若只在两端交换且可逆，可等号。原题“less”应按实际多温流量理解，不能排除经典Carnot等号。
\end{solution}
\begin{exercise}[原题8(a)--(c)，选做]\label{p:333-1-8}
近隔离稀薄气云经引力形成太阳和行星。(a)轨道和结构更有序为何不违反第二定律？(b)太阳核反应从质子形成较重元素是否因此降低总熵？(c)地球生命与智能的组织增长如何与第二定律相容？
\end{exercise}
\begin{solution}
(a)可见空间结构并非总微观熵。引力收缩把势能释放为运动、热与辐射，粒子与辐射的状态数可增长；自引力系统的长程非加性还使普通气体直觉的均匀平衡不直接适用。(b)核组成的某些自由度减少不决定总熵，必须把释能、光子及其他产物和传输一起计入，恒定粒子种类的简单混合熵不覆盖核反应。(c)地球接受高温太阳光并向较冷空间辐射，吸能与排能的熵不同，开放子系统可形成局部组织同时向环境排更多熵。第二律约束包括热源与辐射环境的总熵，不以“复杂程度”替代$S$。
\end{solution}

\originalchapter{8.333作业2：概率与涨落}
\psource{Mehran Kardar，8.333 Fall 2013，PS2。}
\begin{exercise}[原题1(a)--(b)]\label{p:333-2-1}
独立各列沉积金原子，平均每秒$d$层。(a)时间$t$后某位置有$m$原子的概率、未覆盖比例。(b)厚度方差。
\end{exercise}
\begin{solution}
若各小时间段独立且单次到达率$d$，Poisson参数$\lambda=dt$，$P(m)=\e^{-\lambda}\lambda^m/m!$，未覆盖比例$\e^{-dt}$，均值与方差皆$dt$；原子层实际高度为$a$时厚度方差$a^2dt$。只给平均率却不假定独立Poisson到达不能唯一确定分布，题目的独立沉积模型在此明确采用。
\end{solution}
\begin{exercise}[原题2(a)--(d)]\label{p:333-2-2}
二维链$N$段$t_i$、$|t_i|=a$，能$H=-\kappa\sum t_i\cdot t_{i+1}=-\kappa a^2\sum\cos\phi_i$，$R=\sum t_i$。(a)求$\avg{t_m\cdot t_n}$指数相关及持久长。(b)$N\gg1$求$\avg{R^2}$。(c)求大$N$端点分布。(d)拉力$F$使权重乘$\e^{\beta F\cdot R}$，求$\avg R=K^{-1}F+O(F^3)$与弹簧常数。
\end{exercise}
\begin{solution}
相邻转角独立且密度$\propto\e^{u\cos\phi}$、$u=\beta\kappa a^2$。$c=\avg{\cos\phi}=I_1(u)/I_0(u)$，$\avg{\sin\phi}=0$，迭代给$\avg{t_m\cdot t_n}=a^2c^{|m-n|}$，$\xi=-1/\log c,\ell_p=a\xi$（$\kappa>0$；反向偏好需含振荡符号）。双和
\begin{equation}
\avg{R^2}=a^2\left[N+2\sum_{s=1}^{N-1}(N-s)c^s\right]\simeq Na^2\frac{1+c}{1-c}
\end{equation}
要求$N\gg\xi$，不仅$N\gg1$。中央极限给$P(R)\simeq[\pi\avg{R^2}]^{-1}\e^{-R^2/\avg{R^2}}$，两坐标方差各一半。对拉力展开配分，$\avg{R_i}=\beta\avg{R_iR_j}F_j=\beta\avg{R^2}F_i/2$，$K=2\kb T/\avg{R^2}$。大拉力伸展到$aN$附近Gaussian与线性Hooke律失效。
\end{solution}
\begin{exercise}[原题3(a)--(b)]\label{p:333-2-3}
二维随机觅食从巢原点扩散，系数$D$。(a)离巢$t$时距$r$分布。(b)每段寻找时间独立指数平均$\tau$，返回与停留可忽略，长时间后的空间分布及大$r$鞍点渐近。
\end{exercise}
\begin{solution}
向量密度$P(\vct r,t)=(4\pi Dt)^{-1}\e^{-r^2/(4Dt)}$，径向密度$p_r(r,t)=2\pi rP$，二者的单位不同。指数重置的平稳年龄密度$\tau^{-1}\e^{-t/\tau}$，所以
\begin{equation}
P_\infty(r)=\frac1{4\pi D\tau}\int_0^\infty\frac{\dd t}{t}\e^{-t/\tau-r^2/(4Dt)}=\frac1{2\pi D\tau}K_0(r/\sqrt{D\tau}).
\end{equation}
指数的鞍点$t_*=r\sqrt{\tau/(4D)}$，作用$r/\sqrt{D\tau}$，Gaussian宽度给$K_0(z)\sim\sqrt{\pi/(2z)}\e^{-z}$。向量密度$\propto r^{-1/2}\e^{-r/\sqrt{D\tau}}$，径向密度$\propto r^{1/2}\e^{-r/\sqrt{D\tau}}$。原点的对数峰可积，不是有限概率点质量，因为停巢时间被取零。
\end{solution}
\begin{exercise}[原题4(a)--(b)]\label{p:333-2-4}
(a)由凸函数切线下界证明Jensen不等式。(b)以它证明$D(p\Vert q)=\int p\log(p/q)\dd x\ge0$，说明零概率条件。
\end{exercise}
\begin{solution}
在$y=\avg X$取$f(x)\ge f(y)+f'(y)(x-y)$并平均，线性项0，故$\avg{f(X)}\ge f(\avg X)$；非可微凸函数用支撑线同理。取$f=-\log$、随机量$q/p$且按$p$平均，$D\ge-\log\int_{p>0}q\ge0$。若$p>0$而$q=0$，相对熵$+\infty$；若都是归一且支撑相同、函数严格凸，等号要求$q/p$恒定即$p=q$。
\end{solution}
\begin{exercise}[原题5(a)--(e)，(e)选做]\label{p:333-2-5}
IID序列$x_n$，$x_n$大于此前全部值时记录指标$R_n=1$，$R_1=1$。(a)求记录率与分布形式的关系。(b)记录总数$S_N$期望及大$N$，每年尝试数翻倍时增长。(c)证明不同记录指标协方差0。(d)给所有矩与前三累积量，判断中心极限。(e)已知第三记录$n_3=9$，求第二记录$n_2$与第四$n_4$期望。采用连续分布或随机打破平局；有原子概率时原题的分布无关结论须修改。
\end{exercise}
\begin{solution}
无平局时前$n$个次序均等，每个位置成为最大概率$1/n$，所以$P(R_n=1)=1/n$。$\avg S_N=H_N\sim\log N+\gamma_E$，若$N=N_02^t$，记录均数约$t\log2+\log N_0+\gamma_E$。随机排列的“新元素在前$n$中排名”依次独立且均匀，故记录指标独立，特别$\avg{R_mR_n}=1/(mn)$。

矩母函数$M(t)=\prod_{n=1}^N[1+(\e^t-1)/n]=\Gamma(N+\e^t)/[\Gamma(\e^t)N!]$，所有$r$阶矩由$M^{(r)}(0)$给出，也可写$\avg{S_N^r}=\sum_k k^r c(N,k)/N!$，$c(N,k)$为第一类无符号Stirling数。前三累积量$\kappa_1=H_N$、$\kappa_2=H_N-H_N^{(2)}$、$\kappa_3=H_N-3H_N^{(2)}+2H_N^{(3)}$。独立非同分布的有界指标满足Lindeberg条件，方差$\sim\log N\to\infty$，故标准化后趋Gaussian，不能说IID中心极限直接适用于原$R_n$。

(e)条件$n_3=9$要求2--8只有一个记录，位置$r$相对权重$[1/r]/[1-1/r]=1/(r-1)$，所以$P(n_2=r|n_3=9)=1/[(r-1)H_7]$、$\avg{n_2}=1+7/H_7\simeq3.6997$。已出现9记录后，下一记录在$m>9$的概率$9/[m(m-1)]$，均值$9\sum_{m>9}1/(m-1)=+\infty$。中位等待有限却均值发散，不能把$\avg{n_4}$猜成18。
\end{solution}
\begin{exercise}[原题6(a)--(d)]\label{p:333-2-6}
系统初始正则平衡，按协议做外功$W$，Hamilton微观态经一一且体积保持的演化到新$H'$；正反协议初态各为对应平衡。(a)证明$P_f(W)/P_b(-W)=\e^{\beta(W-\Delta F)}$。(b)推Jarzynski$\avg{\e^{-\beta W}}=\e^{-\beta\Delta F}$。(c)推$\avg W\ge\Delta F$。(d)循环若净抽出功$\omega>0$，给违例概率上界。题页净功符号文字有歧义，此处规定$W$为输入系统的功。
\end{exercise}
\begin{solution}
对应路径的初态权重比$[\e^{-\beta H}/Z]/[\e^{-\beta H'}/Z']=\e^{\beta(W-\Delta F)}$。时间反演给反路径，Liouville的单位Jacobian使按相同$W$积分后仍得(a)；仅“确定性”不保证Jacobian1，题中需Hamilton或等价微观可逆过程。

乘$P_f\e^{-\beta W}=P_b(-W)\e^{-\beta\Delta F}$积分得(b)。Jensen给$\e^{-\beta\avg W}\le\avg\e^{-\beta W}$，得(c)。循环$\Delta F=0$，以非负变量$Y=\e^{-\beta W}$、$\avg Y=1$用Markov不等式，$P(W\le-\omega)=P(Y\ge\e^{\beta\omega})\le\e^{-\beta\omega}$。稀有轨迹可抽功，平均仍服从第二律；用有限样本估计指数均值时罕见事件的重要性不可忽略。
\end{solution}
\begin{exercise}[原题7(a)--(b)，选做]\label{p:333-2-7}
仅知骰子6面概率为1面的两倍。(a)求除此约束外最大熵的“无偏”概率。(b)这条信息相对公平骰子减少多少比特熵？
\end{exercise}
\begin{solution}
令约束$p_6-2p_1=0$，最大$-\sum p_i\log p_i$的乘子给$p_1=q2^{-2/3},p_6=q2^{1/3},p_2=\cdots=p_5=q$，$q=[4+2^{-2/3}+2^{1/3}]^{-1}$。若将“其他面与1相同”另作假设会得1/7与2/7，但那是新增约束，不是本题最大熵。记$Z=1/q$，约束使熵中乘子项相消，$H=\log_2Z$，信息量$I=\log_26-H=\log_2(6/Z)$，约0.0267 bit。
\end{solution}
\begin{exercise}[原题8(a)--(d)，选做]\label{p:333-2-8}
平衡涨落变量$x_i$，$C_{ij}(t)=\avg{x_i(t)x_j(0)}$。(a)用时间平移与反演得互易关系。(b)若分布$P\propto\exp[-x^TKx/2]$、$K$正定，求$C(0)$。(c)共轭力$J=-\nabla\log P$证明$\avg{J_ix_j}=\delta_{ij}$。(d)平均松弛$\dot x=-\mu Kx$，由$\dot C(0)$推出Onsager$\mu$对称，说明变量的时间反演奇偶条件。
\end{exercise}
\begin{solution}
平稳是同时平移两时刻，$C_{ij}(t)=C_{ji}(-t)$；不是把时间差任意移仍相同。若$x_i$皆时间反演偶，反演又给$C_{ji}(-t)=C_{ji}(t)$。Gaussian完成平方得$C(0)=K^{-1}$，$J_i=K_{ij}x_j$，所以(c)$KC=I$。$\dot C(0^+)=-\mu KC(0)=-\mu$，互易$C(t)=C(t)^T$令$\mu=\mu^T$。奇变量时应为Onsager--Casimir带奇偶符号关系，不能强行对称。真正平稳随机松弛还需FDT噪声；所给确定性方程描述平均线性响应，噪声对这个右导数与过去初态无相关贡献。
\end{solution}

\originalchapter{8.333作业3：动力学}
\psource{Mehran Kardar，8.333 Fall 2013，PS3，截止2013年10月18日。}
\begin{exercise}[原题1(a)--(b)]\label{p:333-3-1}
(a)证明沿Hamilton轨道$\dot O=\{O,H\}$。(b)若任意可测光滑观测$O$的系综平均$\avg{\{O,H\}}=0$，证明$\{H,\rho\}=0$，注明边界条件。
\end{exercise}
\begin{solution}
链式法则与Hamilton方程给(a)。对任意紧支撑$O$或有消失表面项的周期/衰减边界，分部积分$\int\rho\{O,H\}=\int O\{H,\rho\}$。左侧全为0，测试函数任意使$\{H,\rho\}=0$按分布意义；若$\rho$可微则逐点成立。只对一个$O$或有限几个矩为0不够推出全密度平稳。
\end{solution}
\begin{exercise}[原题2(a)--(f)]\label{p:333-3-2}
局域三维Maxwell分布$f_1^0=n(2\pi m\kb T)^{-3/2}\exp[-(p-mu)^2/(2m\kb T)]$。(a)求压力张量$P_{ij}=mn\avg{c_ic_j}$与热流$h_i=nm\avg{c_ic^2}/2$，$c=p/m-u$。(b)代入粒子、动量与能量守恒求$n,u,T$的零阶流体方程。(c)证明$D_t\log(nT^{-3/2})=0$。(d)积分动量写$H_0=\int f_1^0\log f_1^0$。(e)求$\dot H_0$。(f)解释平衡趋近。
\end{exercise}
\begin{solution}
Gaussian二阶矩给$P_{ij}=n\kb T\delta_{ij}$，奇矩$h_i=0$。故$D_tn=-n\nabla\cdot u$、$mnD_tu=-\nabla(n\kb T)$、$D_tT=-(2/3)T\nabla\cdot u$。前后两对数导数相消得(c)。积分$\avg{mc^2/(2\kb T)}=3/2$，
\begin{equation}
H_0=\int\dd^3r\,n\left[\log\frac n{(2\pi m\kb T)^{3/2}}-\frac32\right].
\end{equation}
常相空间单位仅加与$N$有关常数。由(c)方括号沿流线守恒，再以连续性求$\partial_t(n\mathcal A)+\nabla\cdot(nu\mathcal A)=0$，无边界流时$\dot H_0=0$。零阶Euler不产生耗散，局域Maxwell但非均匀态可持续运动；趋平衡需要梯度修正的黏性、热导或其他不可逆机制。
\end{solution}
\begin{exercise}[原题3(a)--(c)]\label{p:333-3-3}
两板距$w$，下板静止、上板沿$x$速度$u$。零阶分布为均匀$n,T$、局域流速$u_x=\alpha y$的Maxwell，$\alpha=u/w$。(a)用单碰撞时间$\tau_\times$求一阶$f_1=f_0-\tau_\times(v\cdot\nabla)f_0$。(b)求穿过$y$面$x$动量流$\Pi_{xy}$。(c)求板受力与黏度。
\end{exercise}
\begin{solution}
令$c_x=p_x/m-\alpha y,c_y=p_y/m$，$\partial_yf_0=(m\alpha c_x/\kb T)f_0$，所以$\delta f=-\tau_\times m\alpha c_xc_yf_0/(\kb T)$。$\Pi_{xy}=\int mc_xc_y\delta f\dd p=-n\kb T\tau_\times\alpha$，因$\avg{c_x^2c_y^2}=(\kb T/m)^2$。下板所受沿流方向的力密度$F_x=-\Pi_{xy}$，上板反向，$\eta=n\kb T\tau_\times$，与$y$无关。要求$|\alpha|\tau_\times\ll1$且板间有流体窗口$w\gg\ell$。
\end{solution}
\begin{exercise}[原题4(a)--(e)]\label{p:333-3-4}
固定两态原子$a_0,a_1$相差$\varepsilon$，与$\omega=\varepsilon/\hbar$光子作用。自发$a_1\to a_0+\gamma$率$\sigma_{sp}$，吸收截面$\sigma_{ad}$、受激发射$\sigma_{st}$均角无关。(a)求$n_1/n_0$。(b)写光子密度$f(r,p,t)$动力学。(c)求$f_{eq}$。(d)用$ f_{eq}=h^{-3}/(\e^{\beta\hbar\omega}-1)$求三截面关系。(e)介质占$x>0$，沿$+x$入射光的非平衡通量如何衰减？求穿透长。
\end{exercise}
\begin{solution}
(a)等简并两态$n_1/n_0=\e^{-\beta\varepsilon}$。选原题相空间密度归一，自发源每态为$n_1\sigma_{sp}$，吸收与受激的速度/几何常数吸入其率系数，$\partial_tf+c\hat p\cdot\nabla f=n_1\sigma_{sp}+(n_1\sigma_{st}-n_0\sigma_{ad})f$。稳态$ f_{eq}=n_1\sigma_{sp}/(n_0\sigma_{ad}-n_1\sigma_{st})=\sigma_{sp}/[\e^{\beta\varepsilon}\sigma_{ad}-\sigma_{st}]$。

(d)对所有温度匹配Planck式，需$\sigma_{ad}=\sigma_{st}$且$\sigma_{sp}=\sigma_{ad}/h^3$于上述密度归一；若截面用单位实际光通量，$c$与态密度常数须恢复，两Einstein B系数相等的内容不变。(e)沿入射方向$c\partial_x(f-f_{eq})=-(n_0\sigma_{ad}-n_1\sigma_{st})(f-f_{eq})$，所以穿透长$\ell_\gamma=c/(n_0\sigma_{ad}-n_1\sigma_{st})$，正温平衡$n_0>n_1$使衰减。若$n_1>n_0$人口反转，系数负，变增益而非正穿透长。
\end{solution}
\begin{exercise}[原题5(a)--(b)]\label{p:333-3-5}
外势$U(q)$中经典气体。(a)在归一与固定平均总能下极小一体$H=N\int\rho_1\log\rho_1\dd p\dd q$，求平衡密度。(b)两种质量$m_a,m_b$的气体共享固定总能，求两密度并说明每粒子平动能可定义共同温度。原题声称平稳解的必要性须检查其他守恒量。
\end{exercise}
\begin{solution}
(a)乘子变分$\log\rho_1+1+\lambda+\beta[p^2/(2m)+U]=0$，得$\rho_1=Z_1^{-1}\e^{-\beta[p^2/(2m)+U]}$。严格凸$\rho\log\rho$使给定这些约束下极小唯一。(b)各组归一乘子不同，能量乘子$\beta$相同，所以$\rho_\alpha\propto\e^{-\beta[p^2/(2m_\alpha)+U_\alpha]}$，每粒子平动能$3/(2\beta)=3\kb T/2$相同。质量改变速度方差，不改共同能量温标。

H定理的零碰撞态可含流速；在周期系统守恒总动量或有整体旋转等额外约束时，平稳态还可为广义canonical分布。题目所列零流唯一必要性只在无额外宏观守恒、固定静止边界或相应约束已排除时成立，不能将熵极小解当任意平稳解的无条件必要式。
\end{solution}
\begin{exercise}[原题6(a)--(c)，选做]\label{p:333-3-6}
经典电子气被相对真空深度$\varphi$的势阱束缚。(a)平坦表面逃逸电子内外速度关系。(b)温度$T$的速度分布。(c)密度$n$时热发射电流密度。
\end{exercise}
\begin{solution}
切向速度守恒。法向速度满足$v_z'^2=v_z^2-2\varphi/m$，因此仅$v_z>\sqrt{2\varphi/m}$的电子向外逃逸。设Maxwell速度密度为
\begin{equation}
P(v)=\left(\frac{m}{2\pi\kb T}\right)^{3/2}\e^{-mv^2/(2\kb T)}.
\end{equation}
电流大小由通量积分求出：
\begin{equation}
j=en\int_{v_z>v_0}v_zP(v)\dd^3v
=en\sqrt{\frac{\kb T}{2\pi m}}\e^{-\varphi/\kb T}.
\end{equation}
用入射到表面的速度做通量积分，不能用外部变小的速度却不换密度Jacobian。题目是经典模型；真实简并金属Richardson--Dushman前因子需量子态填充。
\end{solution}

\originalchapter{8.333作业4：非相互作用粒子}
\psource{Mehran Kardar，8.333 Fall 2013，PS4，截止2013年11月13日。}
\begin{exercise}[原题1(a)--(d)]\label{p:333-4-1}
$N$个相同经典原子在球形谐陷阱$H=\sum[p_i^2/(2m)+Kr_i^2/2]$。(a)证明每粒子角动量守恒，说明加入旋转对称相互作用后只保证总角动量。(b)广义分布原题取$\exp[-\beta(H+\Omega L_z)]$，$|\Omega|<\sqrt{K/m}$，求$Z$。(c)求$\avg{L_z}$。(d)求位置密度与$\avg{x^2},\avg{y^2},\avg{z^2}$。
\end{exercise}
\begin{solution}
中心力使$\dot L=r\times(-Kr)=0$。相互作用两体中心力的转矩对总量相消，却在粒子之间交换角动量。完成平方
\begin{equation}
H_1+\Omega L_z=\frac{(p_x-m\Omega y)^2+(p_y+m\Omega x)^2+p_z^2}{2m}+\frac{K-m\Omega^2}{2}(x^2+y^2)+\frac K2z^2.
\end{equation}
令$\omega_0=\sqrt{K/m}$，$Z=Z_1^N/N!$，$Z_1=(\kb T/\hbar\omega_0)^3/[1-\Omega^2/\omega_0^2]$。原题加号给$\avg{L_z}=-\beta^{-1}\partial_\Omega\log Z=-2Nm\Omega\kb T/(K-m\Omega^2)$；通常转动热库用$H-\Omega L$时符号反转，不能暗改题。

位置Gaussian密度$P(r)=[\beta/(2\pi)]^{3/2}(K-m\Omega^2)\sqrt K\exp\{-\beta[(K-m\Omega^2)(x^2+y^2)+Kz^2]/2\}$，方差横向$\kb T/(K-m\Omega^2)$、纵向$\kb T/K$。$|\Omega|\to\omega_0$失去横向约束，配分函数发散，是稳定条件的真实边界。
\end{solution}
\begin{exercise}[原题2(a)--(e)]\label{p:333-4-2}
$N$个双原子分子固定在方格金属表面，每个可沿$x,y$平躺（能0）或沿$z$站立（能$\varepsilon>0$）。(a)基态数与最高微观能。(b)能量$E$的态数及熵。(c)热容并画温度变化。(d)特定分子站立概率。(e)正温内能的最大值。
\end{exercise}
\begin{solution}
(a)基态数$2^N$，最高$N\varepsilon$的全部站立态只有1种。(b)$M=E/\varepsilon$须为0至$N$整数，$\Omega=\binom NM2^{N-M}$、$S=\kb\log\Omega$；大$N$以Stirling给$S/(N\kb)=-x\log x-(1-x)\log[(1-x)/2]$，$x=M/N$。

canonical单站$z=2+\e^{-\beta\varepsilon}$，站立$q=1/(1+2\e^{\beta\varepsilon})$，$E=N\varepsilon q$，$C=N\kb(\beta\varepsilon)^2q(1-q)$。两端$T\to0$指数小，$T\to\infty$按$T^{-2}$小，中间有Schottky峰。以下为$C/(N\kb)$按$k_BT/\varepsilon$的矢量曲线。
\begin{center}\begin{tikzpicture}\begin{axis}[width=10cm,height=4.6cm,xlabel={$\kb T/\varepsilon$},ylabel={$C/(N\kb)$},xmin=0,xmax=3,ymin=0]
\addplot[structurecolor,thick,domain=.06:3,samples=180]{(1/x)^2*2*exp(1/x)/(1+2*exp(1/x))^2};\end{axis}\end{tikzpicture}\end{center}
正温$q<1/3$，最大热能的上确界$N\varepsilon/3$；高于此能的微正则宏观态可对应负温，不是“最高微观能”等于正温均能。
\end{solution}
\begin{exercise}[原题3(a)--(d)]\label{p:333-4-3}
$N$个量子自旋$s$，磁矩$m_z=\mu m,m=-s,\ldots,s$，外場$B$。(a)Gibbs配分。(b)小$B$自由能至$B^2$。(c)Curie磁化率。(d)小场证明$C_B-C_M=cB^2/T^2$，$\chi=c/T$。
\end{exercise}
\begin{solution}
$x=\beta\mu B$，$Z=[\sum_me^{xm}]^N=[\sinh((2s+1)x/2)/\sinh(x/2)]^N$，$G=-N\kb T\log z$。零場均匀$m$有$\avg{m^2}=s(s+1)/3$，所以$G(B)=G(0)-N\mu^2s(s+1)B^2/(6\kb T)+O(B^4)$，$\chi=N\mu^2s(s+1)/(3\kb T)=c/T$。由$\dd G=-S\dd T-M\dd B$的Maxwell关系，$C_B-C_M=T(\partial_TM|_B)^2/(\partial_BM|_T)$，代$M=cB/T+O(B^3)$给$cB^2/T^2+O(B^4)$。题示式是Curie小场阶，并非整个Brillouin函数的有限场精确关系。
\end{solution}
\begin{exercise}[原题4(a)--(f)]\label{p:333-4-4}
理想气与有$N_s$个单占据吸附位置的催化表面接触，每吸附粒子能$-\varepsilon$。(a)求气体$\mu(T,p)$。(b)表面巨配分。(c)占据率$\theta=p/[p+p_0(T)]$及$p_0$。(d)粒子数特征函数与全部累积量相对巨势的导数。(e)证明数涨落响应。(f)求相对涨落。
\end{exercise}
\begin{solution}
$\mu=\kb T\log[p\lambda_T^3/(\kb T)]=\kb T[\log(p/T^{5/2})+A_0]$，后式的常数包含使对数无量纲的单位。表面$\Xi_s=(1+z\e^{\beta\varepsilon})^{N_s}$，$\theta=z\e^{\beta\varepsilon}/(1+z\e^{\beta\varepsilon})$，所以$p_0=(\kb T/\lambda_T^3)\e^{-\beta\varepsilon}$。

$\avg{\e^{-\ii kN}}=\Xi_s(\mu-\ii\kb Tk)/\Xi_s(\mu)$，巨势$\mathcal G=-\kb T\log\Xi_s$，$\kappa_m=-(\kb T)^{m-1}\partial_\mu^m\mathcal G$，固定$T$求导。$\operatorname{Var}N=\kb T\partial_\mu\avg N=N_s\theta(1-\theta)$，除$\avg N^2=(N_s\theta)^2$得$(1-\theta)/(N_s\theta)$。题中固定位置数与随机粒子数都写$N$易混，本解用$N_s$区分。
\end{solution}
\begin{exercise}[原题5(a)--(d)]\label{p:333-4-5}
O$_2$分子自旋1、$S_z=-1,0,1$，忽略其他内部自由度，$H=\sum[p_i^2/(2m)-\mu BS_{zi}]$，体积$V$，平动经典。(a)配分。(b)各自旋概率。(c)磁化每体积。(d)零场总磁化率。
\end{exercise}
\begin{solution}
$x=\beta\mu B$，$Z=(V/\lambda_T^3)^N(1+2\cosh x)^N/N!$，$P_m=\e^{xm}/(1+2\cosh x)$。$M/V=(N\mu/V)2\sinh x/(1+2\cosh x)$，$\chi_{\rm total}=2N\mu^2/(3\kb T)$，单位体积响应再除$V$。磁场自然势已包含$-BM$，不能再加一次磁功造成重复。
\end{solution}
\begin{exercise}[原题6(a)--(e)]\label{p:333-4-6}
链由$N$个盘状分子构成，长轴方向贡献长度$2a$能0，短轴方向长度$a$能$\varepsilon$，$H=\varepsilon U$，$U$为短轴数。(a)配分。(b)两方向概率。(c)平均长。(d)长度方差。(e)大数分布。
\end{exercise}
\begin{solution}
$Z=(1+\e^{-\beta\varepsilon})^N$，短方向$q=1/(1+\e^{\beta\varepsilon})$、长$1-q$。$U$为Binomial$(N,q)$，$L=a(2N-U)$，故$\avg L=aN(2-q)$、$\operatorname{Var}L=a^2Nq(1-q)$。若$Nq(1-q)\gg1$，中心窗口Gaussian；极低温$q$小到$Nq=O(1)$时应使用稀有短轴的Poisson或准确二项，而不是未经条件的全温Gaussian。
\end{solution}
\begin{exercise}[原题7(a)--(e)]\label{p:333-4-7}
极性刚棒惯性矩$I$、偶极$\mu$，外电场$E$。经典转动$H=(p_\theta^2+p_\phi^2/\sin^2\theta)/(2I)-\mu E\cos\theta$。(a)转动配分。(b)极化。(c)零场极化率。(d)转动能及高低温。(e)热容曲线。
\end{exercise}
\begin{solution}
Gaussian动量积分给$(2\pi I/\beta)\sin\theta$，角积分$4\pi\sinh x/x$，$x=\beta\mu E$，所以$Z_{\mathrm{rot}}=8\pi^2I\kb T\sinh x/(h^2x)$。极化$P=\mu(\coth x-1/x)$、$\chi_T=\mu^2/(3\kb T)$。$U_{\mathrm{rot}}=2\kb T-\mu E\coth x$，$C_E/\kb=2-x^2\operatorname{csch}^2x$。高温$x\ll1$，$U\simeq\kb T-(\mu E)^2/(3\kb T)$、$C\to\kb$；形式低温$x\gg1$，$U\simeq-\mu E+2\kb T$、$C\to2\kb$，因对齐附近两角方向也有谐振势能。真正$T\to0$须量子转动冻结，经典结果不能违背第三律。
\begin{center}\begin{tikzpicture}\begin{axis}[width=10cm,height=4.6cm,xlabel={$\kb T/(\mu E)$},ylabel={$C_E/\kb$},xmin=0,xmax=3,ymin=.9,ymax=2.1]
\addplot[structurecolor,thick,domain=.04:3,samples=150]{2-(1/x)^2/(sinh(1/x))^2};\end{axis}\end{tikzpicture}\end{center}
\end{solution}
\begin{exercise}[原题8(a)--(c)，选做]\label{p:333-4-8}
无序玻璃有独立双态缺陷，能$\epsilon_i,\epsilon_i+\delta_i$。(a)求$Z,E,C$。(b)激发差态密度$\rho(\delta)=N/\Delta$于$0<\delta<\Delta$，求热容与低高温。(c)什么低能态密度特征使$C\propto T$？
\end{exercise}
\begin{solution}
$Z=\prod_i\e^{-\beta\epsilon_i}(1+\e^{-\beta\delta_i})$，$E=\sum_i[\epsilon_i+\delta_i/(\e^{\beta\delta_i}+1)]$，$C=\kb\sum_i(\beta\delta_i)^2\e^{\beta\delta_i}/(1+\e^{\beta\delta_i})^2$。均匀密度给$C=N\kb(\kb T/\Delta)\int_0^{\Delta/\kb T}x^2\e^x/(1+\e^x)^2\dd x$。低温积分$\pi^2/6$，$C\simeq\pi^2N\kb^2T/(6\Delta)$；高温$C\simeq N\Delta^2/(12\kb T^2)$。低能$\rho(0)$有限非零且平滑给线性；$\rho\propto\delta^a$则$C\propto T^{a+1}$。整体基能$\epsilon_i$不改变热容。
\end{solution}
\begin{exercise}[原题9(a)--(b)，选做]\label{p:333-4-9}
canonical相空间任意坐标$X_i$与函数$f$。(a)证明$\avg{\partial_if}=\beta\avg{f\partial_iH}$。(b)分别$f=q_i,p_i$得到virial形式，说明表面项。
\end{exercise}
\begin{solution}
对$\partial_i(f\e^{-\beta H})$积分，在边界项0时得恒等式。$f=X_i$给$\avg{X_i\partial_iH}=\kb T$；动量无穷远Gaussian使边界0，二次动量每自由度$\avg{p_i^2/m}=\kb T$。坐标若硬壁有限且$f=q_i$在壁不为0，必须保留表面项，等价将墙力纳入$H$。三维中心两体势的压力virial为$pV=N\kb T-\frac13\avg{\sum_{i<j}r_{ij}v'(r_{ij})}$。无墙的齐次约束势$v\propto r^a$则相应$2\avg K=a\avg U$。不能对自由硬盒忽略墙项再由$\partial_qH=0$导出$0=\kb T$。
\end{solution}

\originalchapter{8.333作业5：既有热统卷的逐题接口}
\source{Mehran Kardar，8.333 Fall 2013，PS5，截止2013年11月22日，12题。此前热统卷第33章已包含完整题面与独立解答，本卷按协调范围引用，不复制这12题的重复页面。}
本项目既有《热学与统计物理》第33章（PDF物理页61--65）已完整写出下表各题及全部明确小问，包括原选做3、5、7、11。以下接口指向固定提交中的中文成果，而不是把英文原件当中文交付：
\begin{center}\begin{tabular}{lll}\toprule 原题&中文位置&内容\\\midrule
1(a)--(f)&习题33.1&表面活性剂凝聚\\
2(a)--(d)&习题33.2&三阶状态方程临界点\\
3(a)--(e)&习题33.3&Manning转变（选做）\\
4(a)--(d)&习题33.4&点吸引能量累积量与塌缩\\
5(a)--(d)&习题33.5&二维硬杆几何近似与熵比较（选做）\\
6(a)--(c)&习题33.6&电子自旋密度矩阵\\
7(a),(b)&习题33.7&平面量子转子（选做）\\
8(a)--(c)&习题33.8&量子熵、最大熵与幺正演化\\
9(a)--(c)&习题33.9&零点能的渐近匹配\\
10(a)--(f)&习题33.10&振动、转动高温量子修正\\
11(a)--(d)&习题33.11&氢的ortho/para统计（选做）\\
12&习题33.12&Bohr--van Leeuwen定理\\\bottomrule\end{tabular}\end{center}
\url{https://github.com/xhc144/mit-ocw-notes/blob/eef678cd6cbfa047d6bfb7697334db7be33b6368/subjects/thermal-statistical-physics/thermal-statistical-physics-zh.pdf}

本卷第\ref{ch:ensemble}、\ref{ch:clusters}章系统重新建立这些前置理论，后面的复习题仅在原条件完全相合时引用这一接口；有差异的小问另外补写。例如Manning题的势零点随容器变化必须先恢复固定基准再求侧壁压力；硬杆驻点曲线折点不能当全局相变，须比较两支熵。既有卷已经保留这两项限定。固定源TeX、PDF与页数的SHA-256见\texttt{qa/prior-volume-reference.json}，正常编译本卷不需要下载或输入其他卷文件。

\originalchapter{8.333作业6：量子气体与低能激发}
\psource{Mehran Kardar，8.333 Fall 2013，PS6，截止2013年12月6日。选做题5、9--13也收入。}
\begin{exercise}[原题1(a)--(d)]\label{p:333-6-1}
原题提供代表表（He3能量栏本身漏印系数，保留缺口）：
\begin{center}\begin{tabular}{lllll}\toprule 粒子&$n$/m$^{-3}$&$m$/kg&$E_F$/eV&$T_F$/K\\\midrule 电子&$10^{29}$&$9\times10^{-31}$&4.4&$5\times10^4$\\ 核子&$10^{44}$&$1.6\times10^{-27}$&$10^8$&$1.1\times10^{12}$\\ 液He3&$2.6\times10^{28}$&$4.6\times10^{-27}$&$\text{缺系数}\times10^{-3}$&$10^1$\\\bottomrule\end{tabular}\end{center}
另给He3原子体积46.2 Å$^3$，与代表密度未严格一致；本题所求估算不使用缺印数值。
估算：(a)典型金属室温电子与声子热容比（表中电子$T_F\simeq5\times10^4$ K）；(b)室温中子热波长与晶体声子最短波长；(c)室温、常压H$_2$、He、O$_2$的$n\lambda^3$及出现量子效应的温度；(d)He4在1 K以下每质量热容$20.4T^3$ J kg$^{-1}$ K$^{-1}$，仅一支无简并低能激发，求$E(k)$。
\end{exercise}
\begin{solution}
(a)取每原子一个导电电子，电子$C_e=\pi^2N\kb T/(2T_F)$，高温晶格$C_{\rm ph}\simeq3N\kb$，故比$\pi^2T/(6T_F)\simeq0.0099$。若声子尚未到Dulong--Petit区，应使用实际Debye温度，不能当普适金属常数。

(b)$\lambda=h/\sqrt{2\pi m\kb T}$，300 K中子约1.00 Å；声子布里渊区边界$k\sim\pi/a$，最短波长约$2a$，若$a\sim2$ Å则约4 Å。比较尺度即可，不需要为“典型晶体”虚构晶格参数。(c)$n=P/(\kb T)$，取300 K、101325 Pa，三者依次$\lambda=(0.710,0.502,0.178)$ Å、$n\lambda^3=(8.76\times10^{-6},3.10\times10^{-6},1.37\times10^{-7})$。定密度冷却使$n\lambda^3\propto T^{-3/2}$，量子尺度依次$0.128,0.0638,0.00797$ K；定压则$T^{-5/2}$，依次$2.85,1.88,0.539$ K。实际气体可能先液化或冻结，这些是理想气外推。

(d)三维单支各向同性$E\propto k^s$使$C\propto T^{3/s}$，所以$s=1$。设$E=\hbar ck$，对Bose分布积分得$C/V=2\pi^2\kb^4T^3/(15\hbar^3c^3)$。还需质量密度；采用液He4约145 kg m$^{-3}$，$c=[2\pi^2\kb^4/(15\hbar^3\rho\,20.4)]^{1/3}\simeq240$ m/s。热容系数本身不足以无条件确定密度未知材料的声速；以上额外物性假设明确列出。
\end{solution}
\begin{exercise}[原题2(a)--(d)]\label{p:333-6-2}
太阳中心$T=1.6\times10^7$ K，氢质量密度$6\times10^4$、氦$10^5$ kg m$^{-3}$，完全电离。(a)电子、质子、$\alpha$粒子热波长；(b)理想近似下是否简并；(c)总气体压强；(d)辐射压强及主要支撑来源。
\end{exercise}
\begin{solution}
$n_p=\rho_H/m_p$，$n_\alpha=\rho_{He}/(4m_p)$，电中性$n_e=n_p+2n_\alpha$。热波长为$1.863\times10^{-11},4.349\times10^{-13},2.174\times10^{-13}$ m。按电子、质子自旋简并2，$n\lambda^3/g$为$0.213,1.48\times10^{-6},1.54\times10^{-7}$；核完全经典，电子弱简并，可有约百分之几修正，主阶仍可用经典压力。

$P_{\rm gas}=\kb T(n_e+n_p+n_\alpha)=2.58\times10^{16}$ Pa。$P_{\rm rad}=a_RT^4/3=1.65\times10^{13}$ Pa，$a_R=7.5657\times10^{-16}$ J m$^{-3}$ K$^{-4}$。气体约大1558倍，在给定中心条件主要物质压强支撑；这不取代整个恒星的静力结构方程。
\end{solution}
\begin{exercise}[原题3(a)--(e)]\label{p:333-6-3}
三维非相互作用同种粒子，自旋简并$g$，$\varepsilon=p^2/(2m)$、密度$n$。保留$n^2$项求(a)压力、(b)能量密度、(c)$C_V/N$、(d)$\alpha_P=V^{-1}(\partial_TV)_{P,N}$、(e)$C_P/N$，比较Bose/Fermi。
\end{exercise}
\begin{solution}
令$\eta=+1$为Bose、$-1$为Fermi，$f_s^\eta(z)=\sum_{l\ge1}\eta^{l-1}z^l/l^s$，$n=g\lambda^{-3}f_{3/2}^\eta$、$P=g\kb T\lambda^{-3}f_{5/2}^\eta$。消去$z$得到$B_2=-\eta\lambda^3/(g2^{5/2})$。记$x=nB_2$，$P=n\kb T(1+x)$、$U/V=3P/2$，$x\propto nT^{-3/2}$。定$V$微分给
\begin{equation}
C_V/N=\tfrac32\kb(1-x/2),\qquad
\alpha_P=T^{-1}(1-5x/2).
\end{equation}
后式由$\alpha_P=P_T/(nP_n)$且$P_T=n\kb(1-x/2)$、$P_n=\kb T(1+2x)$。$C_P-C_V=TV\alpha_P^2/\kappa_T=N\kb(1-3x)$，故$C_P/N=\tfrac52\kb(1-3x/2)$。Bose$x<0$提高这些首阶热容，Fermi相反；只在$|x|\ll1$成立。
\end{solution}
\begin{exercise}[原题4(a)--(c)]\label{p:333-6-4}
某中子星核候选相有两支Fermi激发$E_\pm(k)=\pm\hbar^2(|k|-k_F)^2/(2M)$，自旋$g=2$。令$q=|k|-k_F$，用$\dd^3k\simeq4\pi k_F^2\dd q$：(a)证明$\Delta U=(2gVk_F^2/\pi^2)\int_0^\infty E_+(q)[\e^{\beta E_+(q)}+1]^{-1}\dd q$；(b)闭式积分；(c)低温热容。
\end{exercise}
\begin{solution}
粒子与空穴各贡献同一能量，再由$q$正负两侧得到另一个2，合并$4\pi/(2\pi)^3$即题中系数。设$a=\hbar^2/(2M)$，$y=\beta aq^2$，积分变$\tfrac12 a^{-1/2}(\kb T)^{3/2}\int_0^\infty y^{1/2}/(\e^y+1)\dd y$。用Fermi积分，
\begin{equation}
\Delta U=\frac{gVk_F^2}{\pi^2\sqrt a}\Gamma(3/2)(1-2^{-1/2})\zeta(3/2)(\kb T)^{3/2}.
\end{equation}
$C_V$等于系数乘$\tfrac32\kb(\kb T)^{1/2}$。二次零能壳的态密度$\rho(E)\propto E^{-1/2}$使$C\propto\sqrt T$，区别普通光滑Fermi面$C\propto T$；壳近似要求热展宽$q\ll k_F$。
\end{solution}
\begin{exercise}[原题5(a)--(f)，选做]\label{p:333-6-5}
单层石墨烯$E_\pm=\pm t_iak$，双层$E_\pm=\pm(t_i^2a^2/2t_\perp)k^2$，$g=4$，面积$A_s$，未掺杂。(a)求$\mu(T)$；(b)证明粒子空穴能量积分$\Delta U=2gA_s\int\dd^2k/(2\pi)^2\,E_+(k)/(\e^{\beta E_+}+1)$；(c),(d)双层能与热容；(e)面积各半的混合样品双层电子热容何时主导；(f)声速约$2\times10^4$ m/s的声子与单层电子低温热容比较。
\end{exercise}
\begin{solution}
粒子空穴对称固定半填充，$f(E)+f(-E)=1$，故$\mu=0$；有限带宽对称也成立。升温正能电子与负能空穴两项相同给(b)。记$a_2=t_i^2a^2/(2t_\perp)$、$b=t_ia$，径向积分为
\begin{align}
\Delta U_2/A_s&=\frac{g}{2\pi a_2}\int_0^\infty\frac{E\dd E}{\e^{\beta E}+1}=\frac{g\pi(\kb T)^2}{24a_2},&C_2/A_s&=\frac{g\pi\kb^2T}{12a_2},\\
\Delta U_1/A_s&=\frac{g}{\pi b^2}\int_0^\infty\frac{E^2\dd E}{\e^{\beta E}+1}=\frac{3g\zeta(3)(\kb T)^3}{2\pi b^2},&C_1/A_s&=\frac{9g\zeta(3)\kb^3T^2}{2\pi b^2}.
\end{align}
相等条件$\kb T_\times=\pi^2t_\perp/[27\zeta(3)]\simeq0.304t_\perp$，更低温双层主导，且须低能抛物近似仍有效。二维线性色散声子也给$T^2$，但系数$\propto\sum_s c_s^{-2}$；$v_e=t_ia/\hbar$通常远大声速，声子系数大。自由膜弯曲支$\omega\propto k^2$还给$T$，张力可改变它。故不宜仅凭电子态密度就判定石墨烯总热容由电子主导。
\end{solution}
\begin{exercise}[原题6(a)--(e)]\label{p:333-6-6}
单粒子态密度$\rho(\epsilon)$、基态0、$N=\int_0^\infty\rho(\epsilon)/(\e^{\beta(\epsilon-\mu)}-\eta)\dd\epsilon$。(a)能量；(b)Bose临界温隐式式；(c)Fermi低温$\mu-E_F$；(d)能量增量；(e)热容。可用Sommerfeld式$\int gf=\int_0^\mu g+\pi^2g'(\mu)/(6\beta^2)+\cdots$。
\end{exercise}
\begin{solution}
$U=\int\epsilon\rho(\epsilon)f_\eta\dd\epsilon$。Bose激发容量$N=\int_{0^+}^\infty\rho(\epsilon)/(\e^{\epsilon/(\kb T_c)}-1)\dd\epsilon$，要在低能收敛，且基态宏观占据单独处理。

Fermi固定$N$先展开数积分，$\rho_F\delta\mu+\pi^2(\kb T)^2\rho'_F/6=0$。能积分含$E_F\rho_F\delta\mu$，恰好抵消其导数项中的$E_F\rho'_F$，因此
\begin{equation}
\delta\mu=-\frac{\pi^2(\kb T)^2}{6}\frac{\rho'_F}{\rho_F},\quad
\Delta U=\frac{\pi^2}{6}\rho_F(\kb T)^2,\quad C_V=\frac{\pi^2}{3}\rho_F\kb^2T.
\end{equation}
须$\rho_F$有限且在$\kb T$窗口平滑；前题二次零壳不满足，不能套用。
\end{solution}
\begin{exercise}[原题7(a)--(e)]\label{p:333-6-7}
无自旋量子点粒子的弱吸引把每态能降低$uN^2/(2V)$，$Z=Z_0\exp[\beta uN^2/(2V)]$、$u>0$。(a)关联压力；(b)证明$P_n=\kb T f_{3/2}^\eta/f_{1/2}^\eta-un$，$n\lambda^3=f_{3/2}^\eta$；(c)稀薄不稳定临界$u_c$首阶量子修正；(d)Fermi零温；(e)Bose降低温度时的行为。
\end{exercise}
\begin{solution}
$F=F_0-uN^2/(2V)$，$P=-F_V=P_0-un^2/2$。理想气$\partial_zP_0=(\kb T/\lambda^3z)f_{3/2}$、$\partial_zn=(\lambda^{-3}/z)f_{1/2}$，即(b)。机械稳定$P_n>0$，故$u_c=(\kb T/n)f_{3/2}/f_{1/2}$。低密度消去$z$得
\begin{equation}
u_c=\frac{\kb T}{n}[1-\eta n\lambda^3/2^{3/2}+O((n\lambda^3)^2)].
\end{equation}
Fermi零温$P_0=2nE_F/5$、$E_F\propto n^{2/3}$，故$u_c=2E_F/(3n)$；Pauli压力抗塌缩。Bose接近$z=1$时$f_{1/2}\to\infty$而$f_{3/2}$有限，$u_c\to0$。凝聚相理想压力独立$n$，任意吸引使负压缩率。该点吸引模型无短程排斥，不能由不稳定条件直接证明稳定液气共存，可能整体塌缩。
\end{solution}
\begin{exercise}[原题8(a)--(f)]\label{p:333-6-8}
$d$维各向异性谐陷阱，$U=m\sum_\alpha\omega_\alpha^2x_\alpha^2/2$，$\hbar\omega_\alpha\ll\kb T$，忽略零点与离散性。(a)累积态数及态密度；(b)巨正则粒子数；(c)能量；(d)高温每粒子能；(e)零温$\mu$；(f)低温热容尺度。按每个自旋组分计。
\end{exercise}
\begin{solution}
相空间椭球或振子占据数单纯形$\sum\hbar\omega_\alpha n_\alpha<E$的体积给$\mathcal N(E)=E^d/[d!\prod\hbar\omega_\alpha]$、$\rho(E)=E^{d-1}/[(d-1)!\prod\hbar\omega_\alpha]$。积分得到
\begin{equation}
N=f_d^-(z)\prod_\alpha\frac{\kb T}{\hbar\omega_\alpha},\quad U=d\kb T f_{d+1}^-(z)\prod_\alpha\frac{\kb T}{\hbar\omega_\alpha}.
\end{equation}
经典$z\ll1$比值1，$U/N=d\kb T$，每维动能和势能各$\kb T/2$。零温$E_F=[d!N\prod\hbar\omega_\alpha]^{1/d}$。代$\rho(E_F)=dN/E_F$入前题，$C_V=\pi^2dN\kb^2T/(3E_F)$，且需$\hbar\omega\ll\kb T\ll E_F$才能同时使用连续谱与低温展开。
\end{solution}
\begin{exercise}[原题9(a)--(d)，选做]\label{p:333-6-9}
密度$\rho_t$的光晶格陷位，每位空或最多一个局域粒子，能$-w<0$，对游离理想Bose谱$p^2/(2m)$无影响，总密度$n$。(a)陷位占据密度；(b)游离密度（自旋$g$）；(c)凝聚临界总密度并画小$w$情形；(d)凝聚相能量密度。
\end{exercise}
\begin{solution}
单陷位$\Xi_t=1+\e^{\beta(\mu+w)}$，故$n_t=\rho_t/[1+\e^{-\beta(\mu+w)}]$；硬占据约束使公式类似Fermi，与游离Bose统计不冲突。$n_g=g\lambda^{-3}\operatorname{Li}_{3/2}(z)$。凝聚$\mu\to0$给
\begin{equation}
n_*(T)=\frac{\rho_t}{1+\e^{-w/(\kb T)}}+g\lambda^{-3}\zeta(3/2),\quad
U/V=-wn_t(0,T)+\tfrac32g\kb T\lambda^{-3}\zeta(5/2).
\end{equation}
基态游离凝聚粒子能0。小$w$时低温陷位几乎满$n_*(0)=\rho_t$，升温陷位逐渐释放可使临界密度先降，再由游离激发容量$T^{3/2}$增长，允许再入凝聚。图取$g\lambda(T_\rho)^{-3}\zeta(3/2)=\rho_t$、$w/(\kb T_\rho)=0.01$。
\begin{center}\begin{tikzpicture}\begin{axis}[width=10cm,height=4.5cm,xlabel={$T/T_\rho$},ylabel={$n_*/\rho_t$},xmin=0,xmax=.8,ymin=.45,ymax=1.3]
\addplot[thick,structurecolor,domain=.0005:.8,samples=300]{1/(1+exp(-.01/x))+x^1.5};\end{axis}\end{tikzpicture}\end{center}
\end{solution}
\begin{exercise}[原题10(a)--(d)，选做]\label{p:333-6-10}
经典吸引的Mayer函数$f(r)=\exp(-\pi r^2/\lambda^2)$，傅里叶$f(\omega)=\lambda^3\exp[-\lambda^2\omega^2/(4\pi)]$。(a)环积分$R_l=\int\prod_i\dd^3q_i/V\prod_if(q_i-q_{i+1})$；(b)证明仅保留环时$\log Z=\log Z_0+V[\operatorname{Li}_{5/2}(x)-x+(1-2^{-5/2})x^2]/(2\lambda^3)$，$x=n\lambda^3$；(c)压力；(d)压缩率预示的凝聚不稳定。
\end{exercise}
\begin{solution}
(a)一个平移积分为$V$，其余卷积令各傅里叶动量相同，$R_l=V^{1-l}\int\dd^3\omega/(2\pi)^3[f(\omega)]^l=\lambda^{3l}/[l^{3/2}\lambda^3V^{l-1}]$。

(b)$l\ge3$环的顶点循环和反转对称数$2l$给$N^lR_l/(2l)$；$l=2$只有一条连边，精确项$N^2\lambda^3/(2V)$，不能把双边环重复计。合成题示多对数；$l=1$扣去不存在的自环。令$c_0=1-2^{-5/2}$，定$N$取体积导数得
\begin{equation}
P=n\kb T+\frac{\kb T}{2\lambda^3}[\operatorname{Li}_{5/2}(x)-\operatorname{Li}_{3/2}(x)-c_0x^2],\quad
\frac{P_n}{\kb T}=1-\frac{\operatorname{Li}_{1/2}(x)-\operatorname{Li}_{3/2}(x)}{2x}-c_0x.
\end{equation}
$x\to1^-$时$\operatorname{Li}_{1/2}$发散，因此必在收敛区间出现负压缩率；数值零点约0.70436。原题称凝聚转变，这里严格能推出环近似的机械不稳定；纯吸引无稳定核的原始经典模型可能塌缩，环求和本身不能证明有限密度稳定液相。
\end{solution}
\begin{exercise}[原题11(a)--(g)，选做]\label{p:333-6-11}
$d$维均匀无自旋Bose气$\epsilon=p^2/(2m)$。(a)巨势、密度用$f_s^+(z)=\operatorname{Li}_s(z)$表示；(b)$PV/U$；(c)$T_c(n)$；(d)凝聚以下热容；(e)全温热容图；(f)最大热容与经典极限比及$d=3$值；(g)$d\to2$与适用维数。
\end{exercise}
\begin{solution}
令$s=d/2$。对$-\log(1-z\e^{-\beta\epsilon})$作动量积分或逐项积分，$\mathcal G=-\kb TV\lambda^{-d}\operatorname{Li}_{s+1}(z)$、$n=\lambda^{-d}\operatorname{Li}_s(z)$，$U=sPV$，比$2/d$与经典相同。

只在$d>2$时$\zeta(s)$有限，$T_c=(2\pi\hbar^2/m\kb)[n/\zeta(s)]^{1/s}$。以下$z=1$，$C/(N\kb)=s(s+1)\zeta(s+1)/\zeta(s)(T/T_c)^s$。以上固定$N$求$z_T$得到
\begin{equation}
\frac C{N\kb}=s(s+1)\frac{\operatorname{Li}_{s+1}}{\operatorname{Li}_s}-s^2\frac{\operatorname{Li}_s}{\operatorname{Li}_{s-1}},\qquad
(T/T_c)^s=\zeta(s)/\operatorname{Li}_s(z).
\end{equation}
$d=3$的峰在$T_c$，$C_{\max}/C_\infty=(5/2)\zeta(5/2)/\zeta(3/2)=1.28378$。下面曲线直接用上式消去$z$，不是手工连接。
\begin{center}\begin{tikzpicture}\begin{axis}[width=10cm,height=4.5cm,xlabel={$T/T_c$},ylabel={$C/(N\kb)$},xmin=0,xmax=4,ymin=0,ymax=2.1]
\addplot[thick,structurecolor] coordinates {(0,0) (0.0144927536,0.00335976272) (0.0289855072,0.009502844) (0.0434782609,0.0174578392) (0.0579710145,0.0268781017) (0.0724637681,0.0375632891) (0.0869565217,0.0493782259) (0.101449275,0.0622236763) (0.115942029,0.076022752) (0.130434783,0.0907135933) (0.144927536,0.106245026) (0.15942029,0.122573795) (0.173913043,0.139662713) (0.188405797,0.157479358) (0.202898551,0.175995134) (0.217391304,0.195184576) (0.231884058,0.215024814) (0.246376812,0.235495161) (0.260869565,0.256576788) (0.275362319,0.278252457) (0.289855072,0.300506313) (0.304347826,0.323323706) (0.31884058,0.346691048) (0.333333333,0.370595687) (0.347826087,0.395025807) (0.362318841,0.41997034) (0.376811594,0.445418887) (0.391304348,0.471361658) (0.405797101,0.49778941) (0.420289855,0.524693402) (0.434782609,0.552065348) (0.449275362,0.579897381) (0.463768116,0.608182016) (0.47826087,0.636912124) (0.492753623,0.666080902) (0.507246377,0.69568185) (0.52173913,0.725708747) (0.536231884,0.756155634) (0.550724638,0.787016797) (0.565217391,0.818286746) (0.579710145,0.849960206) (0.594202899,0.882032101) (0.608695652,0.914497541) (0.623188406,0.947351811) (0.637681159,0.980590364) (0.652173913,1.01420881) (0.666666667,1.04820289) (0.68115942,1.08256852) (0.695652174,1.11730171) (0.710144928,1.15239861) (0.724637681,1.1878555) (0.739130435,1.22366875) (0.753623188,1.25983486) (0.768115942,1.29635041) (0.782608696,1.3332121) (0.797101449,1.37041669) (0.811594203,1.40796107) (0.826086957,1.44584218) (0.84057971,1.48405705) (0.855072464,1.5226028) (0.869565217,1.56147661) (0.884057971,1.60067573) (0.898550725,1.64019748) (0.913043478,1.68003926) (0.927536232,1.72019851) (0.942028986,1.76067274) (0.956521739,1.80145953) (0.971014493,1.84255649) (0.985507246,1.88396129) (1,1.92567168)};
\addplot[thick,structurecolor] coordinates {(1.0000009,1.92567097) (1.00000109,1.92567083) (1.00000132,1.92567065) (1.00000159,1.92567044) (1.00000193,1.92567018) (1.00000233,1.92566987) (1.00000281,1.92566949) (1.00000339,1.92566904) (1.0000041,1.92566849) (1.00000495,1.92566783) (1.00000598,1.92566703) (1.00000722,1.92566606) (1.00000873,1.92566489) (1.00001054,1.92566348) (1.00001273,1.92566178) (1.00001538,1.92565972) (1.00001858,1.92565724) (1.00002244,1.92565423) (1.00002711,1.92565061) (1.00003274,1.92564623) (1.00003955,1.92564094) (1.00004777,1.92563455) (1.0000577,1.92562683) (1.0000697,1.92561751) (1.00008419,1.92560625) (1.00010169,1.92559265) (1.00012284,1.92557622) (1.00014838,1.92555637) (1.00017923,1.9255324) (1.0002165,1.92550345) (1.00026152,1.92546849) (1.0003159,1.92542626) (1.0003816,1.92537525) (1.00046097,1.92531364) (1.00055684,1.92523924) (1.00067267,1.92514938) (1.00081261,1.92504086) (1.00098169,1.92490982) (1.00118597,1.92475157) (1.00143282,1.92456049) (1.00173111,1.92432977) (1.00209161,1.92405124) (1.00252733,1.923715) (1.00305404,1.92330915) (1.00369085,1.92281935) (1.0044609,1.92222836) (1.00539231,1.92151543) (1.0065192,1.9206556) (1.00788307,1.91961897) (1.00953444,1.91836965) (1.01153492,1.91686469) (1.01395981,1.91505283) (1.01690135,1.91287292) (1.02047285,1.91025233) (1.02481404,1.90710503) (1.03009795,1.90332955) (1.03654001,1.8988068) (1.04441011,1.89339788) (1.05404903,1.88694194) (1.0658912,1.87925439) (1.08049679,1.87012574) (1.09859845,1.85932138) (1.121171,1.84658293) (1.14953935,1.83163168) (1.18555187,1.81417478) (1.21293867,1.80191968) (1.24757113,1.78754758) (1.28012684,1.77506693) (1.31127106,1.7639655) (1.34142163,1.75392468) (1.3708629,1.74473057) (1.39980017,1.73623153) (1.42838874,1.72831553) (1.45675061,1.72089708) (1.48498475,1.71390925) (1.51317378,1.70729852) (1.54138833,1.70102131) (1.56969024,1.69504158) (1.59813471,1.68932916) (1.62677196,1.6838585) (1.65564845,1.67860773) (1.68480785,1.67355798) (1.71429182,1.66869284) (1.74414055,1.66399793) (1.77439334,1.65946058) (1.80508898,1.65506954) (1.83626614,1.6508148) (1.86796368,1.64668739) (1.900221,1.64267924) (1.93307826,1.63878304) (1.96657667,1.6349922) (2.00075876,1.63130071) (2.03566861,1.62770306) (2.07135217,1.62419424) (2.10785744,1.62076963) (2.14523485,1.61742498) (2.18353747,1.61415637) (2.22282138,1.61096019) (2.26314601,1.60783306) (2.3045745,1.60477186) (2.34717409,1.60177369) (2.39101662,1.59883583) (2.43617897,1.59595576) (2.48274366,1.5931311) (2.53079944,1.59035963) (2.58044197,1.58763928) (2.63177463,1.58496807) (2.68490939,1.58234417) (2.73996779,1.57976583) (2.7970821,1.57723141) (2.85639657,1.57473938) (2.91806897,1.57228824) (2.98227227,1.56987663) (3.04919661,1.56750322) (3.11905159,1.56516675) (3.19206894,1.56286606) (3.26850564,1.56059999) (3.34864756,1.55836749) (3.4328138,1.55616752) (3.52136185,1.55399912) (3.61469364,1.55186134) (3.71326295,1.54975331) (3.81758427,1.54767417) (3.92824358,1.5456231)};
\addplot[dashed,gray,domain=0:4]{1.5};\end{axis}\end{tikzpicture}\end{center}


一般$d$必须区分临界值与全温最大值。临界左值比$R_c=(s+1)\zeta(s+1)/\zeta(s)$在$s\downarrow1$趋0，但经典极限比恒为1，因此此时$R_c$显然不能是全温最大值。严格答案为上述两分支热容相对$sN\kb$的上确界；例如$s=1.05$时全温峰约1.00140，在$T>T_c$，不是临界值0.1594。随$d\to2^+$，$T_c\to0$、临界值趋0，而全温最大比趋1。$d\le2$均匀气在有限温无此BEC，必须直接用正常分支。原题若将三维“临界峰”结论外推所有$d$，会产生该矛盾。
\end{solution}
\begin{exercise}[原题12(a)--(d)，选做]\label{p:333-6-12}
三维非简并Bose气$P,T$与吸引二维表面接触，吸附每粒子降低能$u$。(a)气体化学势；(b)表面理想Bose面密度；(c)允许的最大压力$P_*$；(d)$\delta P=P_*-P\to0$的奇异性。
\end{exercise}
\begin{solution}
气相$\mu=\kb T\log[P\lambda^3/(\kb T)]$。面内最低能$-u$，有效逸度$z_s=\e^{\beta(\mu+u)}$，二维积分$n_2=\lambda^{-2}\operatorname{Li}_1(z_s)=-\lambda^{-2}\log(1-z_s)$。需$z_s<1$，所以$P_*=(\kb T/\lambda^3)\e^{-\beta u}$、$n_2\simeq\lambda^{-2}\log(P_*/\delta P)$。原题措辞“完全凝聚至表面”指与无限三维库平衡的极限：二维均匀理想气在有限温没有有限临界面密度，不能据此声称发生通常二维BEC。
\end{solution}
\begin{exercise}[原题13(a)--(d)，选做]\label{p:333-6-13}
平行板以$v$运动，超流激发谱$\epsilon(p)$。(a)由大质量板的动量能量改变得板系激发能$\epsilon-p\cdot v$；(b)激发净动量积分；(c)小$v$定义$P_\alpha=V\rho_nv_\alpha$并求$\rho_n$；(d)声子$\epsilon=cp$的低温贡献。
\end{exercise}
\begin{solution}
板给予$p$后自身动量减少$p$，能变$[(Mv-p)^2-(Mv)^2]/(2M)=-p\cdot v+O(M^{-1})$，总产生能即题式。$P=V\int\dd^3p/(2\pi\hbar)^3\,p\,n_B(\epsilon-p\cdot v)$；零速项奇函数0，线性项$-n_B'(\epsilon)p_\alpha p_\gamma v_\gamma$。角平均给
\begin{equation}
\rho_n=\frac13\int\frac{\dd^3p}{(2\pi\hbar)^3}p^2[-n_B'(\epsilon)]=\frac{2\pi^2(\kb T)^4}{45\hbar^3c^5}.
\end{equation}
最后以$x=\beta cp$积分，$\int_0^\infty x^4\e^x/(\e^x-1)^2\dd x=4!\zeta(4)$。这是质量密度，由$P/V$和$v$量纲验证，超流质量密度为总密度减$\rho_n$；激发数密度不能直接替代正常质量分数。
\end{solution}

\originalchapter{8.333复习包1：界面热力学与概率}
\osource{8.333 Fall 2013课程页公开的Fall 2007 Test1 Review Problems及官方Solutions；9题，非2013年实际测验卷。采用官方热量定义，补充近似条件。}
\begin{exercise}[原题1(a)--(c)]\label{r:333-1-1}
表面张力$\sigma$以功$\delta W=\sigma\dd A$定义。(a)球形水滴半径$R$的内外压差及肥皂泡压差；(b)空气、固体、水三界面张力$\sigma_{aw},\sigma_{sw},\sigma_{sa}$给接触角与完全润湿条件；(c)室温水$\sigma_0\simeq0.07$ N/m，估重力/毛细分界尺度并举例。
\end{exercise}
\begin{solution}
水滴$\Delta p\,4\pi R^2\dd R=\sigma\,8\pi R\dd R$，故$\Delta p=2\sigma/R$。肥皂泡有内外两张界面，薄膜近似$4\sigma/R$。三相线平行于固面的张力平衡$\sigma_{sa}=\sigma_{sw}+\sigma_{aw}\cos\theta$（$\theta$从水内测），$\cos\theta=(\sigma_{sa}-\sigma_{sw})/\sigma_{aw}$。若$\sigma_{sa}\ge\sigma_{sw}+\sigma_{aw}$，以水膜替代干界面不增加自由能，完全润湿；粗糙、滞后、钉扎可改变观察角，不包含于此平衡模型。

比较$\rho g\ell$与$\sigma/\ell$，毛细长$\ell_c=\sqrt{\sigma/(\rho g)}\simeq2.7$ mm。毫米小液滴趋球形，大水面受重力平整；毛细管弯月面、细针或水黾在水面承重用到这一竞争。
\end{solution}
\begin{exercise}[原题2(a)--(d)]\label{r:333-1-2}
表面活性剂$N$分子铺于面积$A$。(a)解释$\sigma\simeq\sigma_0-N\kb T/A$；(b)肥皂触滴表面时接触角变化；(c)给$\sigma_A=N\kb T/(A-Nb)^2-2aN^2/A^3$、$\sigma_T=-N\kb/(A-Nb)$，求$\sigma(A,T)$及$a,b$意义；(d)用$E_A,\sigma,\sigma_T,\sigma_A$表达定张力与定面积热容差。此处热容为$\delta Q/\dd T$，原题紧接定义的抄写易漏机械功，应遵循官方解的热量式。
\end{exercise}
\begin{solution}
二维理想分子铺展压为$N\kb T/A$，自由能随面积增大降低，张力减少该量。若固体两界面张力近似固定，降$\sigma_{aw}$使前题$\cos\theta$增大、滴更铺开；肥皂同时改变固--水界面时应重新比较全部三张力。

积分两个偏导并用稀薄极限定常数得$\sigma=\sigma_0-N\kb T/(A-Nb)+aN^2/A^2$。$b$为排除面积，吸引$a>0$减少分子铺展压、回升张力。原两偏导交叉相等提供可积性检查。

$\delta Q=\dd E-\sigma\dd A=(E_A-\sigma)\dd A+C_A\dd T$，定$\sigma$有$A_T=-\sigma_T/\sigma_A$，所以
\begin{equation}
C_\sigma-C_A=-(E_A-\sigma)\frac{\sigma_T}{\sigma_A}=T\frac{\sigma_T^2}{\sigma_A},
\end{equation}
最后用$E_A=\sigma-T\sigma_T$，来自$\dd F=-S\dd T+\sigma\dd A$。若只求$\partial_TE|_\sigma$，则还差机械功$\sigma A_T$，不能将这导数误当定张力吸热热容。
\end{solution}
\begin{exercise}[原题3(a)--(c)]\label{r:333-1-3}
理想气温标$\Theta$满足$PV=N\kb\Theta$且$E=E(\Theta)$，不能假设$E\propto\Theta$。做Carnot循环：(a)两等温吸放热；(b)任意热容时的绝热体积变化；(c)证明气体温标与热力学温标等价。
\end{exercise}
\begin{solution}
热端等温$Q_H=N\kb\Theta_H\log(V_2/V_1)$，冷端放热绝对值$Q_C=N\kb\Theta_C\log(V_3/V_4)$，因为等温$E$不变。绝热满足$E'(\Theta)\dd\Theta=-N\kb\Theta\dd V/V$，故$\log(V_f/V_i)=-\int_{\Theta_i}^{\Theta_f}E'(\Theta)\dd\Theta/(N\kb\Theta)$。两条反向绝热线的积分相反，得到$V_2/V_1=V_3/V_4$，因此$Q_H/Q_C=\Theta_H/\Theta_C$。Carnot热力学定义$Q_H/Q_C=T_H/T_C$，温标比例固定，以同一参考点校准即$T=\Theta$，全程没有用常数热容。
\end{solution}
\begin{exercise}[原题4(a)--(c)]\label{r:333-1-4}
(a)由$\dd E=T\dd S-P\dd V$证明理想气内能仅依赖温度；(b)给与$E=E(T)$相容的最一般状态方程；(c)证明van der Waals气$C_V$只依赖温度。
\end{exercise}
\begin{solution}
$F$的Maxwell关系$S_V=P_T$，故$E_V|_T=TS_V-P=TP_T-P$。理想$P=N\kb T/V$给0。一般$E_V=0$等价$\partial_T(P/T)=0$，所以$P=T f(V,N)$，并非只允许理想气；稳定、广延性可继续限制$f$。van der Waals$P=N\kb T/(V-Nb)-aN^2/V^2$给$E_V=aN^2/V^2$，积分$E=E_0(T)-aN^2/V$，因此$C_V=E_0'(T)$，只在采用这套温度无关$a,b$的单相状态方程模型中成立。
\end{solution}
\begin{exercise}[原题5(a)--(f)]\label{r:333-1-5}
一摩尔水在共存$(P,T)$吸收潜热$L$变汽，体积差$\Delta V$，绝热降至$P-\dd P$、$T-\dd T$后凝结。(a)证明$W=\Delta V\dd P+O(\dd P^2)$及CC式；(b)指出以$L(T)-L(T-\dd T)=W$推$L\propto T$的错误；(c)恒潜热、理想汽时$P(T)$；(d)飓风热端80°F、冷端$-120$°F，暖水层至少200英尺，每小时凝结9000万吨，潜热$2.3\times10^6$ J/kg，估最高效率与功率；(e)等温理想空气静力$P(h)$；(f)估9 km珠峰水沸点，列出所用空气温度和分子质量。
\end{exercise}
\begin{solution}
微小Carnot在$P,V$图的面积为$\Delta V\dd P$，绝热边的横移只影响二阶；$W/Q_H=\dd T/T$给$\dd P/\dd T=L/(T\Delta V)$。绝热膨胀本身可凝结部分汽，低温等温段不再把完整一摩尔汽凝结；直接以两端同一摩尔潜热相减遗漏这一阶分率，官方解正指出此点。

取汽$\Delta V\simeq RT/P$，$\log[P(T)/P(T_0)]=-(L/R)(1/T-1/T_0)$。飓风给温$299.8,188.7$ K，最高效率约0.371；质量率$9\times10^{10}/3600=2.5\times10^7$ kg/s，供热$5.75\times10^{13}$ W，最大功约$2.13\times10^{13}$ W。200英尺约61 m暖层是持续蒸发的储库条件，并非此瞬时Carnot功率式的乘数；实际远低于可逆上界。

空气静力$\dd P=-\rho g\dd h$、$\rho=mP/(\kb T_a)$，所以$P(h)=P_0\e^{-mgh/(\kb T_a)}$。取等温$T_a=273$ K、$m=29$原子质量单位，尺度约8.0 km，9 km处约0.324大气压。每质量潜热用$R_v=\kb/m_{H_2O}\simeq461.5$ J kg$^{-1}$ K$^{-1}$，$1/T_b=1/373.15-(R_v/L_v)\log(P/P_0)$，约344 K、71°C。这是题设等温大气估计，不使用现实天气；换空气温度会改变估值。
\end{solution}
\begin{exercise}[原题6(a)--(f)]\label{r:333-1-6}
石英慢冷结晶、放潜热$L$，快冷成玻璃，均近似不可压缩，$\dd E=T\dd S+\mu\dd N$。(a)广延性给$\mu$；(b)晶体$C_X=\alpha T^3$、玻璃$C_G=\beta T$，假设都$S(0)=0$求熵；(c)再假设同$E_0$求内能；(d)化学平衡形式给$T_m$；(e)潜热；(f)检验符号及最可疑假设。这里$\alpha,\beta$按整个样品热容计。
\end{exercise}
\begin{solution}
Euler$E=TS+\mu N$，所以$\mu=(E-TS)/N$。积分$C/T$得$S_X=\alpha T^3/3$、$S_G=\beta T$；能$E_X=E_0+\alpha T^4/4$、$E_G=E_0+\beta T^2/2$。形式自由能$F_X=E_0-\alpha T^4/12$、$F_G=E_0-\beta T^2/2$，相等$T_m^2=6\beta/\alpha$，假设熔化潜热$E_G-E_X=-\beta T_m^2=-6\beta^2/\alpha<0$，与所述结晶放热趋势不符。

官方解指出将冻结非遍历玻璃当作唯一平衡基态、套$S_G(0)=0$是可疑步骤。应理解为平衡第三定律的\textbf{适用前提}未满足，不能据此说平衡态第三定律被推翻。近似同$E_0$及将低温热容幂律外推到熔点也未经验证；玻璃转变本身含动力学，不能简单用两平衡相$\mu$相等预测。负潜热暴露了这些假设的共同不相容。
\end{solution}
\begin{exercise}[原题7(a)--(e)]\label{r:333-1-7}
求(a)$p=1/(2a),|x|<a$；(b)$p=\e^{-|x|/a}/(2a)$；(c)$p=a/[\pi(x^2+a^2)]$的特征函数、均值、方差；(d)$p=x\e^{-x^2/(2a^2)}/a^2,x\ge0$；(e)$p=\sqrt{2/\pi}x^2\e^{-x^2/(2a^2)}/a^3,x\ge0$仅求均值、方差。
\end{exercise}
\begin{solution}
用$\varphi(k)=\avg{\e^{\ii kx}}$，(a)$\sin(ak)/(ak)$，均0、方差$a^2/3$；(b)正负半轴指数积分给$1/(1+a^2k^2)$，均0、方差$2a^2$；(c)留数或已知Fourier积分$\e^{-a|k|}$，在0不可微，通常均值与方差都不存在，主值均0不等于期望。后两者用$\int_0^\infty x^m\e^{-x^2/(2a^2)}\dd x=\tfrac12(2a^2)^{(m+1)/2}\Gamma((m+1)/2)$：Rayleigh均$a\sqrt{\pi/2}$、方差$a^2(2-\pi/2)$；Maxwell均$2a\sqrt{2/\pi}$、方差$a^2(3-8/\pi)$。
\end{solution}
\begin{exercise}[原题8]\label{r:333-1-8}
任意概率密度$p(x)$均$\lambda$、有限方差$\sigma^2$，证明$\Pr(|X-\lambda|\ge n\sigma)\le1/n^2$。
\end{exercise}
\begin{solution}
$\sigma^2=\int(x-\lambda)^2p(x)\dd x\ge\int_{|x-\lambda|\ge n\sigma}(x-\lambda)^2p\dd x\ge n^2\sigma^2\Pr(|X-\lambda|\ge n\sigma)$，非零方差除去即得。零方差分布集中于均值，需以正偏离阈值处理，不能拿$0/0$论证。仅有限二阶矩保证，不要求Gaussian。
\end{solution}
\begin{exercise}[原题9(a),(b)]\label{r:333-1-9}
$n$个独立$r_\alpha\in[0,1]$同密度$p$，最大$X$。(a)密度；(b)均匀$p$的均值、方差和大$n$，讨论最优选择减少差异的条件。
\end{exercise}
\begin{solution}
$F_X(x)=F(x)^n$，$p_X=np(x)F(x)^{n-1}$，$F(x)=\int_0^xp$。均匀则$nx^{n-1}$，$\avg X=n/(n+1)$、$\avg{X^2}=n/(n+2)$，方差$n/[(n+1)^2(n+2)]\sim n^{-2}$，相对均值的散布收缩。$n(1-X)$收敛指数分布，仍有非Gaussian剩余波动。这解释此独立、有限上界模型的选择收缩，不是所有人群变异都由最优选择决定的结论。
\end{solution}

\originalchapter{8.333复习包2：相空间、等离子体与输运}
\osource{8.333 Fall 2013页面公开的Fall 2007 Test2 Review及Solutions，7题。对Coulomb势、截面角度和Lorentz气Wick系数作独立符号与量纲复核，修正处明确写出。}
\begin{exercise}[原题1(a)--(d)]\label{r:333-2-1}
一维自由气粒子初始$\rho(q,p,0)=\delta(q)f(p)$，$f=\exp[-p^2/(2m\kb T)]/\sqrt{2\pi m\kb T}$。(a)Liouville解和相图；(b)两方差；(c)在$q=\pm Q$加硬壁，长时细密度；(d)有限相空间分辨率的粗粒化密度及其平稳性。
\end{exercise}
\begin{solution}
自由特征$q(t)=q_0+pt/m$、$p$常量，$\rho=\delta(q-pt/m)f(p)$，支持是一条随$t$倾斜的直线而非扩散成二维Gaussian。$\avg{q^2}=\kb Tt^2/m$、$\avg{p^2}=m\kb T$。硬壁反射$\pm p$，将直线按周期$4Q$折回盒子，得到细密度越来越细的折线，仍保存每一相点信息和$|p|$。

对固定小$q,p$单元平均，足够宽的速度小区间在长时覆盖许多反射周期，弱/粗粒化极限$\widetilde\rho=f(p)/(2Q)$（初始$f$偶）。它在盒内$q$均匀且满足反射对称，故Liouville平稳；精细密度本身不点态收敛。初始理想$\delta$的微分熵需先有限宽正则化，不能以无穷负熵作有限数值热力学熵。
\end{solution}
\begin{exercise}[原题2(a)--(d)]\label{r:333-2-2}
相空间规范化密度$S=-\int\rho\log\rho$。(a)Hamilton--Liouville演化$\dot S$；(b)固定$\avg H=E$的最大熵；(c)其平稳性；(d)如何与观察到趋平衡的熵增相容。
\end{exercise}
\begin{solution}
$\dot\rho=-\{\rho,H\}$，Hamilton流无散度且边界通量0，$\dot S=\int(1+\log\rho)\{\rho,H\}=\int\{\rho\log\rho,H\}=0$。对规范化和能量加乘子，变分$-(1+\log\rho)-a-\beta H=0$，故$\rho_{\max}=\e^{-\beta H}/Z$，$\beta$由$E$定；熵对$\rho$凹，确保最大。$\{f(H),H\}=0$，canonical平稳，但不意味着任何初态自动变成它。

粗粒化丢弃细结构，由凹性其熵不小于细密度熵；局部混合、与外界交换、或Boltzmann的分子混沌假设可使可观察熵增加。完全隔离精密可逆Hamilton演化并没有违反(a)，而coarse熵也未由Liouville单独保证每一时刻严格单调。
\end{solution}
\begin{exercise}[原题3(a)--(c)]\label{r:333-2-3}
Vlasov适用于相互作用范围$\lambda$内粒子数$n\lambda^3\gg1$，忽略相对弱的短时两体碰撞项。(a)$\rho_N=\prod_i\rho_1(x_i)$的$f_s$与规范化；(b)从BBGKY得$\partial_tf_1+(p/m)\cdot\nabla_qf_1-\nabla_qU_{\rm eff}\cdot\nabla_pf_1=0$，$U_{\rm eff}=U+\int V(q-q')f_1(x')\dd x'$；(c)均匀盒内任意动量$g(p)$为何平稳而无热化。
\end{exercise}
\begin{solution}
取$\int\rho_1=1$、约化密度按计数规范化，$f_s=N!/(N-s)!\prod_{i=1}^s\rho_1(x_i)$，$\int f_s=N!/(N-s)!$。$f_2=(1-1/N)f_1f_1$，代BBGKY的一体式得到平均场势$(1-1/N)\int Vf_1$，在大$N$忽略自作用的相对$1/N$差便是题式；各$s$乘积按同一一体流因子分离。这个无相关闭合不是任意有限$N$两体系统的精确演化。

均匀$f_1=N g(p)/V$（$\int g=1$）的密度势为空间常数，流项全0，所以任意$g$平稳。若原题把归一化单粒子密度写$g/V$，与计数$f_1$差一个$N$。没有碰撞核就没有要求$g$成为Maxwell的H定理；可有集体相混合，却不等于完整碰撞热平衡。
\end{solution}
\begin{exercise}[原题4(a)--(f)]\label{r:333-2-4}
中性等离子体各$N$个$\pm e$、质量$m_\pm$、密度$n_0$。(a)两组Vlasov及平均Coulomb势；(b)$f_\pm=g_\pm(p)n_\pm(q)$的Poisson式；(c)Boltzmann密度与Debye长；(d)屏蔽Green解；(e)Vlasov自洽条件及平均间距；(f)$\tau\sim\lambda/c_{\rm th}$为何不随$T$变、关联何物理频率。原题$\nabla^2C=+4\pi\delta$的势约定与正$1/r$Green式不相容，以下先统一标准物理势$\phi$。
\end{exercise}
\begin{solution}
取$C(r)=1/r$、$\nabla^2C=-4\pi\delta$、粒子势能$q_s\phi$，则$[\partial_t+(p/m_s)\cdot\nabla-q_s\nabla\phi\cdot\nabla_p]f_s=0$，$\phi=\phi_{ext}+e\int C(f_+-f_-)\dd x'$。规范化$\int g_s\dd p=1$给$\nabla^2\phi=-4\pi[\rho_{ext}+e(n_+-n_-)]$。

热平衡$n_\pm=n_0\e^{\mp\beta e\phi}$，弱势线性$e(n_+-n_-)\simeq-2\beta n_0e^2\phi$，于是$(\nabla^2-\lambda^{-2})\phi=-4\pi\rho_{ext}$、$\lambda^2=\kb T/(8\pi n_0e^2)$。解$\phi(r)=\int\e^{-|r-r'|/\lambda}\rho_{ext}(r')/|r-r'|\dd^3r'$，因为$(\nabla^2-\lambda^{-2})(\e^{-r/\lambda}/r)=-4\pi\delta(r)$；仅在$r\ne0$验证微分式会漏掉决定符号的点源。若使用原题$\Phi=-\phi$、$C=-1/r$则Vlasov正负流及$n_\pm=n_0\e^{\pm\beta e\Phi}$转换一致，但Green要取负号。

$n_0\lambda^3\gg1$等价$\lambda\gg\ell=n_0^{-1/3}$，亦即Coulomb能$e^2/\ell\ll\kb T$（常数阶略去），很多粒子共同形成弱平均场。$c_{\rm th}\sim\sqrt{\kb T/m}$，$\lambda/c_{\rm th}\sim\sqrt{m/(n_0e^2)}\sim\omega_p^{-1}$，温度抵消。这是集体响应/振荡时间尺度，碰撞热化或Landau阻尼时间未必同一个常数。
\end{solution}
\begin{exercise}[原题5(a)--(i)]\label{r:333-2-5}
二维经典电子均匀磁场，单体$H=(p-q_eA)^2/(2m)\pm\mu_B|B|$、$A=B\times q/2$，分自旋上、下，守自旋碰撞截面$\sigma_{\uparrow\uparrow}=\sigma_{\downarrow\downarrow}=\sigma,\sigma_{\uparrow\downarrow}=\sigma_\times$。(a)轨道；(b)两Boltzmann式；(c)总H单调；(d)碰撞零项条件；(e)流项的守恒量函数；(f)角动量碰撞内外；(g)圆对称约束的最一般平衡分布；(h)磁/非磁杂质影响；(i)自旋与角动量是否给额外流体方程。
\end{exercise}
\begin{solution}
机械动量$\pi=p-q_eA$，$\dot q=\pi/m$、$\dot\pi=q_e(\pi/m)\times B$，横向回旋频率$|q_eB|/m$，均匀Zeeman常数不改轨道。每组的流算子$L_s$加$C_{ss}$与$C_{s\bar s}$；具体$C_{ss'}=\int\dd\pi_2\dd\Omega\,|v-v_2|\sigma_{ss'}[f_s'f_{s',2}'-f_sf_{s',2}]$，上、下互换得第二式，局域同位置碰撞。

乘$1+\log f_s$求和，把反向散射和两入射粒子交换平均，$\dot H=-\tfrac14\sum_{ss'}\int W(A-B)\log(A/B)\le0$。等号要求$\log f_s$为碰撞不变量线性组合，即两组各自数乘子、共享总动量与能量系数。若还要求平稳流，乘子组合需是单体Hamilton守恒量；任何$f(H_s,L_z)$的Poisson流可为0，但只有合适指数线性组合同时满足碰撞条件。

圆对称时\textbf{总canonical}$L_z=\sum(q\times p)_z$守恒；单个电子碰撞中可交换角动量，不能声称每个$L_i$均守。局域碰撞总机械动量守恒、且$A(q)$相同，故也保持总canonical动量矩和$L_z$。一般热平衡$f_s\propto\exp[-\beta(H_s-\Omega L_z)-a_s]$，两数各固定、共同$\beta,\Omega$；若容器还平移对称可加相应总动量乘子。

固定杂质一般吸收动量与角动量，$\Omega$不再守恒；非磁杂质仍保两自旋数，磁散射可使自旋翻转，相关化学势差驰豫。自旋数差给额外连续/扩散模式，角动量守恒通常由动量方程乘位置导出，不独立增加一个流体变量；在有外部力矩或自旋--轨道耦合时需重新列守恒量。
\end{solution}
\begin{exercise}[原题6(a)--(j)]\label{r:333-2-6}
二维Lorentz气，固定硬圆散射体密度$n_0$、半径$a$。(a)角$\theta\in[0,2\pi]$微分截面$a\sin(\theta/2)\dd\theta/2$与总截面；(b)仅杂质Boltzmann；(c)任意$\chi(|p|)$的矩方程；(d)质量连续；(e)能量方程用$\epsilon=\rho\avg{c^2}/2,c=p/m-u,h=\rho\avg{cc^2}/2,P_{ij}=\rho\avg{c_ic_j}$；(f)$f_0=n\e^{-p^2/(2m\kb T)}/(2\pi m\kb T)$的零阶矩；(g)单碰撞时间一阶$f$；(h)质量流；(i)$\chi_{ij}=\partial u_i/\partial F_j$；(j)压力、热流、一阶流体式。
\end{exercise}
\begin{solution}
冲击参数有符号$b=a\cos(\theta/2)$，全角$|\dd b|=(a/2)\sin(\theta/2)\dd\theta$，总$2a$。若只用无符号偏转$0\le\theta\le\pi$，两侧合计为$a\sin(\theta/2)\dd\theta$，不能把两个约定再双倍。

$Lf=n_0|p|/m\int\dd\theta\,(\dd\sigma/\dd\theta)[f(p')-f(p)]$。因$|p'|=|p|$，任意径向$\chi$的碰撞积分0，分部积分给$\partial_t(n\avg\chi)+\nabla\cdot[n\avg{(p/m)\chi}]=nF\cdot\avg{\nabla_p\chi}$。$\chi=1$得$\partial_t\rho+\nabla\cdot\rho u=0$；能量得
\begin{equation}
\partial_t(\epsilon+\rho u^2/2)+\nabla\cdot[u(\epsilon+\rho u^2/2)+h+P\cdot u]=nF\cdot u.
\end{equation}
纯Lorentz气还保每个$|p|$分布，故需弱粒子碰撞才可用仅数、能的局部Maxwell闭合。零阶$u=h=0$、$P_{ij}=n\kb T\delta_{ij}$、$\epsilon=n\kb T$，各局域数能暂时无驰豫。

令$C=-(f-f_0)/\tau$，准静态一次梯度$f=f_0[1-\tau(p/m)\cdot(\nabla\log\rho-\nabla\log T+p^2\nabla T/(2m\kb T^2)-F/(\kb T))]$。Gaussian矩$\avg{p_ip_j}=s\delta_{ij}$、$\avg{p_ip_jp^2}=4s^2\delta_{ij}$、$\avg{p_ip_jp^4}=24s^3\delta_{ij}$，$s=m\kb T$。所以
\begin{equation}
\rho u=n\tau[F-\kb T\nabla\log(\rho T)],\quad \chi_{ij}=\tau\delta_{ij}/m,\quad P_{ij}=n\kb T\delta_{ij}+O(\nabla^2).
\end{equation}
实验室能流$Q=2n\kb Tu-2n\tau\kb^2T\nabla T/m$。转到题定义的\textbf{peculiar}热流要减$2n\kb Tu$，故$h=-2n\tau\kb^2T\nabla T/m$。把这些代上述质量、能量式就是一阶流体闭合，无独立保动量Euler方程。

官方解题6(j)把二维Wick六阶系数写22并给出有量纲问题的热流；独立以$\avg{p^6}=2\cdot4\cdot6\,s^3=48s^3$、分量各半验证应为24，且必须减对流焓流。本题据矩积分修复，完整有限数值核对随QA给出，不能将官方误植再当正确结果。
\end{solution}
\begin{exercise}[原题7(a)--(e)]\label{r:333-2-7}
三维经典气夹于$y=0,w$两板，温$T_1,T_2$，无宏观速，$f_0=n(y)\e^{-p^2/(2m\kb T(y))}/[2\pi m\kb T(y)]^{3/2}$。(a)$n,T$必要关系；(b)证明$\avg{p^2}=3s,\avg{p^4}=15s^2,\avg{p_y^2p^4}=35s^3$，可用$\avg{p^6}=105s^3$；(c)一碰撞时长$\tau_K$的$f_1$；(d)热导率；(e)稳态$T(y)$。
\end{exercise}
\begin{solution}
无力动量平衡要求压强$n\kb T$常数，因此$n'/n=-T'/T$。独立Gaussian三分量的Wick矩或$\chi^2_3$径向积分给3、15、105，旋转均分六阶为35。流$f_1=f_0-\tau_K(p_y/m)\partial_yf_0=f_0[1-\tau_K(p_y/m)(T'/T)(p^2/(2m\kb T)-5/2)]$，该5/2保证质量流0。

$h_y=\int(p_y/m)(p^2/2m)f_1\dd^3p$，五阶项奇为0，六阶和四阶差给$h_y=-5n\tau_K\kb^2T\,T'/(2m)$，热导$\kappa=5n\tau_K\kb^2T/(2m)$。稳态热流常量，所以$\int_{T_1}^{T(y)}\kappa(T)\dd T=(y/w)\int_{T_1}^{T_2}\kappa(T)\dd T$。若题模型$\tau_K$为常数且$nT$常数，$\kappa$恒定，温度线性；若硬球实际$\tau_K\propto1/(n\sqrt T)$，则$\kappa\propto\sqrt T$，$T(y)^{3/2}$线性。原题未规定真实碰撞截面的温变时，不能不加条件宣称唯一线性或非线性剖面。
\end{solution}

\originalchapter{8.333复习包3：关联气体与量子正则系综}
\osource{8.333 Fall 2013页面公开的Fall 2007 Test3 Review及Solutions，10题。题3、5对应已有卷内完整题解；同卷重复编号保留映射，不重复抄页。}
\begin{exercise}[原题1(a)--(g)]\label{r:333-3-1}
$N$电子及均匀正背景，$U_Q=\sum_{i<j}V(q_i-q_j)$、$V(r)=e^2/(4\pi r)-c$，背景令$\int V=0$。(a)Fourier非零模$e^2/q^2$、零模0；(b)证明$l$环积分$R_l=V^{1-l}\int_q\widetilde V(q)^l$；(c)累积量中环计数$S_l\simeq(l-1)!N^l/2$；(d)环求和；(e)积分得$\log Z-\log Z_0=V\kappa^3/(12\pi)$，$\kappa^2=\beta e^2N/V$；(f)压强；(g)两粒子间无环链的求和屏蔽势。
\end{exercise}
\begin{solution}
Coulomb的$-\nabla^2(1/4\pi r)=\delta$给$e^2/q^2$；背景专门去掉$q=0$，没有正负长程项残余。环卷积将各边动量固定为同一$q$，一个整体位置积分$V$，即(b)。选$l$不同顶点$\binom Nl$、无向环排列$(l-1)!/2$，累积量展开还含$l!$次边次序，结合给$N!/(N-l)!\,(l-1)!/2$；$l=2$重复同一对边直接验证$N(N-1)/2$。

$\sum_{l\ge2}(-x)^l/(2l)=[x-\log(1+x)]/2$，故
\begin{equation}
\Delta\log Z=\frac V2\int_q\left[\frac{\kappa^2}{q^2}-\log\left(1+\frac{\kappa^2}{q^2}\right)\right].
\end{equation}
把$\kappa$微分，$\partial_\kappa(\Delta\log Z/V)=\kappa\int_q[1/q^2-1/(q^2+\kappa^2)]=\kappa^2/(4\pi)$；在$\kappa=0$定积分常数0即$\kappa^3/(12\pi)$。固定$N,T$有$\kappa\propto V^{-1/2}$，压力修正$\Delta P=\kb T\partial_V\Delta\log Z=-\kb T\kappa^3/(24\pi)$，非解析$n^{3/2}$，不可用有限次数整数virial展开替代。

链边的几何级数$\widetilde V_{scr}=\widetilde V/(1+\beta n\widetilde V)=e^2/(q^2+\kappa^2)$，逆变换$V_{scr}=e^2\e^{-\kappa r}/(4\pi r)$。若题末写无$e^2$，那是势除以耦合后的归一化；本解保持原能量单位。环近似用弱Coulomb耦合，不声称所有密度下精确。
\end{solution}
\begin{exercise}[原题2(a)--(d)]\label{r:333-3-2}
$d$维阶梯势$V=\infty$（$r<a$）、$-\varepsilon$（$a<r<b$）、0（$r>b$）。(a)$B_2$及高低温；(b)压缩率首修正；(c)高温改写van der Waals式和参数；(d)$b=a,d=1$的$B_3$。
\end{exercise}
\begin{solution}
球体积$\Omega_dR^d$、$\Omega_d=\pi^{d/2}/\Gamma(d/2+1)$。$f=-1$于核、$\e^{\beta\varepsilon}-1$于壳，故$B_2=\Omega_d[a^d-(b^d-a^d)(\e^{\beta\varepsilon}-1)]/2$。高温趋核体积半值，低温吸引使它大负；温度低到关联团簇多，稀薄截断失控。

$P=n\kb T(1+B_2n+\cdots)$，$\kappa_T=[n\kb T(1+2B_2n)]^{-1}\simeq[1-2B_2n]/(n\kb T)$。高温$B_2=b_{vdW}-a_{vdW}/(\kb T)$，$b_{vdW}=\Omega_da^d/2$、$a_{vdW}=\Omega_d\varepsilon(b^d-a^d)/2$，重求和$P\simeq n\kb T/(1-b_{vdW}n)-a_{vdW}n^2$仅匹配到$n^2$。一维硬棒精确$P=n\kb T/(1-na)$，$B_3=a^2$，见习题\ref{p:334-4-4}。
\end{solution}
\originalsection{原题3：Dieterici临界点的完整引用}
原题$P(v-b)=\kb T\exp[-a/(\kb Tv)]$三问：临界压缩因子、$v_c$上压缩率、临界等温线最低非零项，恰为习题\ref{p:334-1-1}(e)--(g)。所有条件及解答均在该处，复习编号保留为333-R3-3。
\begin{exercise}[原题4(a)--(e)]\label{r:333-3-4}
二维$N$个$+c_0$、$N$个$-c_0$经典粒子，面积$A=L^2$，$H=\sum_{1}^{2N}p_i^2/(2m)-\sum_{i<j}c_ic_j\log(|q_i-q_j|/r_0)$。(a)吸、斥对数量；(b)$Z$积分并化成距离幂；(c)用$q'=q/L$证明$Z\propto A^{2N-\beta Nc_0^2/2}$；(d)二维压力与低温问题；(e)加入硬核$a$，由$N=1$估其重要温度与相的性质。
\end{exercise}
\begin{solution}
同号斥对$2\binom N2=N(N-1)$，异号吸对$N^2$。动量给$\lambda^{-4N}$，$Z=[(N!)^2\lambda^{4N}]^{-1}\int_A\prod\dd^2q_i\prod_{i<j}(|q_i-q_j|/r_0)^{\beta c_ic_j}$，固定$r_0$只是指定能零。总对积电荷$\sum_{i<j}c_ic_j=[(\sum c)^2-\sum c_i^2]/2=-Nc_0^2$。坐标缩放使测度$L^{4N}$、距离幂$L^{-\beta Nc_0^2}$，即题示面积幂，前提积分收敛。

因此$P=N(2\kb T-c_0^2/2)/A$，高温回到$2N$粒子理想压，形式$\kb T<c_0^2/4$变负。但没有核的异号吸引在更早的短距条件$\beta c_0^2\ge2$便使两体积分$\int_0^Lr^{1-\beta c_0^2}\dd r$发散，形式负压分支并不存在稳定配分函数。

硬核下$\int_a^Lr^{1-\beta c_0^2}\dd r$在$\beta c_0^2<2$由大距控制、$>2$由核控制，边界对数。所以单对估计温标$\kb T_*=c_0^2/2$，低温偏好束缚中性对、高温自由电荷。有限密度的集体涡旋/Coulomb气转变需屏蔽与逸度RG，单对核重要尺度不等于已经精确证明多体BKT温度。
\end{solution}
\originalsection{原题5：最近邻一维气的完整引用}
排序粒子硬核最近邻势，原题(a)--(f)逐项对应习题\ref{p:334-4-4}：Hamilton量、间距域、Laplace鞍点、边界间隙积分、平均长度、硬棒方程与$B_l$。原复习题另问排除体积比较：一维$b_{vdW}=a$与精确硬棒结果相合；有限维一般van der Waals排除重求和并非精确。
\begin{exercise}[原题6(a)--(d)]\label{r:333-3-6}
$N+1$质量$m$粒子以$N$弹簧$K$、原长$a$连成链，端点固定相距$Na$，位移$u_0=u_N=0$，$H=\sum_{1}^{N-1}p_i^2/(2m)+K\sum_{i=0}^{N-1}(u_{i+1}-u_i)^2/2$。(a)正弦正规模与频率；(b)经典$Z$；(c)$\avg{|\widetilde u_k|^2},\avg{u_i^2}$并画位置曲线；(d)仅左端固定、右端自由如何改变。
\end{exercise}
\begin{solution}
$u_i=\sqrt{2/N}\sum_{r=1}^{N-1}\widetilde u_r\sin(\pi ri/N)$是实正交变换；频率$\omega_r=2\sqrt{K/m}\sin(\pi r/(2N))$。独立模$H_r=\widetilde p_r^2/(2m)+m\omega_r^2\widetilde u_r^2/2$，$Z=\prod_r\kb T/(\hbar\omega_r)=[\kb T/(\hbar\sqrt{K/m})]^{N-1}/\sqrt N$，乘积恒等式$\prod_{r=1}^{N-1}2\sin(\pi r/(2N))=\sqrt N$。

$\avg{\widetilde u_r^2}=\kb T/(m\omega_r^2)$；离散Laplace逆矩阵给$\avg{u_iu_j}=(\kb T/K)[\min(i,j)-ij/N]$，所以$\avg{u_i^2}=(\kb T/K)i(N-i)/N$，中部最大，随链长线性增长。
\begin{center}\begin{tikzpicture}\begin{axis}[width=9cm,height=4cm,xlabel={$i/N$},ylabel={$K\avg{u_i^2}/(N\kb T)$},xmin=0,xmax=1,ymin=0]\addplot[thick,structurecolor,domain=0:1]{x*(1-x)};\end{axis}\end{tikzpicture}\end{center}
自由右端保留$u_1,\ldots,u_N$，独立弹簧伸长$\delta_i=u_i-u_{i-1}$，$\avg{\delta_i\delta_j}=\kb T\delta_{ij}/K$，故$\avg{u_iu_j}=\kb T\min(i,j)/K$、$Z=[\kb T/(\hbar\sqrt{K/m})]^N$。频率$2\sqrt{K/m}\sin[(2r-1)\pi/(4N+2)]$，$r=1,\ldots,N$。原提示说$N-1$伸长是沿用固定端自由度的索引；若终点真自由，共有$N$个移动粒子与伸长，不能漏最后一个模。
\end{solution}
\begin{exercise}[原题7(a)--(g)]\label{r:333-3-7}
黑洞熵$S=\kb c^3A/(4G\hbar)$。(a)牛顿逃逸速设为$c$得半径；(b)两个太阳质量$M_\odot\simeq2\times10^{30}$ kg黑洞合并熵增、比特数；(c)$E=Mc^2$给温度；(d)黑体辐射功；(e)蒸发时长和太阳质量估值；(f)与题给2.7 K背景平衡质量；(g)半径$R$球内全息熵上界。采用题设无自旋黑体理想模型，不加未给灰体和额外粒子种类。
\end{exercise}
\begin{solution}
$v_{esc}^2=2GM/R$，$R_s=2GM/c^2$，$A=16\pi G^2M^2/c^4$，所以$S=4\pi\kb GM^2/(\hbar c)$。忽略合并辐射，末质量$2M$，熵增$8\pi\kb GM^2/(\hbar c)$，约$3.1\times10^{77}$ bits（除$\kb\log2$）；真实计算应包括辐射能与熵，不能把无损质量相加作为精确天体结果。

$1/T=\partial S/\partial E$给$T_H=\hbar c^3/(8\pi\kb GM)$，太阳质量约$6.1\times10^{-8}$ K。用黑体$P=\sigma_{SB}AT_H^4=\hbar c^6/(15360\pi G^2M^2)$，$\dot Mc^2=-P$积分得到$t_{ev}=5120\pi G^2M_0^3/(\hbar c^4)$，太阳质量约$2.2\times10^{67}$年。末期黑体半经典模型失效，此巨大尺度只作原题估算。

$T_H=2.7$ K给$M\simeq4.5\times10^{22}$ kg。半径$R$所容最大熵按面积界为$S_{max}=\pi\kb c^3R^2/(G\hbar)$，体现面积而非普通体积增长；界的引力和几何前提属于题设全息原理，不由经典相空间计数单独证明。
\end{solution}
\begin{exercise}[原题8(a)--(h)]\label{r:333-3-8}
量子谐振子$H=p^2/(2m)+m\omega^2q^2/2$。(a)$Z,U$；(b)本征基密度矩阵；(c)非对易算符指数导数与迹导数；(d)用$\partial_mZ=0$证动势能相等；(e)$\avg{q^2}$及链低温量子修正；(f)高温坐标核；(g)低温核；(h)完整精确坐标核。
\end{exercise}
\begin{solution}
$b=\beta\hbar\omega$，$Z=\sum_{n\ge0}\e^{-b(n+1/2)}=[2\sinh(b/2)]^{-1}$，$U=(\hbar\omega/2)\coth(b/2)$，$\rho=Z^{-1}\sum_n\e^{-\beta E_n}\ket n\bra n$。

任意$A(x)$有$\partial_x\e^A=\int_0^1\e^{(1-s)A}A'\e^{sA}\dd s$，由指数乘积展开或Duhamel式导出，只有$[A,A']=0$才可写$A'\e^A$。迹内循环却始终给$\partial_x\Tr\e^A=\Tr(A'\e^A)$。固定$\omega$时$Z$与$m$无关，所以$0=\avg{\partial_mH}=\avg{-p^2/(2m^2)+\omega^2q^2/2}$，两能各$U/2$，$\avg{q^2}=\hbar\coth(b/2)/(2m\omega)$。

链的每模将$\kb T/(m\omega_r^2)$换为$\hbar\coth(\beta\hbar\omega_r/2)/(2m\omega_r)$，再按正弦模相加；零温仍有零点涨落，一维长波$1/|k|$给对数链长而非经典线性增长。

完整核可从热核方程推导：设$\langle q|\e^{-\beta H}|q^\prime\rangle$的非对角一般形式为$A_\beta\exp[-a_\beta(q^2+q^{\prime2})+b_\beta qq^\prime]$。代$-\partial_\beta K=H_qK$并比较多项式系数，得$a_\beta'=m\omega^2/2-2\hbar^2a_\beta^2/m$、$b_\beta'=-2\hbar^2a_\beta b_\beta/m$、$A_\beta'/A_\beta=-\hbar^2a_\beta/m$。$\beta\to0$的自由Gaussian $\delta$核选出$a=m\omega\coth b/(2\hbar)$、$b_\beta=m\omega/(\hbar\sinh b)$、$A=\sqrt{m\omega/(2\pi\hbar\sinh b)}$。再除$Z$、重写中点和差坐标，给
\begin{equation}
\rho(q',q)=\sqrt{\frac{m\omega}{\pi\hbar}\tanh\frac b2}\exp\left[-\frac{m\omega}{\hbar}\tanh\frac b2\left(\frac{q+q'}2\right)^2-\frac{m\omega}{4\hbar}\coth\frac b2(q-q')^2\right].
\end{equation}
高温$b\ll1$是$\sqrt{\beta m\omega^2/(2\pi)}\exp[-\beta m\omega^2(q+q')^2/8-m(q-q')^2/(2\beta\hbar^2)]$到领先阶；低温$b\gg1$是$\psi_0(q')\psi_0(q)$，$\psi_0=(m\omega/\pi\hbar)^{1/4}\e^{-m\omega q^2/(2\hbar)}$。沿对角积分1、对角二阶矩等于(e)，为独立归一化检查。
\end{solution}
\begin{exercise}[原题9(a)--(f)]\label{r:333-3-9}
三维$N$正、$N$负相对论量子粒子$H=c\sum_i|p_i|+\sum_{i<j}e_ie_j/|r_i-r_j|$，立方边$L$。(a)Schrödinger式与Bose/Fermi对称；(b)$E_n(L)$缩放；(c)$Z$的$T,V$组合；(d)$E,P$精确关系；(e)$d$维Coulomb$r^{2-d}$配非相对论$p^2/(2m)$何维同样齐次及关系；(f)为何密集真实Coulomb混合物可能不服这些式。
\end{exercise}
\begin{solution}
$[c\sum_i\sqrt{-\hbar^2\nabla_i^2}+\sum e_ie_j/r_{ij}]\Psi_n=E_n\Psi_n$，同正电荷组、同负组分别按全粒子含自旋的Bose对称或Fermi反对称，不对不同电荷强行交换。$r_i=Lr_i'$、规范化$\Psi_L=L^{-3N}\Psi_1$使每动能、势能均$L^{-1}$，故$E_n(L)=E_n(1)/L$，前提Hamilton量与边界确无另一长度且稳定谱存在。

$Z=\sum_n\exp[-E_n(1)/(L\kb T)]$仅依赖$LT$，$E=-\partial_\beta\log Z$，$P=\kb T\partial_V\log Z=E/(3V)$。非相对论动能$L^{-2}$，须$d-2=2$即$d=4$势也$L^{-2}$，于是$E=dPV/2=2PV$。

\textbf{同一个稳定齐次Hamilton模型}不会仅因变液体或固体就丧失这个形式恒等式。真实密集材料的核尺寸、电子质量/非相对论转换、排斥核、有效相互作用与环境引入额外尺度；异号点吸引也可能导致塌缩、无有限$Z$。是这些前提变化限制缩放，不是“高密度”三个字自动使数学证明失效。
\end{solution}
\begin{exercise}[原题10(a)--(d)]\label{r:333-3-10}
一般$H(p,q)$。(a)经典正则配分在一对$q_1\to\lambda q_1,p_1\to p_1/\lambda$的canonical变换下不变；(b)量子对应；(c)$H=\sum p_i^2/(2m)+V(q)$证明$\avg{p_1^2/m}=\avg{q_1\cdot\nabla_1V}$；(d)虚构G-8.333星系边缘速200 km/s，估质量/尺度比。
\end{exercise}
\begin{solution}
相空间全空间或连同边界一并变换时Jacobian$\lambda^d\lambda^{-d}=1$，$Z[H_\lambda]=Z[H]$。量子用幺正dilation$U_\lambda$保持$[q,p]=\ii\hbar$，$H_\lambda=U_\lambda H U_\lambda^{-1}$谱不变，$\Tr\e^{-\beta H_\lambda}$不变。不能在固定硬盒里缩坐标却忽略域/墙项；这时边界压力会贡献virial。

在$\lambda=1$微分$\log Z$，迹Duhamel式适用于量子非对易，$0=-\beta\avg{-p_1^2/m+q_1\cdot\nabla_1V}$，即题式。星系$Mv^2\sim GM^2/R$给$M/R\sim v^2/G\simeq6.0\times10^{20}$ kg/m，约$9.3\times10^9M_\odot$每kpc。质量分布、轨道各向异性与边缘速定义带来常数阶差，题给速度只能估尺度，不确定完整质量剖面。
\end{solution}

\originalchapter{8.333期末复习包：七个历史年份的量子与相变题}
\osource{Fall 2013页面的2007 Final Review及官方Solutions，内含Fall1998、1999、2000、2003、2004、2005、2006各三题，共21个原编号。没有2001、2002卷，不补造；这些是包内标明年份的历史考试，不是2013年实际期末卷。}
\originalsection{1998年：激子与离子凝聚}
\begin{exercise}[1998题1(a)--(e)]\label{f:333-1998-1}
半导体光激发$N_e$电子、$N_h$空穴，质量$m_e,m_h$，可结合成质量$M=m_e+m_h$、降低能$\varepsilon$的激子，均按经典非相互作用气、忽略内部简并。(a)自由电子空穴自由能；(b)$N_p$激子自由能；(c)三化学势；(d)平衡条件；(e)高温、总自由激发密度$n\simeq n_e+n_h$时的激子密度，取电中性。
\end{exercise}
\begin{solution}
$\lambda_j=h/\sqrt{2\pi m_j\kb T}$，$F_e+F_h=\kb T\sum_{j=e,h}N_j[\log(n_j\lambda_j^3)-1]$，$F_p=N_p\{\kb T[\log(n_p\lambda_p^3)-1]-\varepsilon\}$。所以$\mu_e=\kb T\log(n_e\lambda_e^3)$、$\mu_h$类同、$\mu_p=-\varepsilon+\kb T\log(n_p\lambda_p^3)$。

电子加空穴变一激子，平衡$\mu_e+\mu_h=\mu_p$，即$n_p=n_en_h(\lambda_e^3\lambda_h^3/\lambda_p^3)\e^{\beta\varepsilon}=n_en_h\lambda_\mu^3\e^{\beta\varepsilon}$，$\mu=m_em_h/(m_e+m_h)$为约化质量。高温解离多，$n_e=n_h\simeq n/2$，$n_p\simeq n^2\lambda_\mu^3\e^{\beta\varepsilon}/4$。更低温需用激发对守恒$n_{pair}=n_e+n_p=n_h+n_p$解二次方程，不能保持$n_e=n/2$近似到全结合区。
\end{solution}
1998题2(a)--(e)Manning题逐项在既有习题33.3。此历史题用固定$\log(r/L)$而PS5用$\log(r/R)$：只把单粒子$z$的前因子由$R^{2\xi}$换成$L^{2\xi}$，径向概率、均值、凝聚反馈不变，物理壁压按固定能零接触式计算。1998题3(a)--(g)逐项对应本卷习题\ref{p:333-6-11}；热容全温峰与二维极限的修正同样适用，官方三维临界峰公式不能外推成所有维数的全温最大值。
\originalsection{1999年：Fermi与Bose磁性}
\begin{exercise}[1999题1(a)--(e)]\label{f:333-1999-1}
密度$n=N/V$的spin1/2自由电子。(a)基态、$k_F,E_0/V$；加入$U=\alpha N_+N_-/V$、$\alpha=4\pi\hbar^2a/m$模拟偏好平行自旋。(b)两Fermi波数；(c)$n_\pm=n/2\pm\delta$的小偏极动能；(d)相互作用与Stoner$\alpha_c$；(e)自发磁化随$\alpha$及图。
\end{exercise}
\begin{solution}
基态每自旋占满动量球，$k_F=(3\pi^2n)^{1/3}$、$E_F=\hbar^2k_F^2/(2m)$，$E_0/V=3nE_F/5$。偏极$k_{F\pm}=(6\pi^2n_\pm)^{1/3}$，动能密度$\tfrac35[\hbar^2(6\pi^2)^{2/3}/(2m)](n_+^{5/3}+n_-^{5/3})$，至$\delta^2$增加$4E_F\delta^2/(3n)$。相互作用$\alpha(n^2/4-\delta^2)$，故$\alpha_c=4E_F/(3n)$。

用$m_s=2\delta/n$，能至四阶$E/V=\mathrm{const}+(nE_F/3)(1-\alpha/\alpha_c)m_s^2+(nE_F/81)m_s^4+\cdots$，近临界$m_s\simeq\sqrt{(27/2)(\alpha/\alpha_c-1)}$。大耦合需最小化完整$5/3$次式，磁化最后饱和$|m_s|=1$；饱和边界$\alpha/\alpha_c=3\,2^{2/3}/4\simeq1.19055$。下图用完整方程，未将小磁化展开外推到饱和值。
\begin{center}\begin{tikzpicture}\begin{axis}[width=10cm,height=4.5cm,xlabel={$\alpha/\alpha_c$},ylabel={$|m_s|$},xmin=.8,xmax=1.4,ymin=0,ymax=1.05]\addplot[thick,structurecolor] coordinates {(0.8,0) (0.803351955,0) (0.806703911,0) (0.810055866,0) (0.813407821,0) (0.816759777,0) (0.820111732,0) (0.823463687,0) (0.826815642,0) (0.830167598,0) (0.833519553,0) (0.836871508,0) (0.840223464,0) (0.843575419,0) (0.846927374,0) (0.85027933,0) (0.853631285,0) (0.85698324,0) (0.860335196,0) (0.863687151,0) (0.867039106,0) (0.870391061,0) (0.873743017,0) (0.877094972,0) (0.880446927,0) (0.883798883,0) (0.887150838,0) (0.890502793,0) (0.893854749,0) (0.897206704,0) (0.900558659,0) (0.903910615,0) (0.90726257,0) (0.910614525,0) (0.91396648,0) (0.917318436,0) (0.920670391,0) (0.924022346,0) (0.927374302,0) (0.930726257,0) (0.934078212,0) (0.937430168,0) (0.940782123,0) (0.944134078,0) (0.947486034,0) (0.950837989,0) (0.954189944,0) (0.957541899,0) (0.960893855,0) (0.96424581,0) (0.967597765,0) (0.970949721,0) (0.974301676,0) (0.977653631,0) (0.981005587,0) (0.984357542,0) (0.987709497,0) (0.991061453,0) (0.994413408,0) (0.997765363,0) (1.00111732,0.12245661) (1.00446927,0.242773006) (1.00782123,0.318360244) (1.01117318,0.377207494) (1.01452514,0.42635704) (1.01787709,0.468914866) (1.02122905,0.50658661) (1.02458101,0.540432723) (1.02793296,0.571168223) (1.03128492,0.599303719) (1.03463687,0.625219856) (1.03798883,0.64921001) (1.04134078,0.671506405) (1.04469274,0.692296914) (1.04804469,0.711736305) (1.05139665,0.729954045) (1.0547486,0.74705987) (1.05810056,0.763147843) (1.06145251,0.7782994) (1.06480447,0.792585657) (1.06815642,0.806069196) (1.07150838,0.818805468) (1.07486034,0.830843905) (1.07821229,0.842228815) (1.08156425,0.85300011) (1.0849162,0.863193903) (1.08826816,0.872843001) (1.09162011,0.881977314) (1.09497207,0.890624206) (1.09832402,0.89880878) (1.10167598,0.906554133) (1.10502793,0.913881562) (1.10837989,0.92081075) (1.11173184,0.927359921) (1.1150838,0.933545975) (1.11843575,0.939384606) (1.12178771,0.944890405) (1.12513966,0.950076942) (1.12849162,0.954956849) (1.13184358,0.959541879) (1.13519553,0.963842963) (1.13854749,0.967870256) (1.14189944,0.971633169) (1.1452514,0.975140398) (1.14860335,0.97839994) (1.15195531,0.981419091) (1.15530726,0.984204443) (1.15865922,0.986761844) (1.16201117,0.989096351) (1.16536313,0.991212125) (1.16871508,0.993112282) (1.17206704,0.994798638) (1.17541899,0.99627127) (1.17877095,0.997527735) (1.18212291,0.99856146) (1.18547486,0.999357859) (1.18882682,0.999880604) (1.19217877,1) (1.19553073,1) (1.19888268,1) (1.20223464,1) (1.20558659,1) (1.20893855,1) (1.2122905,1) (1.21564246,1) (1.21899441,1) (1.22234637,1) (1.22569832,1) (1.22905028,1) (1.23240223,1) (1.23575419,1) (1.23910615,1) (1.2424581,1) (1.24581006,1) (1.24916201,1) (1.25251397,1) (1.25586592,1) (1.25921788,1) (1.26256983,1) (1.26592179,1) (1.26927374,1) (1.2726257,1) (1.27597765,1) (1.27932961,1) (1.28268156,1) (1.28603352,1) (1.28938547,1) (1.29273743,1) (1.29608939,1) (1.29944134,1) (1.3027933,1) (1.30614525,1) (1.30949721,1) (1.31284916,1) (1.31620112,1) (1.31955307,1) (1.32290503,1) (1.32625698,1) (1.32960894,1) (1.33296089,1) (1.33631285,1) (1.3396648,1) (1.34301676,1) (1.34636872,1) (1.34972067,1) (1.35307263,1) (1.35642458,1) (1.35977654,1) (1.36312849,1) (1.36648045,1) (1.3698324,1) (1.37318436,1) (1.37653631,1) (1.37988827,1) (1.38324022,1) (1.38659218,1) (1.38994413,1) (1.39329609,1) (1.39664804,1) (1.4,1)};\end{axis}\end{tikzpicture}\end{center}

这是题给平均场接触模型的失稳，不是实际电子关联体系的精确相图。
\end{solution}
\begin{exercise}[1999题2(a)--(g)]\label{f:333-1999-2}
spin1理想Bose气$\epsilon_{k,s}=\hbar^2k^2/(2m)-\mu_0sB$，$s=-1,0,+1$，忽略轨道磁性。(a)占据；(b)各$N_s$；(c)$M,\chi(T,\mu)$；(d)零场稀薄$z(n)$至$n^2$与磁化率修正；(e)$T_c(n,0)$和上侧$\chi$；(f)凝聚以下小$B$的$\mu$及凝聚态；(g)$B\downarrow0$自发磁化。
\end{exercise}
\begin{solution}
$z_s=z\e^{s\beta\mu_0B}$，$\avg{n_{k,s}}=[z_s^{-1}\e^{\beta\epsilon_k}-1]^{-1}$，$N_s=V\lambda^{-3}\operatorname{Li}_{3/2}(z_s)$（凝聚时另加最低态）。$M=\mu_0(N_+-N_-)$，$\chi_{B=0}=2\beta\mu_0^2V\lambda^{-3}\operatorname{Li}_{1/2}(z)$。

$x=n\lambda^3/3$，数式$x=z+z^2/2^{3/2}+\cdots$给$z=x-x^2/2^{3/2}$，所以$\chi=2N\mu_0^2[1+x/2^{3/2}+\cdots]/(3\kb T)$。$T_c=(2\pi\hbar^2/m\kb)[n/(3\zeta(3/2))]^{2/3}$，$z\to1$时$\operatorname{Li}_{1/2}$发散；数式$\operatorname{Li}_{3/2}(z)=\zeta(3/2)-2\sqrt\pi\sqrt{-\log z}+\cdots$给固定密度$\chi\propto(T-T_c)^{-1}$。

$B>0$最低态为$k=0,s=+1$，凝聚$\mu=-\mu_0B$。先热力学极限再$B\downarrow0$，三个激发组相同而凝聚全在$+$，$M=\mu_0N[1-(T/T_c)^{3/2}]$。严格有限体系先置$B=0$的对称混合给$M=0$；极限次序和固定粒子数条件必须注明。
\end{solution}
1999题3(a)--(d)全部为习题\ref{r:333-3-10}的经典/量子virial、星系估算，条件及每个小问相同。
\originalsection{2000年：氦3熔化与Bose占据}
\begin{exercise}[2000题1(a)--(f)]\label{f:333-2000-1}
He3熔化线低于0.3 K斜率负，0.1 K约$-30$ atm/K。(a)固相每核spin1/2、忽略核自旋作用的熵；(b)液相理想Fermi气每原子体积46 Å$^3$，估$T_F$；(c)低温热容；(d)熵与两相比较；(e)由$\mu$相等证明CC式；(f)$v_\ell-v_s=3$ Å$^3$每原子时低温斜率。
\end{exercise}
\begin{solution}
无耦合核spin各2态，$s_s=\kb\log2$。液体$n=1/(46\times10^{-30}\mathrm{m^3})$，裸He3质量约$5.0\times10^{-27}$ kg给$T_F=\hbar^2(3\pi^2n)^{2/3}/(2m\kb)\simeq6.1$ K。$C_V/N=(\pi^2/2)\kb T/T_F$、$s_\ell=(\pi^2/2)\kb T/T_F$，足够低温液体熵小于核自旋未冻结的固体。

$\dd\mu=-s\dd T+v\dd P$、两相$\dd\mu$相等，$\dd P/\dd T=(s_\ell-s_s)/(v_\ell-v_s)$。在所述低温窗口忽略液体熵，斜率$-\kb\log2/(3\times10^{-30}\mathrm{m^3})\simeq-3.19\times10^6$ Pa/K，约$-31.5$ atm/K，与题给尺度相合。核交换尺度以下固相spin也不能保持$\kb\log2$，故此“低温”是低于$T_F$且高于极小spin作用尺度，不可用于真实严格$T=0$。
\end{solution}
\begin{exercise}[2000题2(a)--(f)]\label{f:333-2000-2}
理想Bose单粒子态$q$、能$E_q$、$z_q=\e^{\beta(\mu-E_q)}<1$。(a)所有$n_q$的联合概率；(b)一模特征函数；(c),(d)均值、方差及方差对均值；(e)以均占据重写联合概率；(f)熵与零温。
\end{exercise}
\begin{solution}
每模几何分布，$P(\{n_q\})=\prod_q(1-z_q)z_q^{n_q}$、$\varphi_q(k)=(1-z_q)/(1-z_q\e^{\ii k})$，导数得$\bar n_q=z_q/(1-z_q)$、$\operatorname{Var}n_q=z_q/(1-z_q)^2=\bar n_q(1+\bar n_q)$。改写$P=\prod_q\bar n_q^{n_q}/(1+\bar n_q)^{1+n_q}$，熵
\begin{equation}
S=\kb\sum_q[(1+\bar n_q)\log(1+\bar n_q)-\bar n_q\log\bar n_q].
\end{equation}
无宏观激发时体熵趋0；凝聚基态的巨正则几何数涨落可仍巨大，单模熵约$\kb\log N_0$而每粒子熵0。固定总$N$的纯基态熵0。必须区分巨正则的基态粒子数不确定与有限温激发熵，不能由温度0无条件断言任何混合密度矩阵熵都0。
\end{solution}
2000题3(a)--(d)二维硬杆恰为既有热统习题33.5。临界角由等熵与驻值共同给，$\theta_c=1.19918$、$n_cl^2=0.265566$；驻值折点$1.72067,0.239624$为亚稳界，不误作全局转变。
\originalsection{2003、2006年：氦4；2003、2005年：Dirac激发}
\begin{exercise}[2003题1及2006题1，均(a)--(e)]\label{f:333-he4}
He4低温压致固化，熔化压力在0.8 K附近比0 K低约20 Pa。(a)CC式；(b)液相一支声速$c$的声子热容；(c)固相纵$c_L$及两横$c_T$的热容；(d)$c\simeq c_L\simeq c_T$、$v_\ell\simeq v_s=v$时熵差与哪个熵高；(e)恒小$\delta v=v_\ell-v_s$下熔化曲线与较高密度相。
\end{exercise}
\begin{solution}
CC式已在前题从$\mu$证明。每支三维声子态密度$V\omega^2/(2\pi^2c_s^3)$，$U=\int\hbar\omega\rho/(\e^{\beta\hbar\omega}-1)\dd\omega$给$C/N=2\pi^2v\kb^4T^3/(15\hbar^3c_s^3)$。固相将$c_s^{-3}$换成$c_L^{-3}+2c_T^{-3}$；熵积分$C/T$把系数除3再乘$T^3$。

若各声速和体积近同，记单支热容系数$a=2\pi^2v\kb^4/(15\hbar^3c^3)$，$s_\ell-s_s=-2aT^3/3$，固相比液体声子支更多、熵较高。因此$P(T)=P(0)-aT^4/(6\delta v)$。题述初始下降要求$\delta v>0$，固体体积小、密度高。该纯声子模型只解释低温下降趋势，不预测曲线在0.8 K反转的全部机制，也不在未给$c,v,\delta v$时伪算20 Pa幅度。
\end{solution}
2003题2(a)--(f)表面活性剂与既有习题33.1完全相同，原条件、完整解、塌缩/相分离近似及定张力热容均在那里。
\begin{exercise}[2003题3(a)--(f)；2005题1(a)--(e)]\label{f:333-dirac}
2003三维spin1/2 Dirac谱$E_\pm(k)=\pm\sqrt{m^2c^4+\hbar^2c^2k^2}$，零温负能满、正能空。(a)占据$\mu+\delta$与空穴$\mu-\delta$概率；(b)$\mu(T)$；(c)粒子空穴能积分；(d),(e)无质量能与热容；(f)有质量低温。2005石墨烯二维谱$E_\pm=\pm\hbar v|k|$、题给自旋$g=2,v=10^6$ m/s：(a)--(d)重复二维能和热容；(e)声速$2\times10^4$ m/s的声子与电子比较。按该题印出的$g=2$计，不再擅加另一个谷因子。
\end{exercise}
\begin{solution}
Fermi函数$f(\mu+\delta)=1-f(\mu-\delta)$，对称谱固定半填充$\mu=0$。粒子空穴因子2、自旋$g=2$，$\Delta U=2gV\int\dd^3k/(2\pi)^3\,E_+/(\e^{\beta E_+}+1)$。无质量径向积分$\int_0^\infty x^3/(\e^x+1)\dd x=7\pi^4/120$，
\begin{equation}
\Delta U/V=\frac{7g\pi^2}{120}\frac{(\kb T)^4}{\hbar^3c^3},\quad C_V/V=\frac{7g\pi^2}{30}\frac{\kb^4T^3}{\hbar^3c^3}.
\end{equation}
有隙$m c^2=\Delta$的低温粒子数$\propto T^{3/2}\e^{-\Delta/\kb T}$，能领先$\Delta$乘粒子数，热容$\propto T^{-1/2}\e^{-\Delta/\kb T}$（其余常数按近底有效质量），所以指数冻结而非无质量$T^3$。

二维$\Delta U/A=3g\zeta(3)(\kb T)^3/(2\pi\hbar^2v^2)$，$C/A=9g\zeta(3)\kb^3T^2/(2\pi\hbar^2v^2)$，与习题\ref{p:333-6-5}单层公式在$g=2,b=\hbar v$一致。线性声子也$T^2$但系数$\propto c_s^{-2}$，$v/c_s\sim50$使声子通常大许多；自由薄层还可有二次弯曲支。并非只凭电子无质量就判定电子控制总热容。
\end{solution}
\originalsection{2004年：零能壳与相对论Bose气}
2004题1(a)--(d)二次零能壳谱是习题\ref{p:333-6-4}，新增(a)化学势由上一题$f(E)+f(-E)=1$给$\mu=0$，其余积分、闭式与$\sqrt T$热容逐项在原题处。2004题2(a)--(d)三阶气体临界行为完整对应习题\ref{p:334-1-1}(a)--(d)。
\begin{exercise}[2004题3(a)--(e)]\label{f:333-2004-3}
$d$维守粒子数、无自旋Bose气$\epsilon=c|p|$。(a)巨势、数密度用$\operatorname{Li}_s(z)$；(b)$P,U,U/(PV)$与经典；(c)$T_c(n)$及维数；(d)凝聚以下热容温变；(e)$T_c$热容与经典极限，区分需要的一侧极限。
\end{exercise}
\begin{solution}
态密度$\rho(E)/V=S_{d-1}E^{d-1}/[(2\pi\hbar c)^d]$，定义$C_d=S_{d-1}\Gamma(d)/(2\pi\hbar c)^d$。能量积分或巨势分部积分给
\begin{equation}
\mathcal G/V=-C_d(\kb T)^{d+1}\operatorname{Li}_{d+1}(z),\quad n=C_d(\kb T)^d\operatorname{Li}_d(z),\quad U=dPV.
\end{equation}
经典每粒子$U=d\kb T$，与量子$U/(PV)=d$同。$d>1$才有$\zeta(d)$有限，$\kb T_c=[n/(C_d\zeta(d))]^{1/d}$；这是守粒子数气，不是可任意创灭而$\mu=0$的平衡光子。

凝聚以下与正常以上固定粒子数的热容分别为
\begin{align}
\frac{C_<}{N\kb}&=d(d+1)\frac{\zeta(d+1)}{\zeta(d)}\left(\frac{T}{T_c}\right)^d,\\
\frac{C_>}{N\kb}&=d(d+1)\frac{\operatorname{Li}_{d+1}}{\operatorname{Li}_d}
-d^2\frac{\operatorname{Li}_d}{\operatorname{Li}_{d-1}}.
\end{align}
临界\textbf{左}值相对经典$dN\kb$为$(d+1)\zeta(d+1)/\zeta(d)$；$1<d\le2$左右同值，$d>2$右侧须减$d\zeta(d)/\zeta(d-1)$，有有限热容跃变。例如$d=3$两侧约$10.805N\kb$、$4.228N\kb$，经典$3N\kb$。原题若只写“在$T_c$的热容”，须明确取凝聚侧，不能把左右混为一个数。
\end{solution}
\originalsection{2005、2006年的其余编号}
2005题2(a)--(f)相对论Coulomb气完整对应习题\ref{r:333-3-9}，同一齐次Hamilton量的缩放不由密度本身破坏。2006题3(a)--(e)双层石墨烯完整对应习题\ref{p:333-6-5}(a)--(e)，包括粒子空穴积分与两层电子热容交叉尺度。
\begin{exercise}[2005题3(a)--(f)]\label{f:333-2005-3}
理想Fermi气巨系综、单态$k$能$E_k$。(a)联合占据概率；(b)用平均占据重写；(c)$l$种概率$p_j$熵及最大值；(d)Fermi占据熵、零温；(e)总$N$方差、零温；(f)由$\operatorname{Var}N=Nn\kb T\kappa_T$估金属零温压缩率，单位Å$^3$/eV。
\end{exercise}
\begin{solution}
每态Bernoulli，$p_k=f(E_k)$，$P(\{n_k\})=\prod_kp_k^{n_k}(1-p_k)^{1-n_k}$、$n_k=0,1$。任意离散熵$S=-\kb\sum p_j\log p_j\le\kb\log l$，由凹性或相对均匀分布的相对熵，均匀时取等号。因此$S_F=-\kb\sum_k[p_k\log p_k+(1-p_k)\log(1-p_k)]$。

零温非简并Fermi海每态0或1，熵与数方差0，$\operatorname{Var}N=\sum_kp_k(1-p_k)$。但低温方差$\simeq\kb T\rho(E_F)$，除$T$有有限响应：$\kappa_T=\rho(E_F)/(Vn^2)=3/(2nE_F)$（三维二次色散）。取示例$n=0.1$ Å$^{-3}$、裸电子自由模型$E_F\simeq7.86$ eV，得$1.91$ Å$^3$/eV。准确金属参数各不相同，原题未给$n$故估值必须列假设，不能把$T=0$直接代涨落式得到错误$0/0$。
\end{solution}
\begin{exercise}[2006题2(a)--(f)]\label{f:333-2006-2}
链$N$个长$a$无质量连杆在细管中只能平行/反平行，端距$L$，$x=L/(Na)$。(a)态数；(b)Stirling熵；压挤给原题近似能$E=-\sigma a^2Nx^2/2$，用于每个给定$L$的态。(c)自由能及端力$J$，功$J\dd L$；(d)高低温等温线与$T_c$；(e)稳定与不稳定范围；(f)无外力近$T_c$下长度。
\end{exercise}
\begin{solution}
$N_\pm=N(1\pm x)/2$，$\Omega=\binom N{N_+}$（奇偶允许），$S/(N\kb)=\log2-[(1+x)\log(1+x)+(1-x)\log(1-x)]/2$。$F=E-TS$，力$J=F_L=(\kb T/a)\operatorname{atanh}x-\sigma ax$。$\kb T_c=\sigma a^2$，小$x$力斜率过零；下图纵轴$J/(\sigma a)$。
\begin{center}\begin{tikzpicture}\begin{axis}[width=10cm,height=4.6cm,xlabel={$x$},ylabel={$J/(\sigma a)$},xmin=-1,xmax=1,ymin=-1.7,ymax=1.7,legend pos=north west]\addplot[structurecolor,domain=-.96:.96,samples=180]{1.3*.5*ln((1+x)/(1-x))-x};\addlegendentry{$T/T_c=1.3$}\addplot[gray,domain=-.96:.96,samples=180]{.5*ln((1+x)/(1-x))-x};\addlegendentry{$T/T_c=1$}\addplot[main,domain=-.96:.96,samples=180]{.6*.5*ln((1+x)/(1-x))-x};\addlegendentry{$T/T_c=.6$}\end{axis}\end{tikzpicture}\end{center}
局部稳定$F_{LL}=[\kb T/(1-x^2)-\sigma a^2]/(Na^2)>0$，低温$|x|<\sqrt{1-T/T_c}$内不稳定。全局平衡要比较两个极小/凸包，不把所有稳定支都当全局最优。无力$T/T_c=x/\operatorname{atanh}x$，展开$\operatorname{atanh}x=x+x^3/3+\cdots$，$|L|\simeq Na\sqrt{3(1-T/T_c)}$。高温$J=0$只有$x=0$，低温两种对称伸长方向形成题给平均场分岔。
\end{solution}

\volpart{第四编　8.334公开作业与复习题}
\originalchapter{8.334作业1：相变、薄膜与耦合序参量}
\psource{Mehran Kardar，8.334 Spring 2014，PS1，截止2014年2月21日。}
\begin{exercise}[原题1(a)--(g)]\label{p:334-1-1}
$P=\kb Tn-bn^2/2+cn^3/6$，$b,c>0$不随温度变。(a)临界$T_c,n_c,P_c$与$\kb T_cn_c/P_c$；(b)压缩率及$n_c$上的温度行为；(c)临界等温线；(d)近临界$n_\pm\simeq n_c(1\pm\delta)$，以Maxwell求$\delta$隐式式和指数。另外Dieterici气体$P(v-b)=\kb T\exp[-a/(\kb Tv)]$，$v=V/N$：(e)临界$Pv/(\kb T)$；(f)$v_c$上的压缩率；(g)临界等温线最低非零项。
\end{exercise}
\begin{solution}
前四问与8.333 PS5题2相同，既有热统卷已有完整题解。这里保留用于后续标度的必要推导：$P_n=P_{nn}=0$给$n_c=b/c$、$\kb T_c=b^2/(2c)$、$P_c=b^3/(6c^2)$，比3。$\kappa_T=1/[nP_n]$，在$n_c$为$[n_c\kb(T-T_c)]^{-1}$，由上方正发散、下方负表示均匀模型不稳定；临界$P-P_c=c(n-n_c)^3/6$。共存同时满足$P(n_+)=P(n_-)$和
\begin{equation}
\kb T\log(n_+/n_-)-b(n_+-n_-)+c(n_+^2-n_-^2)/4=0.
\end{equation}
近临界的对称\textbf{领先}代入得$T/T_c\simeq2\delta/\log[(1+\delta)/(1-\delta)]$、$\delta^2\simeq3(1-T/T_c)$。高阶须允许中心密度位移，不可把对称近似当精确全温binodal。

Dieterici式的对数导数$P_v/P=-1/(v-b)+a/(\kb Tv^2)$，再次微分并令临界一二阶为0，得$v_c=2b$、$\kb T_c=a/(4b)$、$P_c=\kb T_c/(b\e^2)$，压缩因子$2/\e^2$。在$v_c$有$\kappa_T=[2P(v_c,T)(1-T_c/T)]^{-1}$，上方仍$\propto(T-T_c)^{-1}$。临界$P'''(v_c)=-P_c/(2b^3)$，所以$P-P_c=-P_c(v-v_c)^3/(12b^3)+O((v-v_c)^4)$。不同非普适系数与压缩因子，保留同一均场指数。
\end{solution}
\begin{exercise}[原题2(a)--(i)]\label{p:334-1-2}
简单立方Heisenberg薄膜有$n$层，层内平方格4近邻，单位三分量自旋；$-\beta H=J\sum_{\alpha,\langle ij\rangle}s_i^\alpha\cdot s_j^\alpha+J_V\sum_{\alpha=1}^{n-1}\sum_i s_i^\alpha\cdot s_i^{\alpha+1}$。试探$-\beta H_0=\sum_{\alpha i}h^\alpha\cdot s_i^\alpha$。(a)求$Z_0,F_0$；(b)用Langevin函数求$m^\alpha$；(c)场平行时$\avg{\beta H}_0$；(d)Gibbs变分方程及数值方法；(e)随$J=\widehat J/(\kb T)$等求无穷层$T_c$与磁化；(f)大$n$的$T_c(n)$；(g)连续三阶方程的$tanh$解及含义；(h)半无穷系统与愈合长；(i)表层磁化临界指数。可用$L^{-1}(m)=3m+9m^3/5+O(m^5)$。
\end{exercise}
\begin{solution}
用球面测度$\dd\Omega$，单站$z(h)=4\pi\sinh h/h$，层面积站点数$N_\parallel$，$Z_0=\prod_\alpha z(h_\alpha)^{N_\parallel}$，$\beta F_0=-N_\parallel\sum\log z$。$m_\alpha=L(h_\alpha)=\coth h_\alpha-1/h_\alpha$，方向随$h_\alpha$。独立站点给$\avg{\beta H}_0/N_\parallel=-2J\sum m_\alpha^2-J_V\sum m_\alpha m_{\alpha+1}$，而$\avg{\beta H_0}_0=-N_\parallel\sum h_\alpha m_\alpha$。

将以上代$F_{\rm var}=F_0+\avg{H-H_0}_0$，以$m$作变量，单站熵Legendre项导数$h$，得
\begin{equation}
L^{-1}(m_\alpha)=4Jm_\alpha+J_V(m_{\alpha-1}+m_{\alpha+1}),\qquad m_0=m_{n+1}=0.
\end{equation}
数值可用有阻尼固定点迭代或Newton解三对角Jacobian，从小非零磁化起步并比較变分自由能；在临界附近需要检查最小Hessian，纯迭代收敛不保证全局最优。体相$3\kb T_c=4\widehat J+2\widehat J_V$。设$\epsilon=4J+2J_V-3>0$，小磁化$m_b^2=5\epsilon/9$，均场$\beta=1/2$。

有限层线性最大本征值$4J+2J_V\cos[\pi/(n+1)]$，故$\kb T_c(n)=[4\widehat J+2\widehat J_V\cos(\pi/(n+1))]/3$，大$n$比体相低$\widehat J_V\pi^2/[3(n+1)^2]$。这是均场预测；真正有限层连续对称、短程二维模型在非零温无此长程序。

以层距作单位，差分展开$J_Vm''=-\epsilon m+9m^3/5$。积分一次并取体边界，$m(x)=m_b\tanh[k(x-x_0)]$，$k=\sqrt{\epsilon/(2J_V)}$；两边$\mp m_b$描述畴壁。半空间选$\tanh[k(x+x_e)]$以满足外推到$x=-x_e$的零边界，愈合长$\lambda=1/k\propto|T-T_c|^{-1/2}$。表层$x=O(1)$有$m_s\sim m_b k\propto|T-T_c|$，普通表面转变均场$\beta_1=1$。增强表面耦合可能先表面有序，不属于此均匀耦合假设。
\end{solution}
\begin{exercise}[原题3(a)--(i)]\label{p:334-1-3}
He4复序参量$\psi$与He3浓度场$c$有
\begin{equation}
\beta H=\int\dd^dx\,[K|\nabla\psi|^2/2+t|\psi|^2/2+u|\psi|^4+v|\psi|^6+c^2/(2\sigma^2)-\gamma c|\psi|^2],
\end{equation}
$K,u,v>0$。(a)积分$c$；(b)鞍点相图及连续/不连续的$\sigma_*$；(c)序参量跃变如何消失；(d)浓度跃变；(e)$t,\sigma$相边界图；(f)证明$\avg{c-\gamma\sigma^2|\psi|^2}=0$；(g)$x\ne y$的$\avg{c(x)c(y)}$关系；(h),(i)无序、有序区的长距离衰减。
\end{exercise}
\begin{solution}
配方$c^2/(2\sigma^2)-\gamma cR^2=(c-\gamma\sigma^2R^2)^2/(2\sigma^2)-\gamma^2\sigma^2R^4/2$，积分给与$\psi$无关常数，$u_e=u-\gamma^2\sigma^2/2$。均匀密度势$f(x)=tx/2+u_ex^2+vx^3$，$x=|\psi|^2\ge0$。

$u_e>0$时$t=0$连续，$t<0$满足$t/2+2u_ex+3vx^2=0$。$u_e<0$共存由$f(x)=f(0),f'(x)=0$解$x_*=-u_e/(2v)$、$t_*=u_e^2/(2v)>0$。三临界$\sigma_*=\sqrt{2u}/|\gamma|$、$t=0$。$|\psi|$跃变$\sqrt{-u_e/(2v)}\propto(\sigma-\sigma_*)^{1/2}$，平均浓度为$\gamma\sigma^2x$，跃变$\gamma\sigma^2x_*$为一阶消失。两边界在$(0,\sigma_*)$相切：连续侧$t=0$，一阶侧$t=u_e(\sigma)^2/(2v)$，有序在其下。
\begin{center}\begin{tikzpicture}[x=3cm,y=2cm]\draw[->](0,0)--(2,0)node[right]{$\sigma/\sigma_*$};\draw[->](0,-.4)--(0,1.4)node[above]{$t$};\draw[thick,structurecolor](0,0)--(1,0);\draw[thick,structurecolor,domain=1:1.7,samples=60]plot(\x,{(\x*\x-1)^2/3});\fill(1,0)circle(1.3pt);\node[below]at(1,0){1};\node at(.55,.65){无序};\node at(1.2,-.25){有序};\end{tikzpicture}\end{center}
条件于$\psi$，$c=\gamma\sigma^2|\psi|^2+\eta$、$\avg\eta=0$，$\avg{\eta(x)\eta(y)}=\sigma^2\delta(x-y)$，直接给(f),(g)，后一式在重合点必须加接触项。无序Gaussian复场$\avg{|\psi(x)|^2|\psi(0)|^2}_c\propto|G(x)|^2$，远距$\e^{-2r/\xi}$乘幂次；有序区非连接相关趋$\avg c^2$，Gaussian振幅模给指数趋近。超出Gaussian阶，无隙相位也可通过复合算符产生幂次尾；不把振幅的有限相关长误称整个有序相的所有相关均指数。
\end{solution}
\begin{exercise}[原题4(a)--(c)]\label{p:334-1-4}
二维薄片用三维位置向量$r(x_1,x_2)$描述，能量对嵌入空间平移、旋转不变。(a)说明最低二次梯度项为$t\sum_\alpha|\partial_\alpha r|^2/2$与$K\sum_\alpha|\partial_\alpha^2r|^2/2$；(b)列两种最低四次项；(c)$K>0$、四次稳定时$t$变号的状态。
\end{exercise}
\begin{solution}
平移排除$r^2$；嵌入旋转只允许三维指标内积，最低导数是切矢$\partial_\alpha r$，其次弯曲$\partial_\alpha^2r$。原题平方格形式保留$q_1^4+q_2^4$；若另要求片内坐标连续旋转，可用$|\nabla^2r|^2$，混合导数积分后同类，但不能把这额外对称默认为原题已要求。

定义诱导度量$g_{\alpha\beta}=\partial_\alpha r\cdot\partial_\beta r$，各向同性最低四次不变量$u(\Tr g)^2/4+v\Tr(g^2)/4$。它们对应切矢长度与彼此夹角，非仅单一$|r|^4$。$t>0$鞍点切矢0、薄片没有平面取向，称皱缩。$t<0$若$v>0$且$2u+v>0$，最小$g=\zeta^2I$、$\zeta^2=-t/(2u+v)$，$r=\zeta(x_1e_1+x_2e_2)$为平直状态，弯曲项稳定慢起伏。若四次项偏好秩1度量，可出现管状各向异性；仅“四次稳定”不独自保证平面，须列出选择平面的系数条件。
\end{solution}

\originalchapter{8.334作业2：质量生成、随机场与长程相互作用}
\psource{Mehran Kardar，8.334 Spring 2014，PS2，截止2014年3月4日。}
\begin{exercise}[原题1(a)--(j)]\label{p:334-2-1}
$n$分量$m$与标量$A$耦合，$\beta H=\int[K(\nabla m)^2/2+tm^2/2+u(m^2)^2+em^2A^2+L(\nabla A)^2/2]\dd^dx$，$K,L,u>0$，为稳定取$e\ge0$。(a)均匀鞍点$m\hat e_\ell,A=0$；(b)$C=\partial_t^2\log Z$；(c)分纵横与$a$模二阶展开；(d)纵横相关长；(e)$a$相关长；(f)$t>0$时$\avg{a(x)a(0)}$；(g)Gaussian自由能修正；(h)热容修正；(i)积分随$\xi,\Lambda,d$行为；(j)鞍点上临界维是否由标量耦合改变。
\end{exercise}
\begin{solution}
$m_0=0$（$t>0$）或$m_0^2=-t/(4u)$（$t<0$），鞍点$f_0=0$或$-t^2/(16u)$。所以$C_0/V=0$或$1/(8u)$，为阶跃。置$m=(m_0+\phi_\ell)e_\ell+\sum\phi_te_t$、$A=a$，Hessian质量$r_\ell=t+12um_0^2$、$r_t=t+4um_0^2$、$r_a=2em_0^2$；没有二次混合项，因为原耦合最少$a^2$。

纵横相关长$\sqrt{K/r_{\ell,t}}$：正侧皆$\sqrt{K/t}$，负侧纵$\sqrt{K/(-2t)}$、横无穷（Goldstone）；$\xi_a$正侧无穷，负侧$\sqrt{2uL/(-et)}$。正侧$G_a(r)=\int_q\e^{\ii qx}/(Lq^2)$，$d>2$远距$\Gamma(d/2-1)/(4\pi^{d/2}Lr^{d-2})$，$d\le2$无限体积绝对场相关需IR规整。

Gaussian积分得
\begin{equation}
\delta\log Z/V=-\tfrac12\int_q[\log(Kq^2+r_\ell)+(n-1)\log(Kq^2+r_t)+\log(Lq^2+r_a)],\quad \int_q=\int_{|q|<\Lambda}\frac{\dd^dq}{(2\pi)^d}.
\end{equation}
省去与$t$无关测度常数。各分支质量随$t$线性，$\delta C/V=\tfrac12\int_q\sum_i(r_i')^2/(K_iq^2+r_i)^2$。正侧为$n\int_q(Kq^2+t)^{-2}/2$；负侧为$[4\int_q(Kq^2-2t)^{-2}+(e/2u)^2\int_q(Lq^2-et/(2u))^{-2}]/2$。横模质量恒0，不产生这阶$t$导数；仅$t=0$不可用分段微分跨越鞍点。

$\int_q(q^2+\xi^{-2})^{-2}$在$d<4$奇异$\propto\xi^{4-d}$，$d=4$为$\log(\Lambda\xi)$，$d>4$UV主导$\Lambda^{d-4}$加较小奇异项，故Gaussian判据仍$d_u=4$。$A$增加奇异项系数，不改变这一维数；不过积分掉标量可生成非解析有效势，某些参数下使连续转变被一阶转变取代，该一阶问题不由二阶鞍点计算自动排除。
\end{solution}
\begin{exercise}[原题2(a)--(i)]\label{p:334-2-2}
标量$\beta H=\int[K(\nabla m)^2/2+tm^2/2+um^4-h(x)m]\dd^dx$，$u>0$，淬火Gaussian场$\overline h=0$、$\overline{h(x)h(y)}=\Delta\delta(x-y)$，等价$\overline{h(q)h(q')}=\Delta(2\pi)^d\delta(q+q')$。(a)均匀鞍点淬火平均；(b)二阶涨落；(c)Fourier能；(d)固定$h$下均值与连接方差；(e)平均散射线形；(f),(g)固定$h$及平均Gaussian自由能；(h)奇异$t$幂次；(i)上临界维。
\end{exercise}
\begin{solution}
均匀假设仅感受$\bar h_V=V^{-1}\int h$，其方差$\Delta/V$，热力学平均消失，故$m_0=0$或$m_0^2=-t/(4u)$、$f_{sp}=0$或$-t^2/(16u)$。但真实随机场局部变形不能由均匀平均排除。

置$m=m_0+\phi$，二次核$D(q)=Kq^2+r$、$r=t+12um_0^2$，$\beta H_2=\tfrac12\int_qD(q)|\phi(q)|^2-\int_q h(-q)\phi(q)$。配方$\phi=h/D+\delta\phi$，得到$\avg{\phi(q)}_h=h(q)/D$、$\avg{\delta\phi(q)\delta\phi(q')}_h=(2\pi)^d\delta(q+q')/D$。淬火后总线形
\begin{equation}
S(q)=D(q)^{-1}+\Delta D(q)^{-2},\quad
\delta f[h]=\tfrac12\int_q\log D-\frac1{2V}\int_q|h(q)|^2/D.
\end{equation}
第二项不是热连接方差；淬火平均$\overline{\delta f}=\tfrac12\int_q\log D-\tfrac\Delta2\int_qD^{-1}$。减去解析UV项后热奇异$r^{d/2}$，随机项$-\Delta r^{(d-2)/2}$，偶数边界出现对数。负側平均场$f_{sp}\sim t^2$；随机修正相对它$\sim\Delta |t|^{(d-6)/2}$，故上临界维6，不能由清洁系统Ginzburg的4代替。此Gaussian判据不等价于全非微扰随机场临界解，也不证明维数约化。
\end{solution}
\begin{exercise}[原题3(a)--(c)]\label{p:334-2-3}
连续自旋除短程项外有铁磁尾$J(r)\propto r^{-d-\sigma}$。(a)二次核的变化及何时主导；(b)由Goldstone涨落求下临界维；(c)鞍点上临界维。
\end{exercise}
\begin{solution}
差能量用$\int\dd^drJ(r)[1-\cos(qr)]$，$0<\sigma<2$作$r=y/q$得到$|q|^\sigma$，系数正；$\sigma>2$二阶矩有限，领先$q^2$，等号有$q^2\log q$。因此Gaussian层长程主导$\sigma<2$。Goldstone方差$\int_0^\Lambda q^{d-1}\dd q/q^\sigma$在$d\le\sigma$IR发散，所以连续对称$d_\ell=\sigma$（若短程主导则2）。场维$(d-\sigma)/2$使$um^4$尺度维$2\sigma-d$，$d_u=2\sigma$；相互作用临界下长短程精确交界可移到$\sigma=2-\eta_{SR}$，不能把Gaussian阈值当所有参数的全阶结论。
\end{solution}

\originalchapter{8.334作业3：标度、耦合场与长程RG}
\psource{Mehran Kardar，8.334 Spring 2014，PS3，截止2014年3月21日。}
\begin{exercise}[原题1(a)--(c)]\label{p:334-3-1}
二维$n$分量非线性$\sigma$模型给$\dot T=(n-2)T^2/(2\pi)$、$\dot h=2h$，点为对$\ell$求导。(a)$T\to0$相关长；(b)$T,h\to0$奇异自由能；(c)$h=0$磁化率。
\end{exercise}
\begin{solution}
先假设$n>2$。积分$1/T(\ell)=1/T_0-(n-2)\ell/(2\pi)$，到$T(\ell_*)=O(1)$为止，$\ell_*=2\pi/[(n-2)T_0]+O(1)$，$\xi\sim a\e^{\ell_*}$。面积自由能$f_s=\e^{-2\ell_*}\Phi(h\e^{2\ell_*})$，即$f_s\sim\xi^{-2}\Phi(h\xi^2)$（指标按题给截断流，忽略高阶场异常维）。二次场导数$\chi\sim\xi^2\sim\exp[4\pi/((n-2)T)]$。

$n=2$的该微扰流恒0，不能除以$n-2$；XY需要涡旋的BKT方程。$n<2$亦不适用此渐近自由指数尺度，原公式没有无条件覆盖所有$n$。
\end{solution}
\begin{exercise}[原题2(a)--(e)]\label{p:334-3-2}
两标量$\beta H=\int[tm^2/2+K(\nabla m)^2/2-hm+L(\nabla^2\phi)^2/2+v\nabla m\cdot\nabla\phi]\dd^dx$。(a)Fourier式；(b)$q'=bq,m'(q')=m(q)/z,\phi'(q')=\phi(q)/y$的递推，忽略常数积分；(c)$K,L$固定点$y_t,y_h,y_v$；(d)$f_s=t^{2-\alpha}g(h/t^\Delta,v/t^\omega)$的指数；(e)$t,L$固定的另一点，哪些扰动相关。
\end{exercise}
\begin{solution}
$\beta H=\frac12\int_q(m(-q),\phi(-q))\left(\begin{smallmatrix}t+Kq^2&vq^2\\vq^2&Lq^4\end{smallmatrix}\right)\binom{m(q)}{\phi(q)}-hm(0)$。消去高动量Gaussian仅贡献自由能常数，低模改变量得$t'=b^{-d}z^2t$、$K'=b^{-d-2}z^2K$、$L'=b^{-d-4}y^2L$、$v'=b^{-d-2}zyv$、$h'=zh$。

取$z=b^{(d+2)/2}$、$y=b^{(d+4)/2}$，$K,L$固定，$y_t=2,y_h=(d+2)/2,y_v=1$。自由能体积缩放$b^{-d}$，选$b=t^{-1/2}$得$\alpha=2-d/2$、$\Delta=(d+2)/4$、$\omega=1/2$。$v$相关且把临界线移到$t=v^2/L$：配方有效$m$核$t-v^2/L+Kq^2$（$q\ne0$），稳定区$t>v^2/L$。不能在$t=0$保留有限$v$后继续使用不稳定Gaussian。

若固定$t,L$，则$z=b^{d/2}$、$y=b^{(d+4)/2}$，$K'=b^{-2}K$无关、$v'=v$边缘、$h'=b^{d/2}h$相关；这是有质量$m$配以弯曲无质量$\phi$的固定描述，$v$只在$t>v^2/L$的稳定区改变有效系数。两套场重标度反映选择不同主导核，不能混用其$y_t$。
\end{solution}
\begin{exercise}[原题3(a)--(j)]\label{p:334-3-3}
$d$维分为$n$个平行、$d-n$个垂直方向，$\beta H_0=\int[K(\nabla_\parallel m)^2/2+L(\nabla_\perp^2m)^2/2+tm^2/2-hm]\dd^nx_\parallel\dd^{d-n}x_\perp$，扰动$U=u\int m^4$。(a)Fourier；(b)$q'_\parallel=bq_\parallel,q'_\perp=cq_\perp,m'=m/z$递推；(c)$K,L$固定的$c,z,y_t,y_h$；(d)自由能、$\alpha,\Delta$；(e)零场$\chi$；(f)Gaussian两点和$\chi_0(q)$；(g)$U$的Fourier表达；(h)上临界维；(i)一阶展开用Wick化简；(j)$\chi(q)$与一阶$t_c$，无需积分。
\end{exercise}
\begin{solution}
记$D(q)=t+Kq_\parallel^2+Lq_\perp^4$，$H_0=\tfrac12\int_qD|m|^2-hm(0)$。设体积因子$B=b^nc^{d-n}$，有$K'=z^2K/(Bb^2)$、$L'=z^2L/(Bc^4)$、$t'=z^2t/B$、$h'=zh$。两梯度同时固定要求$c=b^{1/2}$，定义有效缩放维$D_e=n+(d-n)/2=(d+n)/2$，$z=b^{(D_e+2)/2}$。因此$y_t=2,y_h=(D_e+2)/2$，$f_s=b^{-D_e}f_s(b^2t,b^{y_h}h)$，$2-\alpha=D_e/2$、$\Delta=(D_e+2)/4$、$\chi=t^{-1}$。

$\avg{m(q)m(q')}_0=(2\pi)^d\delta(q+q')/D(q)$，$\chi_0=1/D$。$U=u\int_{q_1q_2q_3q_4}(2\pi)^d\delta(\sum q_i)\prod m(q_i)$。其递推$u'=uz^4/B^3=b^{4-D_e}u$，故$d_u=8-n$（在题给固定$n$的方向划分下）。

连接一阶展开$\avg{m_qm_{q'}}=\avg{m_qm_{q'}}_0-\avg{m_qm_{q'}U}_{0,c}$。顶点四条腿选两条接外腿有$4\times3=12$，另两条自收缩$I_1=\int_kD(k)^{-1}$，故$\chi(q)=D(q)^{-1}-12uI_1D(q)^{-2}+O(u^2)$。重写为$[t+12uI_1+Kq_\parallel^2+Lq_\perp^4]^{-1}$到一阶，$t_c=-12u\int_k[Kk_\parallel^2+Lk_\perp^4]^{-1}+O(u^2)$，需同一UV截断与IR可积条件；低维发散正是Gaussian失效信号。平行$\xi\sim t^{-1/2}$、垂直$\xi\sim t^{-1/4}$，同一个长度指数不能描述两方向。
\end{solution}
\begin{exercise}[原题4(a)--(g)]\label{p:334-3-4}
$N_s$分量自旋有$J(r)\propto r^{-d-\sigma}$长程项和短程LG项。(a)导出$t+K_2q^2+K_\sigma q^\sigma$及$u(m^2)^2$核；(b)$u=0,\sigma>2$递推与固定核；(c)$\sigma<2$的场重标度与$K_2$无关；(d)Gaussian$\nu,\eta,\gamma$、$d>2\sigma$时$u$；(e)微扰RG；(f)$d<2\sigma$的首阶指数，原题记$\epsilon=d-2\sigma$；(g)$J(r)\propto\e^{-r/a}$时行为。
\end{exercise}
\begin{solution}
长程非解析项的推导见习题\ref{p:334-2-3}；Fourier四次项$u\int_{q_1q_2q_3}(m_{q_1}\cdot m_{q_2})(m_{q_3}\cdot m_{-q_1-q_2-q_3})$。固定$q^2$核时$t'=b^2t$、$K_\sigma'=b^{2-\sigma}K_\sigma$，$\sigma>2$流向短程Gaussian。固定$q^\sigma$则Fourier场因子$z=b^{(d+\sigma)/2}$，$t'=b^\sigma t$、$K_2'=b^{\sigma-2}K_2$。于是$\nu=1/\sigma$，$\chi(q)\sim q^{-\sigma}=q^{-2+\eta}$给$\eta=2-\sigma$，$\gamma=1$，$u$维$2\sigma-d$。

令$K_\sigma=1$、$A_d=S_{d-1}/(2\pi)^d$，窄壳积分$I_1=A_d\Lambda^{d-\sigma}\dd\ell$、$I_2=A_d\Lambda^{d-2\sigma}\dd\ell$。为核对原题(e)的系数，将$m=\varphi+\chi$，$u(m^2)^2$中的两快腿部分写作$V_2=\int\chi_aB_{ab}\chi_b$，$B_{ab}=2u\varphi^2\delta_{ab}+4u\varphi_a\varphi_b$。于是$\operatorname{tr}B=2u(N_s+2)\varphi^2$，一次收缩给$\delta t=4(N_s+2)uI_1$；两顶角连接收缩给$-I_2\operatorname{tr}B^2=-4(N_s+8)u^2I_2(\varphi^2)^2$。展开$I_1(t)=A_d\Lambda^{d-\sigma}\dd\ell[1-t/\Lambda^\sigma+\cdots]$并缩放，对$u(m^2)^2$约定（没有$4!$）得到
\begin{align}
\dot t&=\sigma t+4(N_s+2)u A_d\Lambda^{d-\sigma}-4(N_s+2)ut A_d\Lambda^{d-2\sigma}+\cdots,\\
\dot u&=(2\sigma-d)u-4(N_s+8)u^2 A_d\Lambda^{d-2\sigma}+\cdots.
\end{align}
动量非解析项不由局部壳积分重新生成，其系数这一阶无修正。取$\bar\epsilon=2\sigma-d=-\epsilon>0$，平移质量至临界流形后，$u_*A_d\Lambda^{-\bar\epsilon}=\bar\epsilon/[4(N_s+8)]$，
\begin{equation}
\nu^{-1}=\sigma-\frac{N_s+2}{N_s+8}\bar\epsilon+O(\bar\epsilon^2),\qquad \eta=2-\sigma.
\end{equation}
$\gamma=\sigma\nu$在该长程固定点。注意原题$\epsilon$是负数，不可直接套通常正的$4-d$公式。有限程指数尾所有低阶矩有限，Fourier在$q=0$解析，流向短程普适类（非Gaussian时是短程Wilson--Fisher）。长程固定点与短程点的交界需比较$\sigma$与$2-\eta_{SR}$，首阶Gaussian分界只是近似。
\end{solution}

\originalchapter{8.334作业4：转移矩阵与实空间重整化}
\psource{Mehran Kardar，8.334 Spring 2014，PS4，截止2014年4月9日。原卷前四页介绍转移矩阵，理论第\ref{ch:duality}章及本章承接，全部七题包括选做题收入。}
\begin{exercise}[原题1(a)--(d)]\label{p:334-4-1}
一维spin1取$s=-1,0,1$，$-\beta H=K\sum s_is_{i+1}$。(a)转移矩阵；(b)利用对称求最大本征值与$\log Z/N$；(c)相关长和$K\to\infty$；(d)消去隔站后闭合的最简参数空间。
\end{exercise}
\begin{solution}
按$(-1,0,1)$，$T=\left(\begin{smallmatrix}\e^K&1&\e^{-K}\\1&1&1\\\e^{-K}&1&\e^K\end{smallmatrix}\right)$，奇向量$(1,0,-1)/\sqrt2$的$\lambda_o=2\sinh K$。偶空间取$(1,0,1)/\sqrt2,(0,1,0)$，块为$\left(\begin{smallmatrix}2\cosh K&\sqrt2\\\sqrt2&1\end{smallmatrix}\right)$，最大
\begin{equation}
\lambda_+=\frac{2\cosh K+1+\sqrt{(2\cosh K-1)^2+8}}2,\qquad \log Z/N\to\log\lambda_+.
\end{equation}
自旋关联连接偶基态与奇模，$\xi^{-1}=\log(\lambda_+/|2\sinh K|)$；$K\to+\infty$时$\lambda_+-\lambda_o\simeq4\e^{-K}$，$\xi\simeq\e^{2K}/4$。有限温正矩阵最大根单一，无非零温相变。

闭合核可写$g+Jss'+Qs^2s'^2-D(s^2+s'^2)/2$。四个独立对称权重$T_{00},T_{01},T_{11},T_{1,-1}$都可拟合这四参数，平方$T^2$保持全局$s\mapsto-s$对称，却生成晶场$D$与双二次$Q$，仅$J$不能闭合。
\end{solution}
\begin{exercise}[原题2(a)--(d)]\label{p:334-4-2}
$q$态clock自旋$s\in\mathbb Z_q$，近邻$-\beta H=\sum J(s_i-s_j)$，$J$为任意周期实函数。(a)环链$Z=\Tr T^N$；(b)显写循环矩阵并离散Fourier对角化；(c)热力学自由能；(d)以投影$\Pi_\alpha$写$\avg{\delta_{s_i,s_{i+l}}}$，原题声称连接部分$\sim[\lambda(1)/\lambda(0)]^l$，说明条件。
\end{exercise}
\begin{solution}
$w_n=\e^{J(n)}$、$T_{ab}=w_{a-b}$，环乘积求和即迹。Fourier本征向量$q^{-1/2}\e^{-2\pi\ii ka/q}$，$\lambda_k=\sum_{n=0}^{q-1}w_n\e^{2\pi\ii kn/q}$，$
Z=\sum_k\lambda_k^N$。严格正权重给Perron根$\lambda_0=\sum w_n>|\lambda_{k\ne0}|$，$\beta f=-\log\lambda_0$。

$\Pi_\alpha$是对角仅第$\alpha$为1，$\avg{\delta_{s_i,s_{i+l}}}=Z^{-1}\sum_\alpha\Tr(\Pi_\alpha T^l\Pi_\alpha T^{N-l})$。无穷链由均匀稳态得
\begin{equation}
\avg{\delta_{s_i,s_{i+l}}}_c=\frac1q\sum_{k=1}^{q-1}(\lambda_k/\lambda_0)^l.
\end{equation}
一般$J(n)$下主导为非零$k$中的最大模，未必$k=1$；例如$q=4$、$w_0=w_2>w_1=w_3$，$\lambda_2$主导而$\lambda_1=0$。原题末式对通常铁磁cosine clock中$k=\pm1$最慢模成立，但不能要求证明任意周期$J$下的错误断言。
\end{solution}
\begin{exercise}[原题3(a)--(c)]\label{p:334-4-3}
一维XY单位自旋$-\beta H=K\sum\cos(\theta_i-\theta_{i+1})$。(a)转移核与$\e^{\ii m\theta}$本征函数；(b)自由能及$T\propto K^{-1}\to0$热容；(c)相关長与$K\to\infty$。
\end{exercise}
\begin{solution}
用$\dd\theta$测度，$T(\theta,\theta')=\e^{K\cos(\theta-\theta')}$，卷积本征值$\lambda_m=2\pi I_m(K)$，$f=-\kb T\log[2\pi I_0(K)]$。$K=\widehat J/(\kb T)$，能每站$-\widehat J I_1/I_0$，$I_0(K)\simeq\e^K/\sqrt{2\pi K}$给$U/N=-\widehat J+\kb T/2+\cdots$，所以经典自旋$C/N\to\kb/2$，这是独立相对角谐振模；若量子自由度会有低温量子修正。

$s_i\cdot s_{i+l}$关联取$m=\pm1$，$\xi^{-1}=\log(I_0/I_1)$，$I_1/I_0=1-1/(2K)+O(K^{-2})$，$\xi\simeq2K$。与离散Ising的激活式$\e^{2K}$不同，连续角涨落用渐进的小角变化破坏一维有序。
\end{solution}
\begin{exercise}[原题4(a)--(f)，选做]\label{p:334-4-4}
一维经典气有足够硬核使只有最近粒子相互作用，排序$0\le x_1\le\cdots\le x_N\le L$。(a)写$H=\sum p_i^2/(2m)+\sum_{i=2}^NV(x_i-x_{i-1})$；(b)$Z$及$\delta_1=x_1,\delta_i=x_i-x_{i-1}$区域；(c)Gibbs--Laplace配分$\mathcal Z(P)=\int_0^\infty\e^{-\beta PL}Z(L)\dd L$与鞍点压力；(d)$\delta_{N+1}=L-\sum_1^N\delta_i$分离积分；(e)平均长、密度与无凝聚；(f)硬棒最小间距$a$的状态方程及virial系数。
\end{exercise}
\begin{solution}
排序去掉$N!$，动量积分后$Z=\lambda^{-N}\int_{\delta_i\ge0,\sum_1^N\delta_i\le L}\prod_{i=1}^N\dd\delta_i\exp[-\beta\sum_{i=2}^NV(\delta_i)]$。两端外侧没有粒子间势。Laplace鞍点$\partial_L\log Z=\beta P$，就是canonical压力。

最后空隙变量使$L=\sum_1^{N+1}\delta_i$，Jacobian1。两个边界空隙积分各$(\beta P)^{-1}$，$I(P,T)=\int_0^\infty\e^{-\beta[V(\delta)+P\delta]}\dd\delta$，故$\mathcal Z=\lambda^{-N}(\beta P)^{-2}I^{N-1}$。$\avg L=-\beta^{-1}\partial_P\log\mathcal Z=2/(\beta P)+(N-1)\avg\delta$，$\avg\delta=\int\delta\e^{-\beta(V+P\delta)}/I$，热力学$n=1/\avg\delta$。对稳定短程势、正压有限温，正积分解析，无均场所预言的有限温液气奇点；不是对任意不稳定势都成立。

硬棒$I=\e^{-\beta Pa}/(\beta P)$，$\avg\delta=a+\kb T/P$，故$P=n\kb T/(1-na)$，$B_l=a^{l-1}$。有限$N$的边界空隙给$\avg L=(N-1)a+(N+1)\kb T/P$，不能误当无穷体系方程。
\end{solution}
\begin{exercise}[原题5(a)--(i)]\label{p:334-4-5}
一维$q$态Potts环链$-\beta H=J\sum\delta_{s_i,s_{i+1}}$。(a)$Z,\xi$；(b)反铁磁零温是否临界；(c)隔站消去的$J,g$流；(d)固定点稳定；(e)$\log Z$的常数级数；(f)$t=\e^{-1/\xi}$化简；(g)级数回到转移矩阵答案；(h)加入$h\sum\delta_{s_i,1}$生成$K\sum\delta_{s_i,1}\delta_{s_{i+1},1}$的三参数流；(i)铁磁零温的磁本征值与相关长。
\end{exercise}
\begin{solution}
令$x=\e^J$，$T_{ab}=1+(x-1)\delta_{ab}$，$\lambda_0=x+q-1$、$\lambda_1=x-1$（$q-1$重），$Z=\lambda_0^N+(q-1)\lambda_1^N$、$\xi^{-1}=\log(\lambda_0/|\lambda_1|)$。$J\to-\infty$时$|\lambda_1/\lambda_0|=1/(q-1)$，$q>2$相关长有限，$q=2$为严格交替零温临界，不能统一回答所有$q$。

$T^2$非对角$D=2x+q-2$、对角$C=x^2+q-1$，所以$J'=\log(C/D)$、$g'=2g+\log D$。$J=0$稳定，铁磁$J=\infty$零温点不稳定；对$q=2$反铁磁极限消去隔站后映到铁磁点，交错序需另外跟踪。

铁磁取$t=(x-1)/(x+q-1)\ge0$，$t'=t^2$；反铁磁用有符号比$r$才有$r'=r^2$。从$g_0=0$，每原站$\log Z/N=\sum_{j\ge0}2^{-j-1}\log D(J_j)$。$D=q(1+t)/(1-t)$，因此级数$\log q+\sum2^{-j-1}\log[(1+t^{2^j})/(1-t^{2^j})]$，由$1-t^{2^{j+1}}=(1-t^{2^j})(1+t^{2^j})$逐项抵消为$\log[q/(1-t)]$，正是$\log\lambda_0$。

三参数最清楚从对称核$T_{ab}=\exp[g+J\delta_{ab}+K\delta_{a1}\delta_{b1}+h(\delta_{a1}+\delta_{b1})/2]$平方。按$11$、$1a$、$aa$、$ab$（$a,b\ne1,a\ne b$）四权重去掉$\e^{2g}$后分别
\begin{align}
A&=\e^{2(J+K+h)}+(q-1)\e^h,& B&=\e^{h/2}[\e^{J+K+h}+\e^J+q-2],\\
C&=\e^h+\e^{2J}+q-2,&D&=\e^h+2\e^J+q-3.
\end{align}
故$g'=2g+\log D$、$J'=\log(C/D)$、$h'=2\log(B/D)$、$K'=\log(A/D)-J'-h'$。$q=2$的参数存在冗余，可按Ising场直接闭合；$q\ge3$上式四类齐备。

$J\to\infty$有$J'=J-\log2+\cdots$，$\e^{-J}$翻倍，热本征值2、$y_t=1$；$\xi\simeq\e^J/q$。磁小扰动$\binom{h'}{K'}=\left(\begin{smallmatrix}2&1\\0&1\end{smallmatrix}\right)\binom hK+\cdots$，本征值2、1，对应$y_h=1$与额外边缘方向0。有限温一维没有自发破缺，零温磁场敏感不能推出非零$T_c$。
\end{solution}
\begin{exercise}[原题6(a)--(i)]\label{p:334-4-6}
蜂窝Ising近邻$J$。(a)四站星形簇、长度缩放$b=2$；(b)多数规则如何处理平票；(c)中心为平票决胜的簇自旋$+$微观表；(d)相邻簇内部与跨边能；(e)令$x=\e^J$，证明$R_{++}=x^8+2x^6+7x^4+14x^2+17+14x^{-2}+7x^{-4}+2x^{-6}$，$R_{+-}=9x^4+16x^2+13+16x^{-2}+9x^{-4}+x^{-8}$；(f)耦合流；(g)近似临界$J_c,\nu$与精确值比较；(h)零温铁磁的$y_h,y_t$；(i)反铁磁是否适用及破坏的对称。
\end{exercise}
\begin{solution}
星形取中心$s_0$与三个叶$s_1,s_2,s_3$，总和非0用多数符号，等于0用$s_0$。蜂窝适当四站覆盖后两相邻星形有两条跨簇边。$+$微观态8个：中心$+$、叶中负数0、1、2（分别1、3、3种）；中心$-$仅叶全$+$（1种）。$-$态为全体反号，表已穷尽16态。

设两个簇$a,b$，可标跨边$a_1b_1,a_2b_2$，总键和$E=\sum_{j=1}^3(a_0a_j+b_0b_j)+a_1b_1+a_2b_2$，权重$\e^{JE}$。固定簇$+,+$或$+,-$各8乘8种，按$E$计数得到题中两多项式；每式在$x=1$均64，最高键和8验证总8条微观键。随卷数值代码完整穷举验证系数。

局部两簇投影$R(S_A,S_B)=A_0\exp[J'S_AS_B]$，故$J'=\tfrac12\log(R_{++}/R_{+-})$。这一局部权重近似忽略相邻第三簇、多簇有效相互作用，不是整个格点分区函数的精确闭合变换。非零固定点数值$J_c\simeq1.04553$，$\lambda_t=J'(J_c)'\simeq1.81762$，$\nu=\log2/\log\lambda_t\simeq1.1600$；精确蜂窝$J_c=\tfrac12\log(2+\sqrt3)=0.65848$、$\nu=1$。较大非普适误差必须如实报告，不调系数“校准”到精确答案。

铁磁零温$J'\simeq2J-\log3$，以$T\propto1/J$坐标$T'\simeq T/2$，$y_t=-1$；一个饱和四站簇磁功$4hS$，$y_h=2$。若用激活缺陷$\e^{-2J}$作坐标，线性导数为0，与$1/J$坐标不同，不能伪称有限激活本征值。反铁磁Néel态中心与三个叶反号，簇多数投影与中心决胜不与交错翻转规范变换相容，不能保留原蜂窝两亚格子对称；需改交错块变量后另建方案。
\end{solution}
\begin{exercise}[原题7(a),(b)，选做]\label{p:334-4-7}
两个Ising自旋$H=-\widehat K\sigma_1\sigma_2$，$K=\beta\widehat K$。(a)均匀随机选一个自旋做Metropolis，显写四态转移矩阵并验证平衡权重为左本征向量；(b)单翻、双翻以$p,1-p$选取，重复。
\end{exercise}
\begin{solution}
顺序$(++,+-,-+,--)$，令$a=\min(1,\e^{-2K})$、$b=\min(1,\e^{2K})$。行代表初态，
\begin{equation}
P_1=\begin{pmatrix}1-a&a/2&a/2&0\\b/2&1-b&0&b/2\\b/2&0&1-b&b/2\\0&a/2&a/2&1-a\end{pmatrix},\quad
P_2=\begin{pmatrix}0&0&0&1\\0&0&1&0\\0&1&0&0\\1&0&0&0\end{pmatrix}.
\end{equation}
$\pi=(\e^K,\e^{-K},\e^{-K},\e^K)/[4\cosh K]$，$\e^Ka=\e^{-K}b$使每边详细平衡，直接得$\pi P_1=\pi$。双翻不改键能，必接受，且配对权重相等，$\pi P_2=\pi$，故$P=pP_1+(1-p)P_2$同样平稳。$p=0$分成两不连通轨道，虽然平稳式仍真，却不能从任意初态采到完整平衡；平稳、不可约、有效混合是三个不同问题。
\end{solution}

\originalchapter{8.334作业5：对偶、Potts模型与渗流}
\psource{Mehran Kardar，8.334 Spring 2014，PS5，截止2014年4月28日。选做题2与题5(c)保留。}
\begin{exercise}[原题1(a),(b)]\label{p:334-5-1}
平方格Ising，$-\beta H=K\sum_{\langle ij\rangle}\sigma_i\sigma_j$。(a)从自由能对偶导出$U=\avg{\beta H}=-\partial_{\log K}\log Z$关系；(b)求临界$U$精确值。
\end{exercise}
\begin{solution}
热力学每站$\phi(K)=\log Z/N$满足$\phi(K)=\log\sinh(2K)+\phi(K^*)$，$\sinh2K\sinh2K^*=1$；有限环面的不同边界扇区另处理，对能密度无影响。微分得到$\phi'=2\coth2K+\phi'(K^*)\partial_KK^*$。$\partial_KK^*=-1/\sinh2K$，所以
\begin{equation}
\frac{U(K)}N=-2K\coth2K-\frac{K}{K^*\sinh2K}\frac{U(K^*)}N.
\end{equation}
自对偶$\sinh2K_c=1$使$\phi'(K_c)=\sqrt2$，$U/N=-\sqrt2K_c$。若问物理能$H$，则每站$-\sqrt2\widehat J$，而非多乘$\beta$。该推导用能在连续转变两侧相合，不能直接以单支同值论证一阶转变潜热。
\end{solution}
\begin{exercise}[原题2，选做]\label{p:334-5-2}
验证平方格Ising的精确自由能在对偶变换后只差解析函数。
\end{exercise}
\begin{solution}
用第\ref{ch:ising2}章的积分形式
\begin{equation}
\phi(K)=\tfrac12\log(2\sinh2K)+\frac1{2\pi}\int_0^\pi\operatorname{arcosh}[\cosh2K\coth2K-\cos q]\dd q.
\end{equation}
令$s=\sinh2K$，积分的组合$\cosh^22K/\sinh2K=s+s^{-1}$，在$s\mapsto s^{-1}$下不变；首项之差$\tfrac12\log(s/s^{-1})=\log s$，即上一题关系。检查$K\to\infty$时$\phi\to2K$及$K\to0$时$\phi\to\log2$可排除归一化常数错误。
\end{solution}
\begin{exercise}[原题3(a)--(d)]\label{p:334-5-3}
平方格$q$态clock，$-\beta H=\sum_{\langle ij\rangle}J(s_i-s_j)$。(a)改用$2N$有向键$b_{ij}=s_i-s_j\pmod q$，说明每小方格键和0的约束及计数；(b)用$\delta_q(S)=q^{-1}\sum_{n=0}^{q-1}\e^{2\pi\ii nS/q}$求双格模型，$\lambda(k)=\sum_b\e^{J(b)+2\pi\ii kb/q}$；(c)Potts双耦合、$J_c(q)$；(d)各向异性Potts的自对偶线。
\end{exercise}
\begin{solution}
键来自梯度，所以任意合同闭路的有向和0。平面单连通图$E$键、$N$站，独立面约束$F=E-N+1$，剩$N-1$自由键，再加整体$q$次全局平移，恢复$q^N$站配置。环面还需两个非合同回路的零绕数；仅加小面约束会额外计入扭转边界扇区。

逐面引入Fourier整数$n_p$，每键出现相邻两面的$n_p-n_{p'}$，键求和就是$\lambda$。对平面固定外面$n=0$的精确式为$Z=q^{1-F}\sum_{\{n_p\}}\prod_{\rm edges}\lambda(n_p-n_{p'})$。体积领先$F\simeq N$给原题的$q^{-N}$类归一化；环面有限尺寸需显式扇区和，不能把这个bulk式当任意周期边界的精确单扇区等式。定义双键$\widetilde J(k)=\log\lambda(k)$加归一化常数；若$\lambda$不全正，实正Boltzmann模型解释需另外检验。

Potts$w_0=\e^J,w_{b\ne0}=1$，$\lambda_0=\e^J+q-1$、$\lambda_{k\ne0}=\e^J-1$，故$(\e^J-1)(\e^{J^*}-1)=q$，铁磁自对偶$J_c=\log(1+\sqrt q)$。以铁磁模型唯一相边界的课程前提定位临界，不由对偶独自证明转变连续；$q>4$可是一阶。双格$x$边跨原格$y$边，各向异性$v_x=\e^{J_x}-1,v_y=\e^{J_y}-1$的自对偶线$v_xv_y=q$。
\end{solution}
\begin{exercise}[原题4(a)--(c)]\label{p:334-5-4}
蜂窝、三角格Ising。(a)蜂窝一亚格子部分求和的星--三角耦合$K_t(K_h)$；(b)三角的双格蜂窝；(c)结合求两临界耦合。
\end{exercise}
\begin{solution}
中心自旋$s_0$求和$2\cosh[K_h(s_1+s_2+s_3)]=A\exp[K_t(s_1s_2+s_2s_3+s_3s_1)]$。三叶全同及两同一异给$A\e^{3K_t}=2\cosh3K_h$、$A\e^{-K_t}=2\cosh K_h$，于是$\e^{4K_t}=\cosh3K_h/\cosh K_h=4\cosh^2K_h-3$。平面双格关系$\e^{-2K_h^*}=\tanh K_t$。临界将星--三角与对偶映射接合，解得
\begin{equation}
K_t^c=\tfrac14\log3,\qquad K_h^c=\tfrac12\log(2+\sqrt3).
\end{equation}
这也给前章蜂窝近似RG的精确比较基准。有限图边界的解析配分因子不影响临界位置。
\end{solution}
\begin{exercise}[原题5(a)--(c)，(c)选做]\label{p:334-5-5}
蜂窝$q$态Potts，中心部分和生成三角近邻$K_t$和三自旋$L$。(a)求这些耦合，为何Ising无独立$L$；(b)三角模型双格；(c)要求星--双变换的固定值并验$L(K_c)=0$，求临界。
\end{exercise}
\begin{solution}
令$x=\e^{K_h}$，按三叶全同、两同一异、全异的星和为$A=x^3+q-1$、$B=x^2+x+q-2$、$C=3x+q-3$。对$q\ge3$写权重$C\exp[K_t(\delta_{12}+\delta_{23}+\delta_{31})+L\delta_{12}\delta_{23}]$，得$K_t=\log(B/C)$、$L=\log(AC^2/B^3)$。$q=2$不可能三叶全异，且$\delta_{12}\delta_{23}=(\delta_{12}+\delta_{23}+\delta_{31}-1)/2$，三体项可吸收到常数与两体项，没有独立新算符。

普通三角两体Potts的双格为蜂窝，$v_tv_h=q$，$v=\e^K-1$；若三角已有$L$，双格一般不只是单一两体耦合，不能全温宣称自对偶。要求星变换后$\e^{K_t}-1=q/v_h$且$L=0$，代$A,B,C$化简为
\begin{equation}
v_h^3-3qv_h-q^2=0,\qquad v_t^3+3v_t^2=q,\quad v_tv_h=q.
\end{equation}
取正铁磁根，$K_h^c=\log(1+v_h)$、$K_t^c=\log(1+v_t)$；该根确使$AC^2=B^3$，所以在相边界处生成三体项消失。$q=2$换成Ising耦合$K_{\rm Potts}=2K_{\rm Ising}$回收前题。并非所有温度都无三体项。
\end{solution}
\begin{exercise}[原题6(a)--(c)]\label{p:334-5-6}
独立键以$p$占据，双键占据$1-p$。(a)$N$串联键和其$N$并联双键的渗流/非渗流概率；(b)平方格键渗流阈值；(c)三维键对偶为小面渗流，能否通流同时保持固体完整？
\end{exercise}
\begin{solution}
串链通流需每键占据，概率$p^N$；双并联不通需每双键空，概率也$p^N$，双通$1-p^N$。二维平方格自对偶且通流相边界唯一，对偶映射$p\leftrightarrow1-p$定位$p_c=1/2$（证明唯一和sharpness是额外渗流定理，不能仅一句“自对偶”完成其全证明）。

三维堵塞路径的对偶对象是面形成的阻断曲面，而不是另一条平面相交链。孔通道与固体承重通道可绕行并存，因此几何上能通流而材料仍连通；机械刚性还要求约束与弹性，不等于键连通。二维同一截面上路径与双路径互斥的简单图像不能搬到三维。高维阈值也不能直接猜为$1/2$。
\end{solution}

\originalchapter{8.334作业6：超出自旋波、长程畴壁与熔化}
\psource{Mehran Kardar，8.334 Spring 2014，PS6，截止2014年5月7日。原选做题3保留。}
\begin{exercise}[原题1(a)--(h)]\label{p:334-6-1}
单位$n$分量场$s$，$\beta H=\int\dd^dx\dd^dy\,K(|x-y|)s(x)\cdot s(y)$，核Fourier为$Kq^\omega/2$、$0<\omega<2$。(a)长距尾幂；(b)$s=(\pi,\sigma)$、$\sigma=\sqrt{1-\pi^2}$后有效作用；(c)二次Fourier与协方差；(d)消窄壳后$\avg\sigma$及单位长度恢复$\zeta$；(e)利用非解析核不被局部微扰修正求$K'$；(f)$T=1/K$流；(g)$d=\omega+\epsilon$的$T_c,\nu$；(h)小场$-h\cdot\int s$的流和$y_h$。
\end{exercise}
\begin{solution}
尺度$q^\omega$对应实空间$r^{-d-\omega}$尾；正$q^\omega$核的非接触部分有铁磁所需负号，并有对角常数让均匀态能0。不要把Fourier正系数误读为正的反铁磁实空间尾。

在北极坐标$\sigma=\sqrt{1-\pi^2}$，球面测度为$\prod_x\dd^{n-1}\pi/\sqrt{1-\pi^2}$。可写精确有效能为原核代入加$\tfrac12\sum_x\log(1-\pi_x^2)$。展开二次动能$H_0=K\int_q q^\omega\pi(q)\cdot\pi(-q)/2$，非局域四次项$K\int_q q^\omega|(\pi^2)(q)|^2/8$，另有测度顶点；这些顶点与对称约束一起保证不能任意产生Goldstone质量。二次协方差$\avg{\pi_i(q)\pi_j(q')}_0=\delta_{ij}(2\pi)^d\delta(q+q')/(Kq^\omega)$。

定义$I_\omega=\int_{\Lambda/b<q<\Lambda}\dd^dq/(2\pi)^dq^{-\omega}$，壳平均$\avg\sigma_>=1-(n-1)TI_\omega/2+\cdots=\zeta$。慢场$s'=s_</\zeta$，距离縮放给$K'=b^{d-\omega}\zeta^2K$，非解析动能系数无额外局部顶点修正（按题给假设）。窄壳$I_\omega=A_d\Lambda^{d-\omega}\dd\ell$，
\begin{equation}
\dot T=-(d-\omega)T+(n-1)A_d\Lambda^{d-\omega}T^2+\cdots.
\end{equation}
$n>1$、$\epsilon=d-\omega>0$小时$T_c=\epsilon/[(n-1)A_d\Lambda^\epsilon]+O(\epsilon^2)$，临界不稳定斜率$\epsilon$，$\nu=1/\epsilon+O(1)$。小场$h'=b^d\zeta h$，在固定点$y_h=d-\epsilon/2=(d+\omega)/2+O(\epsilon^2)$，对应$\eta=2-\omega$。短程模型有额外动能修正导致$n-2$，这里是$n-1$；不可把两公式混抄。$n=1$没有连续横向自旋波，不能除以$n-1$。
\end{solution}
\begin{exercise}[原题2(a)--(c)]\label{p:334-6-2}
XY在$d=2+\epsilon$的题给流$\dot T=-\epsilon T+4\pi^3y^2$、$\dot y=(2-\pi/T)y$。(a)有限温固定点；(b)最低非零阶本征值；(c)保持最低非平凡$\epsilon$阶估$d=3$超流$\nu,\alpha$。
\end{exercise}
\begin{solution}
非零$y$令$T_* =\pi/2$、$y_*^2=\epsilon/(8\pi^2)$。Jacobian为$\left(\begin{smallmatrix}-\epsilon&8\pi^3y_*\\4y_*/\pi&0\end{smallmatrix}\right)$，本征值$[-\epsilon\pm\sqrt{\epsilon^2+16\epsilon}]/2=\pm2\sqrt\epsilon+O(\epsilon)$。所以领先$\nu=1/(2\sqrt\epsilon)$；用$\alpha=2-d\nu$外推$\epsilon=1,d=3$给$\nu\simeq0.5,\alpha\simeq0.5$。

这是近二维截断流的粗外推，$\epsilon=1$没有小参数控制，不能与精确三维值等同；也不能先保留某个高阶本征值再丢其他同阶流修正，给出貌似更精确的数字。$\epsilon\to0$的幂次展开不还原二维BKT有限$\nu$，二维本来是本质奇点。
\end{solution}
\begin{exercise}[原题3(a)--(c)，选做]\label{p:334-6-3}
XY加$-\beta H_p=h_p\int\dd^2x\cos(p\theta)$，$p=1$磁场，$p=2,3,4,6$可来自离散晶格对称。(a)无涡旋sine--Gordon流；(b)含涡旋时证明$\dot h_p=[2-p^2/(4\pi K)]h_p$、$\dot K^{-1}=-\pi p^2h_p^2K^{-2}/2+4\pi^3y^2$、$\dot y=(2-\pi K)y$；(c)解释$8\pi/p^2<K^{-1}<\pi/2$区间、$p>4,p<4$相图与$p=4$边缘情况。
\end{exercise}
\begin{solution}
Gaussian壳角方差$\avg{\theta_>^2}=\dd\ell/(2\pi K)$，顶点平均$\avg{\cos[p(\theta_<+\theta_>)]}=\e^{-p^2\dd\ell/(4\pi K)}\cos p\theta_<$，体积因子2给场流。二阶连接累积量中两个近邻余弦取差角，$\cos[p(\theta(x)-\theta(y))]\simeq1-p^2[(x-y)\cdot\nabla\theta]^2/2$；角平均与截断壳积分按题给归一化给$\dot K=\pi p^2h_p^2/2$，即角被锁定使刚度增大。涡旋偶极则屏蔽刚度，贡献$\dot K^{-1}=4\pi^3y^2$，将两项合并即(b)。$h_p$归一化改变二阶非普适系数，不改顶点尺度维。

两种扰动都无关需$2-p^2/(4\pi K)<0$且$2-\pi K<0$，所以$8\pi/p^2<K^{-1}<\pi/2$，只有$p>4$非空。这是Gaussian\textbf{稳定区}，不是微扰方程只能在该区推导：区外小初始耦合仍能用流判定其相关，长到$O(1)$后才失控。原题“only valid”须以此条件解释。

$p>4$典型升温为离散锁定序、代数关联、涡旋无序两次转变；$p=2,3$无稳定中间Gaussian相，一般一次离散序转变；$p=1$外场去掉所有对称破缺，不保证有真转变。$p=4$二扰动可边缘竞争，关联到Ashkin--Teller类连续变指数，截断线性流不足以全定相图。
\begin{center}\begin{tikzpicture}[x=1.1cm,y=.9cm]\draw[->](0,0)--(9,0)node[right]{$T$};\node[left]at(0,0){$p>4$};\draw[thick](3,-.15)--(3,.15);\draw[thick](6,-.15)--(6,.15);\node[above]at(1.4,0){离散序};\node[above]at(4.5,0){代数序};\node[above]at(7.5,0){无序};\draw[->](0,-1.4)--(9,-1.4)node[right]{$T$};\node[left]at(0,-1.4){$p=2,3$};\draw[thick](4.5,-1.55)--(4.5,-1.25);\node[above]at(2.3,-1.4){离散序};\node[above]at(6.7,-1.4){无序};\end{tikzpicture}\end{center}
\end{solution}
\begin{exercise}[原题4(a)--(f)]\label{p:334-6-4}
一维标量$\pm1$自旋有逆平方铁磁相互作用，物理能$\beta H=-K\int\dd x\dd y\,s(x)s(y)/(2|x-y|^2)+\int\Phi(s)\dd x$，$\Phi$强制$\pm1$，例如$g(s^2-1)^2,g\gg1$。原卷把$\Phi$加在$-\beta H$中会反向失稳，本题取其所述“偏好$\pm1$”的稳定能量符号。(a)尺度$a$畴壁能；(b)单壁概率给$K_c$下界；(c)交替电荷畴壁对数气；(d)由$K_{\rm eff}=K-2Ky_0^2\int_a^\infty(r\dd r/a^2)(a/r)^{4K}+O(y_0^4)$得壳流；(e)画流及相；(f)推导此屏蔽式。原卷积分未显写$a^{-2}$，等价于所有长度按$a$计的约定。
\end{exercise}
\begin{solution}
单壁两侧相反的能代价$2K\int_0^L\dd u\int_0^L\dd v/(u+v)^2\simeq2K\log(L/a)$；位置熵$\log(L/a)$，单壁概率尺度$(L/a)^{1-2K}$，所以$K<1/2$壁大量出现，不可能保持该铁磁有序，$K_c\ge1/2$。

利用$1/(x-y)^2=\partial_x\partial_y\log|x-y|$并分部积分，$s'(x)=2\sum_iq_i\delta(x-x_i)$，得到$-\beta H_Q=4K\sum_{i<j}q_iq_j\log(|x_i-x_j|/a)+\log y_0\sum_i|q_i|$；电荷必须正负交替，局域壁核贡献逸度。

缩放$a\to ba$，中性$\sum_{i<j}q_iq_j=-N_w/2$给每壁$b^{-2K}$，位置积分$b$给$\dot y_0=(1-2K)y_0$。积分掉$a<r<ba$小壁对给$\dot K=-2Ky_0^2$。置$T=1/K$，$\dot T=2Ty_0^2$；$T<2$且小$y$时流向$y=0$有序固定线，$T>2$时壁相关流向高温无序，分界靠近$(T,y)=(2,0)$，具有类BKT本质尺度，强耦合处截断流不可靠。
\begin{center}\begin{tikzpicture}[x=2cm,y=2cm]\draw[->](0,0)--(2.5,0)node[right]{$T$};\draw[->](0,0)--(0,1.25)node[above]{$y_0$};\node[below]at(1.5,0){2};\draw[dashed](1.5,0)--(1.5,1);\draw[->,thick](.5,.65)..controls(.9,.2)..(1,.04);\draw[->,thick](1.3,.6)..controls(1.7,.75)..(2.2,1.1);\draw[->,thick](1.9,.25)..controls(2,.5)..(2.2,.7);\node at(.55,1){有序};\node at(1.9,1.25){无序};\end{tikzpicture}\end{center}

(f)固定外壁$+,-$于$0,L$，区间内允许小对方向$-,+$于$x,x+r$，相对权重$(y_0^2/a^2)(a/r)^{4K}$。与外壁交互的指数增量为$4K\log[(x+r)/x]+4K\log[(L-x)/(L-x-r)]\simeq4Kr[1/x+1/(L-x)]$。对$x$从核截断$a$积分到$L-a$给$8Kr\log(L/a)$；因此$\Delta\log W=8Ky_0^2\log(L/a)\int(r\dd r/a^2)(a/r)^{4K}$，与原$\log W=-4K\log(L/a)$比较正得题式。体积常数和端点核项吸收于自由能与逸度；交替约束允许一阶$r$贡献，是它比二维任意方向偶极屏蔽更简单的原因。
\end{solution}
\begin{exercise}[原题5(a)--(d)]\label{p:334-6-5}
各向同性弹性$\beta H=\tfrac12\int[2\mu u_{ij}u_{ij}+\lambda u_{ii}u_{jj}]\dd^dx$，$u_{ij}=(\partial_iu_j+\partial_ju_i)/2$。（原题将正稳定弹性能写在$-\beta H$之后，改正此符号。）(a)Fourier模；(b)协方差；(c)晶格截断$a$下$U^2(x)=\avg{|u(x)-u(0)|^2}$；(d)$\mu=\widehat\mu/(\kb T),\lambda=\widehat\lambda/(\kb T)$，由远距$U=c_La$的Lindemann准则求$T_m$及维数依赖。
\end{exercise}
\begin{solution}
纵横投影$P^L_{ij}=q_iq_j/q^2$、$P^T=I-P^L$，弹性核$q^2[\mu P^T+(\lambda+2\mu)P^L]$。静态本征值横$\mu q^2$（$d-1$支）、纵$(\lambda+2\mu)q^2$；若加质量密度动能，声速$\sqrt{\widehat\mu/\rho},\sqrt{(\widehat\lambda+2\widehat\mu)/\rho}$。稳定需$\mu>0$、体积模量$\lambda+2\mu/d>0$。

协方差为$(2\pi)^d\delta(q+q')q^{-2}[P^T/\mu+P^L/(\lambda+2\mu)]$。故
\begin{equation}
U^2(x)=2\int_q\frac{1-\cos(qx)}{q^2}\left[\frac{d-1}{\mu}+\frac1{\lambda+2\mu}\right].
\end{equation}
$d>2$远距为$2A_d\Lambda^{d-2}\kb T[(d-1)/\widehat\mu+1/(\widehat\lambda+2\widehat\mu)]/(d-2)$，$\Lambda\sim a^{-1}$。所以
\begin{equation}
T_m=\frac{c_L^2a^2(d-2)}{2\kb A_d\Lambda^{d-2}[(d-1)/\widehat\mu+1/(\widehat\lambda+2\widehat\mu)]}.
\end{equation}
这是截断与经验$c_L$决定的估计，不是精确转变温。$d=2$有$\log r$、$d<2$有$r^{2-d}$发散，此\textbf{绝对远距位移}准则会给所有正温$T_m=0$。实际二维可有代数晶序和有限温缺陷熔化，需第\ref{ch:melting}章的相对序与位错方法；不能以此准则错误宣称所有二维材料无固相。
\end{solution}

\originalchapter{8.334复习包1：Landau理论与Gaussian涨落}
\osource{Mehran Kardar，8.334 Spring 2014，Test1 Review及公开Solutions，11题。题4引用既有热统卷，其他全部明确小问在本章。}
\begin{exercise}[原题1(a)--(g)]\label{r:334-1-1}
简单立方二元合金$N_A,N_B$，六近邻，同种键$-J$、异种$+J$，$J>0$。(a)零温基态；(b)独立随机占位$p_A=N_A/N,p_B=N_B/N$的平均能；(c)混合熵；(d)$x=(N_A-N_B)/N$自由能四阶与失凸$T_c$；(e)$T\gtrless T_c$的势及spinodal$x_{sp}$；(f)全局分离组成$\pm x_{eq}$；(g)$T,x$相边界与亚稳边界。
\end{exercise}
\begin{solution}
固定两种总数时基态宏观分区，同种内部所有键有利，界面能为次广延修正。随机每键平均$-J(p_A-p_B)^2=-Jx^2$、3N键，$E=-3NJx^2$；$S=-N\kb[p_A\log p_A+p_B\log p_B]$，不是不同原子都能无关联移动的动力学证明。

$F/N=-\kb T\log2+(\kb T/2-3J)x^2+\kb Tx^4/12+O(x^6)$，$\kb T_c=6J$。近临界四阶模型$f''=\kb(T-T_c)+\kb Tx^2$，spinodal$x_{sp}^2=(T_c-T)/T$；对称共存极小$x_{eq}^2=3(T_c-T)/T$。上方单井、临界四次井、下方双井；$|x|<x_{sp}$局部不稳，$x_{sp}<|x|<x_{eq}$亚稳，总组成$|x|<x_{eq}$由两相杠杆规则实现。下图画近临界两线；低温应回到完整二项熵，四阶式不能使$|x|>1$仍合法。
\begin{center}\begin{tikzpicture}\begin{axis}[width=9cm,height=4.2cm,xlabel={$x$},ylabel={$T/T_c$},xmin=-.7,xmax=.7,ymin=.7,ymax=1.04,legend pos=south east]\addplot[thick,structurecolor,domain=-.7:.7]{3/(3+x*x)};\addlegendentry{共存}\addplot[dashed,domain=-.7:.7]{1/(1+x*x)};\addlegendentry{spinodal}\end{axis}\end{tikzpicture}\end{center}
\end{solution}
\begin{exercise}[原题2(a)--(e)]\label{r:334-1-2}
Ising$H=\tfrac12\sum_{ij}J_{ij}\sigma_i\sigma_j-h\sum\sigma_i$，等效近邻模型$J_{ij}=-J/N$。(a)$E(M,h)$；(b)以$M=Nm$分组$Z$、$F(m,h)$四阶；(c)鞍点适用条件、真假自由能凸性；(d)$h=0$磁化与$T_c$；(e)零场$C,\chi$奇异。
\end{exercise}
\begin{solution}
$E=-N(Jm^2/2+hm)$，每$M$简并$\binom N{N(1+m)/2}$，$Z=\sum_M\exp[-\beta F(m,h)]$、$F=E-\kb T\log\Omega$。Stirling给$f=F/N=-\kb T\log2+(\kb T-J)m^2/2+\kb Tm^4/12-hm$。大$N$指数$N$使$F_{thermo}/N=\min_m f+o(1)$，需极小附近展开且保持有限尺寸临界涨落的适当顺序。

原题说真实$F(h,T)$“凸”，在外场$h$上实际\textbf{凹}：$F_{hh}=-\beta\operatorname{Var}M\le0$；固定磁化的Legendre势才须凸并作双井凸包。$\min_mf$在$h=0,T<T_c$有尖点，不能把有限$N$解析配分直接叫非解析。

$\kb T_c=J$，$m=0$或$m^2=3(J-\kb T)/(\kb T)$（近临界）。奇异$f_{min}\simeq-3\kb(T_c-T)^2/(4T_c)$，$C$每站从上方0到下方$3\kb/2$跳变，忽略平滑背景。$\chi_m^+=1/[\kb(T-T_c)]$、$\chi_m^-=1/[2\kb(T_c-T)]$，$\gamma=1$、幅比2；总$M$响应再乘$N$。
\end{solution}
\begin{exercise}[原题3(a)--(c)]\label{r:334-1-3}
经典气$H=\sum p_i^2/(2m)+\tfrac12\sum V(r_i-r_j)$。(a)积分动量后的巨配分；(b)划分$N_L=V/a^3$小格、每格$n_\alpha=0,1$，积分用$a^3$和替代；(c)$n_\alpha=(1+\sigma_\alpha)/2,V_{\alpha\beta}=-J/N_L$映到上一题，推气体相变。
\end{exercise}
\begin{solution}
连续气体的巨正则配分函数为
\begin{equation}
\Xi=\sum_{N\ge0}\frac{(\e^{\beta\mu}/\lambda^3)^N}{N!}
\int\prod_i\dd^3r_i\exp\!\left[-\beta\sum_{i<j}V(r_{ij})\right].
\end{equation}
每格单占据使选位置的$N!$排列取消。记$\widetilde z=\e^{\beta\mu}a^3/\lambda^3$，则
\begin{equation}
\Xi\simeq\sum_{\{n_\alpha\}}\widetilde z^{\sum n_\alpha}
\exp\!\left[-\frac\beta2\sum_{\alpha\beta}V_{\alpha\beta}n_\alpha n_\beta\right].
\end{equation}

用$M=\sum\sigma$，$\sum n=(N_L+M)/2$，指数中$M^2$系数$\beta J/(8N_L)$、$M$系数$\beta J/4+\log\widetilde z/2$。所以等效Ising$J_I=J/4$、$h_I=J/4+\kb T\log\widetilde z/2$。临界$\kb T_c=J/4$、平均密度$n_c=1/(2a^3)$；$h_I=0$共存化学势$\mu=-J/2-\kb T\log(a^3/\lambda^3)$，低温两密度相。单占格与等效邻近耦合是近似，不表示所有连续势气体都精确同一临界温。
\end{solution}
\originalsection{原题4：表面活性剂凝聚}
原题4(a)--(f)的完整条件、排除面积、吸引积分$\int_a^\infty2\pi rV=-u_0$、张力、临界条件、两种热容，均与8.333 PS5题1相同，引用既有热统卷习题33.1；见上一编逐题接口，不复制第三次。
\begin{exercise}[原题5(a)--(c)]\label{r:334-1-5}
标量LG$\beta H=\int[K(\nabla m)^2/2+tm^2/2+cm^3+um^4]\dd^dx$，$K,c,u>0$。(a)势图显示正$t$的一阶跳变；(b)解共存$m,t$；(c)$K\xi^{-2}=f''(m_{eq})$，画相关长随$t$。
\end{exercise}
\begin{solution}
$c>0$偏好负$m$。$m=0$与负极小共存需$f(m_*)=f(0)$、$f'(m_*)=0$，相除消$m$得到$m_*=-c/(2u)$、$t_*=c^2/(2u)>0$。比它高采用0极小，比它低采用$m=[-3c-\sqrt{9c^2-16ut}]/(8u)$；另一根为不稳定或亚稳支，不能只用$f'=0$定相。

无序$\xi=\sqrt{K/t}$，有序$\xi=\sqrt{K/[t+6cm+12um^2]}$。在共存两侧曲率均$c^2/(2u)$，所以$\xi$有限，最大$\sqrt{2Ku}/c$，导数转折但本例$\xi$值无跳变。若$0$支继续到$t=0$，那是亚稳spinodal发散，不是真相边界的临界发散。势和相关长的精确两支公式已给出，图用$c=u=K=1$。
\begin{center}\begin{tikzpicture}\begin{axis}[width=9cm,height=4.2cm,xlabel={$t$},ylabel={$\xi$},xmin=-.2,xmax=1.4,ymin=0,ymax=1.6]\addplot[thick,structurecolor,domain=.5:1.4]{1/sqrt(x)};\addplot[thick,structurecolor,domain=-.2:.5]{1/sqrt(x+6*(-3-sqrt(9-16*x))/8+12*((-3-sqrt(9-16*x))/8)^2)};\draw[dashed](axis cs:.5,0)--(axis cs:.5,1.4142);\end{axis}\end{tikzpicture}\end{center}
\end{solution}
\begin{exercise}[原题6(a)--(d)]\label{r:334-1-6}
$\beta H=\int[K(\nabla m)^2/2+tm^2/2+um^4+vm^6-hm]\dd^dx$，$v>0$。(a)负$u$一阶；(b)共存$t$与跳变；(c)$u,t$相图；(d)$u=t=0$三临界$\beta,\delta,\gamma,\alpha$。
\end{exercise}
\begin{solution}
$u<0$且$h=0$，以$x=m^2$重复$f=tx/2+ux^2+vx^3$的等值、驻值，$x_*=-u/(2v)$、$t_*=u^2/(2v)$。$u>0$在$t=0$连续，两边界在原点相切；有序在其下，无序在其上，三临界需额外参数把$u$调到0。

$u=0$时$tm+6vm^5=h$，$h=0,t<0$的$m^4=-t/(6v)$给$\beta=1/4$，$t=0$给$\delta=5$，$\chi^{-1}=t+30vm^4$给$\gamma=1$。有序$f_{min}=-|t|^{3/2}/(3\sqrt{6v})$，$C=-f''\propto|t|^{-1/2}$，$\alpha=1/2$。这些都是鞍点指数，涨落是否小由本章末题判据而非四次固定点决定。
\end{solution}
\begin{exercise}[原题7(a)--(d)]\label{r:334-1-7}
旋转不变$n$分量$m$加$-h\cdot\int m$，$f(h)=-\log Z/V$仅依赖$h=|h|$。(a)$m_\alpha$；(b)响应张量；(c)纵横响应；(d)自发磁化时$h\to0$的横发散及纵是否必发散。
\end{exercise}
\begin{solution}
$\partial_{h_\alpha}f=f'(h)h_\alpha/h$，$m_\alpha=-f'h_\alpha/h$。再求导，
\begin{equation}
\chi_{\alpha\beta}=-f''\hat h_\alpha\hat h_\beta-\frac{f'}h(\delta_{\alpha\beta}-\hat h_\alpha\hat h_\beta),\quad \chi_\ell=-f'',\quad\chi_t=m/h.
\end{equation}
自发$m_0>0$使横$\chi_t\sim m_0/h$发散，反映无场整体转动的零代价；不能从旋转对称推出纵同样发散。均场纵有限，下一复习包将展示$d<4$的Goldstone两模涨落如何额外导致纵奇异。
\end{solution}
\begin{exercise}[原题8(a)--(h)]\label{r:334-1-8}
$d$维立方格XY$-\beta H=K\sum_{\langle ij\rangle}s_i\cdot s_j$，$s=(\cos\theta,\sin\theta)$。(a)角积分配分；(b)低温二次式；(c)周期1D$L$站正规模与实/复约束；(d)$d$维；(e)经典自由能、热容；(f)$\avg{s_0\cdot s_x}$远距；(g)横响应随$L$；(h)二维Gaussian模型的响应发散阈值。
\end{exercise}
\begin{solution}
$Z=\int\prod_i\dd\theta_i\exp[K\sum\cos(\theta_i-\theta_j)]$。低温$H_2=\tfrac K2\sum(\theta_i-\theta_j)^2+\mathrm{const}$。取$\theta_x=L^{-d/2}\sum_q\theta_q\e^{\ii qx}$，$q_\alpha=2\pi n_\alpha/L$、$\theta_{-q}=\theta_q^*$，$H_2=\tfrac12\sum_qD(q)|\theta_q|^2$，$D(q)=2K\sum_\alpha(1-\cos q_\alpha)$。每$q,-q$对为两个实模而非两个复自由模；$q=0$为整体角，固定或只积有限$2\pi$。

每独立非零实模贡献$-\log Z=\tfrac12\log[D/(2\pi)]$，物理每模能$\kb T/2$，故$C\simeq(N-1)\kb/2$；不在无质量经典模上添加量子零点能。Gaussian指数给$G(x)=\exp[-\tfrac12\avg{(\theta_x-\theta_0)^2}]$，方差$2\int_q(1-\cos qx)/(Kq^2)$。$d=1$有$G\sim\e^{-|x|/(2K)}$，$d=2$有$G\sim(a/r)^\eta,\eta=1/(2\pi K)$，$d>2$趋非零常数。

二维无长程序时$\chi_t\propto\int_a^Lr\dd r,r^{-\eta}$：$\eta<2$为$L^{2-\eta}$，等号$\log L$，$\eta>2$有限；形式阈值$K=1/(4\pi)$。它是只保自旋波的模型阈值，不是含涡旋真实XY的BKT刚度$2/\pi$。$d=1$响应在大$L$饱和$\sim\xi$；对于未选定方向的有限盒零场对称系综，$d>2$的$G\to m_0^2$给$\chi\sim m_0^2L^d$；若选定纯相并减去$m_0^2$，须区分整体旋转和局部Goldstone连接响应，后者无场横模积分$G_t(q)\sim1/q^2$给均匀场响应发散，不能用减去$m_0^2$后仍称全$L^d$连接磁化涨落。
\end{solution}
\begin{exercise}[原题9(a)--(f)]\label{r:334-1-9}
$d$维空间的平界面由$D=d-1$个坐标$x$上的高度$h(x)$，面积$A=\int\sqrt{1+(\nabla h)^2}\dd^Dx$、$H=\sigma A$。(a)二次作用和泛函配分；(b)Fourier模；(c)破缺对称；(d)高度差方差；(e)$d=4,3,2,1$；(f)小坡度近似条件。
\end{exercise}
\begin{solution}
$H=\sigma A_0+\tfrac\sigma2\int(\nabla h)^2+O((\nabla h)^4)$，$Z\propto\int Dh\exp[-\beta\sigma\int_q q^2|h_q|^2/2]$，固定平均高度移去$q=0$。破缺为沿界面法线的连续平移；位置改变没有体势代价，产生毛细Goldstone模。$\avg{|h_q|^2}=\kb T/(\sigma q^2)$，
\begin{equation}
\avg{[h(x)-h(0)]^2}=\frac{2\kb T}{\sigma}\int_{q<\Lambda}\frac{1-\cos qx}{q^2}\frac{\dd^Dq}{(2\pi)^D}.
\end{equation}
$d=4,D=3$远距趋截断依赖常数；$d=3,D=2$为$\kb T\log(r/a)/(\pi\sigma)$；$d=2,D=1$为$\kb T|r|/\sigma$；$d=1$界面是零维点，没有可比较的远距$x$或延展毛细波谱，其整体位置模需另固定，不能硬套$D$维长距公式。

局部坡度方差$\avg{(\nabla h)^2}\simeq(\kb T/\sigma)A_D\Lambda^D/D\ll1$，$\Lambda\sim1/a$。高度的对数粗糙不等于局部坡度失控：二维界面可局部光滑而长程高度不平，展开条件是坡度和UV尺度。
\end{solution}
\begin{exercise}[原题10(a)--(d)]\label{r:334-1-10}
三维超导LG$\beta H=\int[t|\Psi|^2/2+u|\Psi|^4+K|D_\mu\Psi|^2/2+L(\nabla\times A)^2/2]\dd^3x$，$D_\mu=\partial_\mu-\ii eA_\mu$。(a)局域规范变换；(b)均匀鞍点；(c)$t<0$以$\Psi=(R+\phi)\e^{\ii\theta},A=a,\nabla\cdot a=0$二阶能；(d)$\phi,\theta,a$的Fourier协方差。
\end{exercise}
\begin{solution}
$\Psi\to\e^{\ii\chi}\Psi,A\to A+\nabla\chi/e$时$D\Psi\to\e^{\ii\chi}D\Psi$，旋度不变。$R=0$（$t>0$）或$R^2=-t/(4u)$（$t<0$）、$A=0$。二次
\begin{equation}
H_2=\tfrac12\int[(K(\nabla\phi)^2-2t\phi^2)+KR^2(\nabla\theta-ea)^2+L(\nabla\times a)^2]\dd^3x.
\end{equation}
Coulomb横向$a$与梯度$\theta$的交叉积积分0，所以$\avg{\phi_q\phi_{-q}}=(Kq^2-2t)^{-1}$，$\avg{\theta_q\theta_{-q}}=(KR^2q^2)^{-1}$，$\avg{a_i(q)a_j(-q)}=P^T_{ij}/(Lq^2+Ke^2R^2)$，$\avg{|a|^2}=2/(Lq^2+Ke^2R^2)$。共同$(2\pi)^3\delta$或有限体积因子按Fourier约定补上。横光子获得质量$Ke^2R^2/L$；规范固定的相位坐标不是额外规范不变物理Goldstone粒子，真正可观测量用$\nabla\theta-eA$。
\end{solution}
\begin{exercise}[原题11(a)--(g)]\label{r:334-1-11}
$n$分量$\beta H=\int[K(\nabla m)^2/2+tm^2/2+u(m^2)^2+v(m^2)^3]\dd^dx$，三临界$u=0,v>0$。(a)鞍点热容；(b)纵横Hessian；(c)关联；(d)Gaussian自由能；(e)热容修正；(f)Ginzburg与上临界维；(g)一般$u_{2p}(m^2)^p$多临界点。
\end{exercise}
\begin{solution}
负侧$m_0^4=-t/(6v)$，$f_{sp}=-|t|^{3/2}/(3\sqrt{6v})$，$C_{sp}/V=-f_{sp}''=|t|^{-1/2}/(4\sqrt{6v})$。质量$r_\ell=t+30vm_0^4=-4t$、$r_t=t+6vm_0^4=0$；正侧均为$t$。两点$G_\ell(q)=1/(Kq^2+r_\ell)$、$G_t=1/(Kq^2+r_t)$。

$\delta f=\tfrac12\int_q[\log(Kq^2+r_\ell)+(n-1)\log(Kq^2+r_t)]$，每个独立实模一个半迹。负側$r_\ell'=-4$、横质量恒0，$\delta C/V=8\int_q(Kq^2-4t)^{-2}\sim|t|^{(d-4)/2}$。相对鞍点$\delta C/C_{sp}\sim\sqrt v,|t|^{(d-3)/2}$，所以$d<3$失控，$d=3$边缘对数，三临界$d_u=3$。

一般Gaussian场维$(d-2)/2$使$u_{2p}$的尺度维$d-p(d-2)=2p-(p-1)d$，$d_u=2p/(p-1)$。应调低阶偶次耦合至相应多临界流形，不能保留一个相关四次项却仍称由六次项控制的三临界点。
\end{solution}

\originalchapter{8.334复习包2：标度、Goldstone奇异与各向异性RG}
\osource{Mehran Kardar，8.334 Spring 2014，Test2 Review及公开Solutions，9题。保留题设截断流与现实模型、单一场缩放与复合算符的区别。}
\begin{exercise}[原题1(a)--(d)]\label{r:334-2-1}
流体奇异Helmholtz每粒子$f_s=t^{2-\alpha}g(\delta\rho/t^\beta)$，$t=|T-T_c|/T_c$、$\delta\rho=\rho-\rho_c$。沿临界等容线$Q\sim t^x$、临界等温线$Q\sim|\delta\rho|^y$，求(a)能与熵；(b)$C_V,C_P$；(c)$\kappa_T,\alpha_P$，验证$C_P-C_V=TV\alpha_P^2/\kappa_T$；(d)共存潜热。先按题给\textbf{只保奇异自由能}的模型，说明缩放函数导数不消失的假设。
\end{exercise}
\begin{solution}
对$t$每微分降低一次幂，对密度每微分降低$\beta$次幂；一个齐次量沿$t=0$指数为$x/\beta$。$S_s=-f_{s,T}$、$U_s=f_s+TS_s$领先$t^{1-\alpha}$；$P_s=\rho^2f_{s,\rho}\sim t^{2-\alpha-\beta}$，$P_{s,T}\sim t^{1-\alpha-\beta}$、$P_{s,\rho}\sim t^{2-\alpha-2\beta}$。于是纯奇异模型给
\begin{center}\begin{tabular}{lll}\toprule 量&$x$&$y$\\\midrule $U_s/N,S_s/N$&$1-\alpha$&$(1-\alpha)/\beta$\\ $C_V,C_P$&$-\alpha$&$-\alpha/\beta$\\ $\kappa_T$&$\alpha+2\beta-2$&$(\alpha+2\beta-2)/\beta$\\ $\alpha_P$&$\beta-1$&$(\beta-1)/\beta$\\\bottomrule\end{tabular}\end{center}
$\alpha_P=P_T/(\rho P_\rho)$，$C_P-C_V$指数$2(\beta-1)-(\alpha+2\beta-2)=-\alpha$，吻合。共存$\delta\rho_\pm\sim t^\beta$使$L_s=T(S_{s,+}-S_{s,-})\sim t^{1-\alpha}$，前提两侧缩放函数差不为0；完全对称$g$可使该领先差消失。

不能把这个表不加条件用到任意实际流体：若完整状态方程有通常非零解析$P_T$背景，$\kappa_T\sim t^{-\gamma}$、$\gamma=2-\alpha-2\beta$仍成立，而$\alpha_P\sim\kappa_T$、$C_P-C_V\sim\kappa_T$，不再是表中的纯奇异幂。类似CC若共存线斜率有限非零、体积差$t^\beta$，完整潜热可领先$t^\beta$。官方表解是在忽略这些解析背景且相关函数系数非零的题设模型内，缺省前提必须保留。
\end{solution}
\begin{exercise}[原题2(a)--(g)]\label{r:334-2-2}
题给$d=1+\epsilon$ Ising流$\dot T=-\epsilon T+T^2/2$、$\dot h=dh$。(a)流与固定点；(b)$y_t,y_h$；(c)由$\xi$变换得$\nu,\Delta$；(d)奇异自由能、$\alpha$；(e)$\beta,\gamma$；(f)由关联积分得$\eta$；(g)$d=1$沿$h=0,T\to0$相关长。
\end{exercise}
\begin{solution}
$h=0$轴有$T=0$、$T_c=2\epsilon$；$0<T<T_c$向0，$T>T_c$向高温，$h$以$\e^{d\ell}$离开轴。临界$y_t=\partial_T\dot T|_{T_c}=\epsilon$、$y_h=d$。$\xi(t,h)=b\xi(b^{y_t}t,b^{y_h}h)$选$b=|t|^{-1/y_t}$，得$\nu=1/\epsilon$、$\Delta=d/\epsilon$。

体积标度$f_s=b^{-d}f_s(b^{y_t}t,b^{y_h}h)=|t|^{d/\epsilon}\Phi(h/|t|^{d/\epsilon})$，所以$\alpha=2-d/\epsilon$、$\beta=(d-y_h)/y_t=0$、$\gamma=(2y_h-d)/y_t=d/\epsilon$。$\chi\sim\int^{\xi}r^{d-1}r^{-d+2-\eta}\dd r\sim\xi^{2-\eta}$，给$\eta=2-d=1-\epsilon$。

$\beta=0$是\textbf{题给截断流}忽略磁场异常重标度的结果，不能冒称所有真实$d>1$ Ising的精确磁化指数。$\epsilon=0$积分$1/T(\ell)=1/T_0-\ell/2$，到$O(1)$尺度$\xi\sim a\exp(2/T_0)$；一维严格转移矩阵$\xi=-1/\log\tanh(1/T)$回收此指数式。
\end{solution}
\begin{exercise}[原题3(a)--(h)]\label{r:334-2-3}
$n$分量有序LG$H=\int[K(\nabla m)^2/2+tm^2/2+u(m^2)^2]$、$t<0$。(a)置$m=(m_0+\phi_\ell)e_\ell+\pi$展开全部项；(b)将$U=4um_0\int\phi_\ell\pi^2$Fourier化；(c)裸纵横协方差；(d)横四点Wick（$n-1$分量）；(e),(f)二阶纵关联连接展开；(g)略去横断连项给$\chi_\ell(q)$；(h)小$q$奇异。说明有序相存在与$n$的条件。
\end{exercise}
\begin{solution}
$m_0^2=-t/(4u)$。除常数，二次$H_0=\tfrac12\int[K((\nabla\phi_\ell)^2+(\nabla\pi)^2)-2t\phi_\ell^2]$，三次$4um_0\phi_\ell(\phi_\ell^2+\pi^2)$，四次$u(\phi_\ell^2+\pi^2)^2$。题取最低纵横耦合$U=4um_0\sum_a\int_{qk}\phi_\ell(q)\pi_a(k)\pi_a(-q-k)$，其他顶点并非永不存在，只为该IR奇异分离忽略。

$G_\ell^0=(Kq^2-2t)^{-1}$、$G_t^0=(Kq^2)^{-1}$。对两复合$\pi^2$的Wick，断连自收缩系数$(n-1)^2$只移鞍点，两个跨顶点收缩各系数$n-1$，合为$2(n-1)G_tG_t$。纵关联展开$G_\ell=G_\ell^0+\tfrac12[\avg{\phi\phi U^2}_0-G_\ell^0\avg{U^2}_0]+\cdots$，一次项因纵奇数0；两个外腿接不同顶点有2种，消$1/2$，得到
\begin{equation}
\chi_\ell(q)=G_\ell^0(q)+32u^2m_0^2(n-1)[G_\ell^0(q)]^2\int_k\frac1{Kk^2K|k+q|^2}+\cdots.
\end{equation}
$2<d<4$时积分$\propto K^{-2}|q|^{d-4}$，$d=4$对数，故纵$\chi$也奇异。另一检查是长度近固定约束$\phi_\ell\simeq-\pi^2/(2m_0)$，给系数$(n-1)/(2m_0^2)$乘同一泡图，与上式$G_\ell^0(0)=(-2t)^{-1}$一致。

$n=1$没有横模，系数0，不推出Ising有序纵响应这样发散；短程连续对称$d\le2$非零温没有这里假定的普通长程序。$q$趋0时有限阶微扰也会失控，泡图识别机制而非提供所有尺度的均匀收敛展开。
\end{solution}
\begin{exercise}[原题4(a)--(j)]\label{r:334-2-4}
LG加$r m_1^2/2+v m_1^2m^2$、$u>0$，整题令$v=0$。(a)给定$|m|$的易方向；(b)平行/垂直1轴鞍点自由能；(c)$t,r$相图与阶；(d)Goldstone；(e)正$t,r,u=0$的1、2分量协方差；(f)四次Fourier；(g),(h)首阶关联修正；(i)$\chi_{22}^{-1}$及$t_c$；(j)$r$在O(n)点相关性和普适类。
\end{exercise}
\begin{solution}
$r>0$易垂直平面，$r<0$易1轴。平行密度$f_\parallel=-(t+r)^2/(16u)$（$t+r<0$），垂直$f_\perp=-t^2/(16u)$（$t<0$）。无序需$t>0,t+r>0$；$r>0$在$t=0$连续入O$(n-1)$平面序，$r<0$在$t+r=0$连续入Ising轴序。$r=0,t<0$两方向相会，$\partial_rf=m_1^2/2$跳变，为取向一阶线，端点$(0,0)$为双临界。平面序有$n-2$ Goldstone、轴序0、$r=0$的各向同性序有$n-1$。

$D_1=Kq^2+t+r,D_\perp=Kq^2+t$，$G_i=D_i^{-1}$。$U=u\sum_{ij}\int_{q_1q_2q_3}m_i(q_1)m_i(q_2)m_j(q_3)m_j(-\sum q)$。Wick外腿同分量两种附加收缩，令$I_1=\int_qG_1,I_\perp=\int_qG_\perp$，质量修正
\begin{equation}
\delta r_1=4u[3I_1+(n-1)I_\perp],\qquad\delta r_\perp=4u[I_1+(n+1)I_\perp].
\end{equation}
关联首修正$G_i-\delta r_iG_i^2$，$\chi_{22}^{-1}=t+4u[I_1+(n+1)I_\perp]+O(u^2)$。平面支$r>0$的一阶$t_c=-4u\int_q[(Kq^2+r)^{-1}+(n+1)(Kq^2)^{-1}]+O(u^2)$，相同截断，发散时不可当有限移位。

$r$分裂二次质量，O(n)固定点处相关，微小易轴/易平面最终分别流向Ising/O$(n-1)$类，不仅改临界温。$r=0$经细调才留O(n)；若平面为O(2)且有效空间二维，应另考虑BKT而非通常有序相均场。
\end{solution}
\begin{exercise}[原题5(a)--(d)]\label{r:334-2-5}
$n$分量LG四次$u(m^2)^2+v\sum_im_i^4$。(a)易方向及简并；(b)$u,v$稳定区；(c)有正六次稳定项、$u>0$的$t,v$鞍点相图；(d)Goldstone。
\end{exercise}
\begin{solution}
固定$R^2=\sum m_i^2$，$R^4/n\le\sum m_i^4\le R^4$。$v>0$取各分量等模$\pm R/\sqrt n$，共$2^n$条有向对角方向；$v<0$取单一分量$\pm R$，共$2n$条有向轴（不计正负时各减半）。四次稳定$u+v/n>0$（$v\ge0$）或$u+v>0$（$v<0$）。

$u>0$、$v>-u$时$t=0$连续，$v>0$入对角序、$-u<v<0$入轴序；$v=0,t<0$取向切换，恢复O(n)连续简并。$v<-u$轴方向有效四次$a_4=u+v<0$，加$wR^6,w>0$后$t_*=a_4^2/(2w)$、$R_*^2=-a_4/(2w)$为一阶，$(t,v)=(0,-u)$三临界。非零$v$离散易方向无Goldstone，$v=0$有$n-1$；稳定区边界上若没有六次项原四次模型无有限极小，不能把那里画成合法低温相。
\end{solution}
\begin{exercise}[原题6(a)--(d)]\label{r:334-2-6}
证明题给一环流$\dot u=\epsilon u-4C[(n+8)u^2+6uv]$、$\dot v=\epsilon v-4C[12uv+9v^2]$，$\epsilon=4-d$、$C\simeq A_4/K^2$。(b)全部固定点、$n<4,n>4$稳定流；(c)温度流及稳定点$\nu$；(d)涨落对稳定区与小负$v,n>4$相变的影响。
\end{exercise}
\begin{solution}
没有$4!$约定的四腿顶点为$8u(\delta_{ab}\delta_{cd}+\delta_{ac}\delta_{bd}+\delta_{ad}\delta_{bc})+24v\delta_{abcd}$。二阶连接累积量两顶点以两条壳传播子连接，三种外腿配对；壳泡$I_2=C\dd\ell$。投影到$aabb$与$aaaa$不变量，$u$修正来自$4(n+8)Cu^2+24Cuv$，$v$修正来自$48Cuv+36Cv^2$；两个纯$v$顶点不能生成不同分量$aabb$，纯$u$必须保持O(n)所以不生成独立$v$，这两条是计数检查。再加树层维$\epsilon$即题式。
\begin{center}\begin{tikzpicture}[x=1cm,y=1cm]\fill(-1,0)circle(2pt);\fill(1,0)circle(2pt);\draw(-1,0)..controls(-.4,.8)and(.4,.8)..(1,0);\draw(-1,0)..controls(-.4,-.8)and(.4,-.8)..(1,0);\draw(-2,.6)--(-1,0)--(-2,-.6);\draw(2,.6)--(1,0)--(2,-.6);\node[above]at(0,.65){两条壳传播子};\node[left]at(-1.7,.6){$a$};\node[left]at(-1.7,-.6){$b$};\node[right]at(1.7,.6){$c$};\node[right]at(1.7,-.6){$d$};\end{tikzpicture}\end{center}
用$\bar u=4Cu,\bar v=4Cv$，固定点为Gaussian$(0,0)$、独立Ising$(0,\epsilon/9)$、O(n)$(\epsilon/(n+8),0)$、cubic$(\epsilon/(3n),\epsilon(n-4)/(9n))$。Gaussian两方向$+\epsilon$，Ising有$+\epsilon/3$与$-\epsilon$，均不完全稳定。O(n)本征值$-\epsilon,\epsilon(n-4)/(n+8)$，cubic本征值$-\epsilon,\epsilon(4-n)/(3n)$；所以$n<4$稳定O(n)、$n>4$稳定正$v$cubic，$n=4$一环边缘需更高阶。

平移临界质量后$\dot t=[2-(n+2)\bar u-3\bar v]t+\cdots$，来自单壳tadpole$4[(n+2)u+3v]I_1$及$I_1$对$t$的导数。因此
\begin{equation}
\nu_{O(n)}=\tfrac12+\frac{n+2}{4(n+8)}\epsilon+O(\epsilon^2),\quad
\nu_{cubic}=\tfrac12+\frac{n-1}{6n}\epsilon+O(\epsilon^2).
\end{equation}
$n>4$负$v$可离开原四次稳定区、流向大耦合，六次稳定项允许涨落诱发一阶。原均场稳定区因涨落的长尺度要求而缩小；runaway是机制指示，有限阶beta函数不独自证明精确一阶边界，$n=4$阈值也仅是一环结论。
\end{solution}
\begin{exercise}[原题7(a)--(f)]\label{r:334-2-7}
平方格Ising，$b=2$的Niemeijer--van Leeuwen首累积量。(a)四站方块内/跨键划分；(b)块自旋多数规则、2:2平票各一半，列所有微观权重；(c)$\avg U_0$与$K'$；(d)固定点、$y_t$；(e)小$h$的$h',y_h$；(f)精确比较。
\end{exercise}
\begin{solution}
块内四边作$H_0$，相邻块之间两条原键作$U$。配置类及完整数目如下：
\begin{center}\begin{tabular}{llll}\toprule 块内$M$&构型数&键和、权重&给$S'=+$比例\\\midrule 4&1&$4,\e^{4K}$&1\\ 2&4&$0,1$&1\\ 0（相邻两正）&4&$0,1$&$1/2$\\ 0（交替）&2&$-4,\e^{-4K}$&$1/2$\\ $-2$&4&$0,1$&0\\ $-4$&1&$4,\e^{4K}$&0\\\bottomrule\end{tabular}\end{center}
反号给$S'=-$，两条件权重各$Z_b=\e^{4K}+6+\e^{-4K}$。旋转对称使每边站$\avg{\sigma_i}_{S'}=a(K)S'$、$a=(\e^{4K}+2)/Z_b$。首累积量跨块因子化，两条边给$K'=2Ka^2$；新场$h'=4ah$，常数每块$\log Z_b$。

非零固定点$a=1/\sqrt2$。$x=\e^{4K_*}$满足$(\sqrt2-1)x^2-(6-2\sqrt2)x-1=0$，$K_*\simeq0.5187$。$y_t=\log[\partial_K(2Ka^2)|_{K_*}]/\log2\simeq1.006$，$y_h=\log(4/\sqrt2)/\log2=1.5$。精确$K_c=\tfrac12\log(1+\sqrt2)=0.440687$、$y_t=1,y_h=15/8$。热指数接近是该近似的数值巧合，磁指数误差明显；首累积量丢跨块相关不能误称精确RG。随卷代码穷举16态及验证全部表项。
\end{solution}
\begin{exercise}[原题8(a)--(d)]\label{r:334-2-8}
$d$维Potts近邻$K\delta_{s_i,s_j}$。(a)一维$b=2$精确隔站流及固定点；(b)采用\textbf{先移键、后消站}的MK流；(c)$d>1$有非零有限固定点；(d)$d=2$、$q=3,1,0$的$K_*,y_t$，$q=1,0$作随机簇解析延拓。
\end{exercise}
\begin{solution}
一维$F_q(K)=\log[(\e^{2K}+q-1)/(2\e^K+q-2)]$，常数$g'=2g+\log(2\e^K+q-2)$，$K=0$稳定、$\infty$不稳定（$K'=K-\log2$）。先把$2^{d-1}$条键并到一条，再精确消站，$K'=F_q(2^{d-1}K)$。与先消站后并键的另一MK次序不同，本文用官方题的次序；两种都是近似，数字不宜混比。

$q>0$小$K$有$K'\simeq2^{2d-2}K^2/q<K$，大$K$有$K'\simeq2^{d-1}K-\log2>K$（$d>1$），连续性保证中间有不稳定有限根。$d=2,x=\e^{K_*}$，方程$x(2x^2+q-2)=x^4+q-1$。

$q=3$正根$x=2$，$K_*=\log2$、热本征值$\lambda_t=16/9$、$y_t=\log_2(16/9)\simeq0.8301$。$q=1$根$x=(1+\sqrt5)/2=\varphi$，$K_*\simeq0.48121$、$\lambda_t=4(1-\varphi^{-1})$、$y_t\simeq0.6115$。$q\to0$化简$K'=\log[(1+\e^{2K})/2]=K+K^2/2+\cdots$，仅有限根0，$y_t=0$但二阶边缘相关。$q=0$不能当“零种自旋”的普通统计和，它是连通图生成函数的极限；$q=1$对应渗流。原$q>0$的非零固定点论证在奇异$q=0$极限须重新检查。
\end{solution}
\originalsection{原题9：Potts转移矩阵的完整引用}
原题9(a)--(c)周期Potts链的矩阵、最大与次大根、自由能、$\xi(T\to0)$，逐项在习题\ref{p:334-4-5}(a),(b)及其解答。原卷Hamilton式把末键单列又让前和上限写$N$，本卷按每条周期近邻边只计一次处理，不将末键双计。

\originalchapter{8.334复习包3：图展开、矩阵场与粗糙化}
\osource{Mehran Kardar，8.334 Spring 2014，Test3 Review及公开Solutions，12题。修复粗糙化能量符号与正交矩阵说明，并区分官方统一场缩放和复合算符壳平均。}
\begin{exercise}[原题1(a),(b)]\label{r:334-3-1}
蜂窝格O(n)单位向量，改用键权重$\exp f(s_i\cdot s_j)=1+nt\,s_i\cdot s_j$。(a)高温$t$图展开；(b)$n\to0$如何表示单条自避免聚合链，说明使用的生成函数。
\end{exercise}
\begin{solution}
每条边选1或$nt\sum_as_i^as_j^a$。规范化球面积分$\int s^a=0$、$\int s^as^b=\delta_{ab}/n$；蜂窝度3，所以有边的顶点必须度2，不能度1或3。所选图是互不交叉闭环，长度$l$的边$n^lt^l$被$l$个顶点$n^{-l}$抵消，绕圈分量迹剩$n$，故
\begin{equation}
Z=\sum_{\rm disjoint\ loops}t^{N_b}n^{N_{loop}}.
\end{equation}
裸$\lim_{n\to0}Z=1$只剩空图，不能直接说它等于一个聚合物。$\partial_n Z|_0$选一个闭合自避免环；插入$\avg{s_i\cdot s_j}$则两个端点允许一条开放线，其余闭环在$n\to0$消失，得到从$i$到$j$的单条自避免路径生成函数。度3格点避免偶四度顶点的分叉/交叉歧义；模型正权重需相应$t$范围，解析$n$极限作为图生成工具。
\end{solution}
\begin{exercise}[原题2(a)--(f)]\label{r:334-3-2}
$q$态Potts键$\e^{K\delta}=C(K)[1+T(K)g(s,s')]$、$g=q\delta-1$。(a)$C,T$；(b)证明$\sum_sg=0$、$\sum_sg(s_1,s)g(s,s_2)=qg(s_1,s_2)$、$\sum_{ss'}g^2=q^2(q-1)$；(c)一维自由能、$\avg{g(s_m,s_n)}$；(d)平方格高温至$T^4$与低温首项；(e)对偶、唯一边界前提下$K_c$；(f)$q=2$图的区别。
\end{exercise}
\begin{solution}
比较相同/不同自旋给$C=(\e^K+q-1)/q$、$T=(\e^K-1)/(\e^K+q-1)$。矩阵$g=qI-\mathbf1\mathbf1^T$在均匀向量上0、其余$q-1$维本征值$q$，故$g^2=qg$，三恒等式立即由行和和迹得到。

一维$Z=q^NC^{N-1}$（开链），环链改为$Z=q^NC^N[1+(q-1)T^N]$，等价于转移矩阵精确式；体自由能$\log Z/N=\log(qC)$，关联$\avg{g(s_m,s_n)}=(q-1)T^{|m-n|}$。平方格最小非空高温闭图为一个面，权重$(q-1)T^4$，故$Z=q^NC^{2N}[1+N(q-1)T^4+O(T^6)]$（thermodynamic bulk，有限小环面另有绕圈图）。低温基态$q$个，单站改成另$q-1$色破4键，$Z=q\e^{2NK}[1+N(q-1)\e^{-4K}+\cdots]$。

匹配$T(K)=\e^{-K^*}$即$(\e^K-1)(\e^{K^*}-1)=q$，唯一铁磁边界在$K_c=\log(1+\sqrt q)$。$q=2$时$g=\sigma_i\sigma_j$，每顶点奇度平均0，图只有偶度；$q>2$可有三度及更复杂分叉图，不能用同一简单Ising费米环结构完成求和。这是所用解法的结构区别，不是所有非Ising模型一律不可解的定理。
\end{solution}
\begin{exercise}[原题3(a)--(c)]\label{r:334-3-3}
Potts$\e^{K\delta}=1+v\delta,v=\e^K-1$，加场$h\sum_i\delta_{s_i,1}$。(a)证明随机簇式$Z=\sum_Gv^{N_b(G)}\prod_{c\in G}[q-1+\e^{hN_s(c)}]$；(b)$q\to1$独立键渗流与平方阈值；(c)$q\to0$单连通簇生成及lattice animals解释。
\end{exercise}
\begin{solution}
选中$v\delta$的每键强迫两端同色，一个连通簇所有自旋同色；若色1权$\e^{hN_s}$，其余$q-1$色权1，独立簇相乘即式。$h=0$时$Z=\sum_Gv^{N_b}q^{N_c}$，所有图均允许，不用删除带端点的图。

$q=1$有$Z=(1+v)^{N_E}$，规范化图概率$p^{N_b}(1-p)^{N_E-N_b}$，$p=v/(1+v)$，平方格$p_c=1/2$。$q\to0$时裸$Z\to0$，其$q$的领先系数只保一个簇，即\textbf{固定全顶点集}上的连通生成图。若要枚举任意尺寸的单根lattice animals，需再标根或对参与顶点数引入生成变量；不能把有限图的$Z/q$未经说明当任意尺寸簇数。它允许分叉和环，与上一题$n\to0$的单条自避免链不同。
\end{solution}
\begin{exercise}[原题4(a),(b)]\label{r:334-3-4}
平方格Potts。(a)比较高低温或从随机簇证明配分函数对偶、$K_c$；(b)从对偶求$U=\avg{\beta H}=-\partial_{\log K}\log Z$的临界值，说明不连续转变时的限制。
\end{exercise}
\begin{solution}
平面图的已占边与双图未占边对应，Euler关系将簇数换成双簇数，bulk每站$\phi(K)=2\log v-\log q+\phi(K^*)$、$v=\e^K-1,vv^*=q$。有限边界有解析常数和扭转扇区，见习题\ref{p:334-5-3}。自对偶$K_c=\log(1+\sqrt q)$。

微分$\partial_KK^*=-q\e^K/[v(v+q)]$，故$U(K)/N=-2K\e^K/v+(K/K^*)(\partial_KK^*)U(K^*)/N$。自对偶导数$-1$，若能连续，$U_c/N=-K_c(1+1/\sqrt q)$。$q>4$的一阶情况下，它严格固定\textbf{两侧能量平均}，而非使潜热0：$[U_c^-+U_c^+]/(2N)=-K_c(1+1/\sqrt q)$。原题对任意$q$的单值说法须加连续前提；这也是只用自对偶不能决定相变阶的例子。
\end{solution}
\begin{exercise}[原题5(a)--(c)]\label{r:334-3-5}
平方格所有从原点出发随机游走，权$t_x^{l_x}t_y^{l_y}$、$t_x,t_y\ge0$。(a)终点$(x,y)$总权$W$及$t=(t_x+t_y)/2<1/4$；(b)近临界远距等权曲线；(c)全部终点加权平均步数如何发散。
\end{exercise}
\begin{solution}
单步Fourier振幅为$A(q)=2t_x\cos q_x+2t_y\cos q_y$。所有步数作几何和，$\widetilde W=1/(1-A)$，所以
\begin{equation}
W(x,y)=\int_{[-\pi,\pi]^2}\frac{\dd^2q}{(2\pi)^2}
\frac{\e^{\ii(q_xx+q_yy)}}{1-2t_x\cos q_x-2t_y\cos q_y}.
\end{equation}
非负权总和有限需$2(t_x+t_y)<1$，即$t<1/4$。

近临界$\delta=1-4t$，小$q$核$\delta+t_xq_x^2+t_yq_y^2$，逆Fourier$W\simeq K_0(\sqrt\delta\sqrt{x^2/t_x+y^2/t_y})/(2\pi\sqrt{t_xt_y})$。等值是椭圆，$\xi_x=\sqrt{t_x/\delta},\xi_y=\sqrt{t_y/\delta}$；某耦合严格0时降为一维，不能使用两维Jacobian。总生成$\sum_{xy}W=1/(1-4t)$，共同放大两步权的对数导数给$\avg l=4t/(1-4t)$，$\gamma_l=1$。固定终点的条件步数需对该$W(x,y)$另求导，不与全终点平均混为一个量。
\end{solution}
\begin{exercise}[原题6(a)--(c)]\label{r:334-3-6}
平方格各向异性Ising$-\beta H=\sum[K_x\sigma_{x,y}\sigma_{x+1,y}+K_y\sigma_{x,y}\sigma_{x,y+1}]$。(a)改造行转移矩阵/Jordan--Wigner正规模，给体自由能；(b)由奇点求相边界并验$K_x=K_y^*$；(c)奇异自由能及指数、幅比受各向异性影响。
\end{exercise}
\begin{solution}
行内键仍给$\exp[K_x\sum\sigma^z_j\sigma^z_{j+1}]$；行间单站核对角化后将原$K^*$换为$K_y^*$，$\tanh K_y^*=\e^{-2K_y}$，常数每站$(2\sinh2K_y)^{1/2}$。转为Majorana后每$(q,-q)$二模块的双曲转移乘积得到
\begin{equation}
\cosh\gamma(q)=\cosh2K_x\cosh2K_y^*-\sinh2K_x\sinh2K_y^*\cos q,
\end{equation}
其两根$\e^{\pm\gamma}$，选最大根求迹，$\phi=\log Z/N=\tfrac12\log(2\sinh2K_y)+(2\pi)^{-1}\int_0^\pi\gamma(q)\dd q$。这展示哪些各向同性步骤改变，而非直接猜一个替换式。

最小间隙$\gamma(0)=2|K_x-K_y^*|$，所以$\sinh2K_x\sinh2K_y=1$。近临界$\mu=2(K_x-K_y^*)$、$v_q^2=\sinh2K_x\sinh2K_y^*$，小$q$有$\gamma\simeq\sqrt{\mu^2+v_q^2q^2}$，积分奇异$\phi_s=\mu^2\log(C/|\mu|)/(4\pi v_q)$加解析项。$\mu\propto T-T_c$给$\nu=1$与对数热容$\alpha=0$；有限正空间刚度改变尺度和非普适振幅，不改临界类。两侧$\mu^2\log|\mu|$系数同，热容对数幅比1；其他普适幅比比较须使用同一坐标和场归一化，相关长方向之比本身是非普适度量。
\end{solution}
\begin{exercise}[原题7(a)--(d)]\label{r:334-3-7}
二维Ising界面沿$x$，忽略孤岛、悬垂，整数列高$h_n$，周期列边界。(a)证明$\beta H=2K_yL+2K_x\sum|h_{n+1}-h_n|$；(b)列转移核及对角化；(c)界面自由能；(d)零界面代价与原精确临界线比较。
\end{exercise}
\begin{solution}
每水平投影列破一竖键，代价$2K_y$；高度差$|\Delta h|$破这么多横键，每条$2K_x$。核$T(h,h')=\e^{-2K_y-2K_x|h-h'|}$、本征函数$\e^{\ii kh}$，$r=\e^{-2K_x}$，$\lambda(k)=\e^{-2K_y}(1-r^2)/(1-2r\cos k+r^2)$。整体高度无限平移量需固定或除去，不影响每长度自由能。

最大$k=0$，$\beta\sigma_{axis}=-\log\lambda_0=2K_y-\log\coth K_x=2K_y+\log\tanh K_x$。为0时$\e^{-2K_y}=\tanh K_x$，正好$\sinh2K_x\sinh2K_y=1$。这个无悬垂近似所得边界与精确线相合，但不由此证明所有界面构型都已计入；形式负代价区原有序界面解释已失效。
\end{solution}
\begin{exercise}[原题8(a)--(c)]\label{r:334-3-8}
$n$分量磁性LG考虑空间各向异性。(a)一般低阶作用；(b)这种各向异性是否相关；(c)题设La2CuO4经典三分量spin，平面内强平方格反铁磁、相邻平面弱同向耦合，低温态与转变普适类。
\end{exercise}
\begin{solution}
一般$H=\int[t m^2/2+\tfrac12K_{\alpha\beta}\partial_\alpha m\cdot\partial_\beta m+u(m^2)^2-h\cdot m]$，对称正定空间刚度可正交对角化，再以$x_\alpha'=x_\alpha/\sqrt{K_\alpha}$化成等向梯度，体积Jacobian与场单位重定。有限正定刚度的各向异性属度量参数，通常不改变普适指数；若有特征值0、不同$q^2/q^4$阶或长程尾，则不能这样移去，见作业3。

低温每平面$s_{x,y,z}=(-1)^{x+y}\hat n$，弱同向层耦合使同一位置跨层取同向，共用Néel轴。任意有限正层耦合在足够长尺度构成三维O(3)类（在题设经典、无spin各向异性假设内）；很弱时有广的二维到三维 crossover，孤立短程二维连续spin平面本身无非零温普通长程序。不能把空间弱耦合与spin易轴相混后改猜Ising。
\end{solution}
\begin{exercise}[原题9(a)--(c)]\label{r:334-3-9}
单位$n$分量$H=(2T)^{-1}\int[(\nabla s)^2+gs_1^2]$，$d=2+\epsilon$。题给$\dot T=-\epsilon T+(n-2)T^2/(2\pi)$，整体spin缩放$\zeta=1-(n-1)T\dd\ell/(4\pi)$。(a)固定$T,y_T$；(b)各向异性$g$的流、指数；(c)$T,g$相图、$g\to0$的边界形状。区分官方仅以$\zeta^2$处理平方场的近似和直接壳平均。
\end{exercise}
\begin{solution}
$n>2$、$\epsilon>0$，$T_*=2\pi\epsilon/(n-2)$、$y_T=\epsilon$。官方解把$s_1^2$仅视为两个整体场因子，给$y_g^{off}=d+2\log\zeta/\dd\ell=2-\epsilon/(n-2)$。但平方场是复合算符，壳内分量混合不能只用$\avg s$的平方。

直接取慢单位轴$u$和快横向$\pi_a$，$s=u\sqrt{1-\pi^2}+\sum e_a\pi_a$、$\avg{\pi_a\pi_b}=TI\delta_{ab}$，$I=\dd\ell/(2\pi)$。由$\sum_ae_{a1}^2=1-u_1^2$得$\avg{s_1^2}_>=(1-nTI)u_1^2+TI$。常数项移自由能，真正各向异性算符$s_1^2-1/n$系数为$1-nTI$。在固定$T_*$，$g/(2T)$的体积缩放给
\begin{equation}
y_g=d-nT_*/(2\pi)=2-\frac{2\epsilon}{n-2}+O(\epsilon^2).
\end{equation}
这比官方仅场缩放多一个复合算符顶点修正；两者的$2/\epsilon$领先 crossover形状相同，但首阶更正不同，不能冒称官方数字就是全一环算符结果。

$g<0$易轴Ising序，$g>0$易平面O$(n-1)$序，$g=0$轴以下有取向一阶线，上面无序；各连续边界近双临界满足$g\sim|\delta T|^\phi$、$\phi=y_g/y_T=2/\epsilon-2/(n-2)+O(\epsilon)$（官方近似为$2/\epsilon-1/(n-2)$）。这是小$\epsilon$结果，不能把外推$\epsilon=1$所得小或零$y_g$当实际三维结论。
\end{solution}
\begin{exercise}[原题10(a)--(j)]\label{r:334-3-10}
实$N\times N$正交矩阵场$M$，$H=K\int\Tr[\nabla M\cdot\nabla M^T]/4$、$MM^T=M^TM=I$。(a)元素式和基态；(b)对称$\sigma=(M+M^T)/2$、反对称$\pi=(M-M^T)/2$的作用及约束；(c)近$I$到四阶，略测度项；(d)Goldstone数；(e)Fourier协方差；(f)壳$\avg{\Tr\sigma}$与$\zeta$；(g)两种直接梯度修正给$\widetilde K=K+NI/2$；(h)$K'$；(i)$T$流与$d\lessgtr2$、$T_c,\nu$；(j)场$-h\int\Tr M$与$y_h$。$I=\int_{\Lambda/b}^\Lambda\dd^dq/[(2\pi)^dq^2]$。
\end{exercise}
\begin{solution}
元素作用$K\sum_{ij}(\nabla M_{ij})^2/4$，约束$\sum_kM_{ik}M_{jk}=\delta_{ij}$，任意空间常数正交矩阵为基态。$M=\sigma+\pi$给$H=K\int\Tr[(\nabla\sigma)^2-(\nabla\pi)^2]/4$，\textbf{两边}正交式还给$[\sigma,\pi]=0$及$\sigma^2-\pi^2=I$。近$I$正分支$\sigma=\sqrt{I+\pi^2}=I+\pi^2/2+O(\pi^4)=I-\pi\pi^T/2+\cdots$。

代入得到
\begin{equation}
H=-\frac K4\int\Tr(\nabla\pi)^2+\frac K8\int\Tr[(\pi\nabla\pi)^2+\pi^2(\nabla\pi)^2]+O(\pi^6).
\end{equation}
两个四阶项因矩阵不对易不能合成标量单项。$\pi_{ij}=-\pi_{ji}$有$N(N-1)/2$个独立实模，$H_0=K\sum_{i<j}\int(\nabla\pi_{ij})^2/2$，所以$\avg{\pi_{ij}(q)\pi_{kl}(q')}=(2\pi)^d\delta(q+q')(\delta_{ik}\delta_{jl}-\delta_{il}\delta_{jk})/(Kq^2)$。

$\avg{\pi^2}_>=-[(N-1)I/K]I_N$，故$\avg{\Tr\sigma}_>=N[1-(N-1)I/(2K)]$、$\zeta=1-(N-1)I/(2K)$为\textbf{平均}长度/迹的恢复因子。官方旁注说正交矩阵本征值全$\pm1$、近$I$的迹恒$N$，这是错误的：二维转动本征$\e^{\pm\ii\theta}$、迹$2\cos\theta$即反例。本推导不使用该旁注，也不将每个涨落矩阵的迹强行固定$N$。

(g)写慢梯度$A=\nabla\pi_<$、快矩阵$B=\pi_>$，两直接梯度收缩分别$\avg{\Tr(BABA)}=-(I/K)\Tr A^2$、$\avg{\Tr(B^2A^2)}=-(N-1)(I/K)\Tr A^2$，来自上述两Kronecker项逐索引求和。四阶系数$K/8$合成$-NI\Tr A^2/8$，与二次$-K\Tr A^2/4$比较即$\widetilde K=K+NI/2$。快速梯度收缩只移局域/测度相关项，不在此直接动能投影中重复计数。

重标$K'=b^{d-2}\zeta^2\widetilde K=b^{d-2}[K-(N-2)I/2]+\cdots$。$I=A_d\Lambda^{d-2}\dd\ell$，$\dot T=-(d-2)T+(N-2)A_d\Lambda^{d-2}T^2/2$。$N>2$时$d<2$及$d=2$向高温，二维$\xi\sim\exp[4\pi/((N-2)T)]$；$d=2+\epsilon$有$T_c=4\pi\epsilon/(N-2)$、$\nu=1/\epsilon$。$N=2$为角变量，微扰二阶0，还需拓扑缺陷。

$h'=b^d\zeta h$，临界$y_h=d-(N-1)\epsilon/(N-2)=2-\epsilon/(N-2)$。SO(3)局部双覆盖S$^3$关联O(4)场模型的微扰几何，但全局SO(3)有$\mathbb Z_2$基本群、S$^3$没有，拓扑缺陷不在这个一环展开内；$M$场也对应覆盖变量的复合算符。不能从局部微扰等价就无条件宣布实际堆叠三角磁体全属于O(4)普适类。
\end{solution}
\begin{exercise}[原题11(a)--(g)]\label{r:334-3-11}
三维空间中的二维界面，稳定Gaussian$H_0=K\int(\nabla h)^2/2$，晶格锁定$-U=y\int\cos(2\pi h)/a^2$；原题$H_0$负号改为正以保证配分收敛。(a)中性顶点平均，$C(x)=\log(|x|/a)/(2\pi)$；(b)$G_k(x-y)=\avg{\e^{\ii k(h(x)-h(y))}}$的二次导数；(c)锁定的首个非零修正；(d)有效刚度与微扰失效阈值；(e)壳RG；(f)相图；(g)长距粗糙度跳变的普适意义。
\end{exercise}
\begin{solution}
平均高度零模积分令总顶点荷$\sum q_\alpha=0$，否则平均0。Gaussian特征函数为$\exp[-\tfrac12\sum q_\alpha q_\beta G(x_\alpha-x_\beta)]$；中性消去自项后为$\exp[K^{-1}\sum_{\alpha<\beta}q_\alpha q_\beta C(x_\alpha-x_\beta)]$。故$G_k(r)=\e^{-k^2C(r)/K}$，$-\partial_k^2G_k|_0=\avg{[h(r)-h(0)]^2}=2C(r)/K$。

令$p=2\pi$，$D=C(r-x)-C(r-y)-C(x)+C(y)$，首次锁定项因非中性0，二阶连接归一化给
\begin{equation}
\frac{\delta G_k}{G_k^0}=\frac{y^2}{4a^4}\int\dd^2x\dd^2y\,(a/|x-y|)^{p^2/(2\pi K)}[\cosh(kpD/K)-1]+O(y^4).
\end{equation}
减1正是分母配分的断连扣除。为区分完整Green函数与一次RG壳，写$C(r)=C_<(r)+C_>(r)$，其中
\begin{equation}
C_>(r)=\int_{\Lambda\e^{-\dd\ell}<|q|<\Lambda}\frac{\dd^2q}{(2\pi)^2}
\frac{1-\cos(q\cdot r)}{q^2}
=\frac{\dd\ell}{2\pi}[1-J_0(\Lambda r)]+O((\dd\ell)^2).
\end{equation}
它用于重整顶点$y'=b^2\e^{-p^2\langle h_>^2\rangle/2}y$；计算长距响应中的中心积分则要保留低动量Green函数。小偶极$s=x-y$有$D\simeq s\cdot F(R)$，$F(R)=\nabla C(r-R)-\nabla C(R)$。对$R$做分部积分（边界以相同UV/IR截断处理），用$\nabla^2C(R)=\delta(R)$，或等价地用Parseval，可逐步得
\begin{equation}
\int\dd^2R\,|F(R)|^2
=2\int\frac{\dd^2q}{(2\pi)^2}\frac{1-\cos(q\cdot r)}{q^2}
=2C(r).
\end{equation}
相对位移的角积分$\int\dd\varphi(s\cdot F)^2=\pi s^2|F|^2$，代回上式的$y^2/(4a^4)$并与$G_k=\exp[-k^2C/K_{\rm eff}]$对照，得到
\begin{equation}
\delta K=\frac{\pi p^2}{4}y^2\int_a^\infty\frac{s^3\dd s}{a^4}
\left(\frac as\right)^{p^2/(2\pi K)}.
\end{equation}
该积分在$K\ge\pi/2$（$p=2\pi$）IR发散，指示锁定相关。它不是把全部$s$一口气改成RG增量；只有$a<s<a\e^{\dd\ell}$一壳给$\dot K=\pi p^2y^2/4$。

以此显写的字面$y/a^2$及圆形相对位移截断，$\dot y=(2-\pi/K)y$、$\dot K=\pi^3y^2$。作业6原题把二阶流写为$\dot K=\pi p^2h_p^2/2$；对同一个$p=2\pi$余弦顶点，本文圆壳的字面系数$y$与该题二阶规范的振幅关系为$y=\sqrt2h_p$。两种逸度的线性尺度维相同，改变二阶规范不改$K_c$；不可在一个方程里混用$y$的字面定义和$h_p$的二阶系数。$K_R<\pi/2$且逸度流0为粗糙对数相；较大刚度锁到整数高度，为平相；边界类BKT，有限裸逸度使裸临界线偏离单一$\pi/2$。

粗糙侧临界重整化刚度$K_R=\pi/2$，长距方差领先$[2/\pi^2]\log(r/a)$；平相长距饱和。因此普适跃变是
\begin{equation}
\lim_{r\to\infty}\frac{\avg{\Delta h^2}}{\log(r/a)}:
\quad\frac{2}{\pi^2}\longrightarrow0.
\end{equation}
等价地，领先差为$2\log(r/a)/\pi^2$。平相饱和值和有限$r$取极限方式非普适，不能声称高度方差本身有一个与距离无关的普适绝对跳变量。
\end{solution}
\begin{exercise}[原题12(a)--(c)]\label{r:334-3-12}
平方格整数高RSOS，每邻站$\Delta h=0$能0，$\Delta h=\pm1$能$K/2$，更大禁止。(a)改键$n_{ij}=h_i-h_j$，加每面和0约束$\delta_{n,0}=\int_0^{2\pi}\e^{\ii n\theta}\dd\theta/(2\pi)$，求双格相互作用；(b)大$K$双模型与XY及前题一致性；(c)一维高度链（二维空间界面）是否有有限温粗糙化转变。
\end{exercise}
\begin{solution}
键重$w(n)=1$（0）、$t=\e^{-K/2}$（$\pm1$）、0（其他）。面约束Fourier变量$\theta_p$给每键$\sum_{n=-1}^1w(n)\e^{\ii n(\theta_p-\theta_{p'})}=1+2t\cos(\theta_p-\theta_{p'})$，所以双键能$\widetilde v(\Delta\theta)=-\log[1+2t\cos\Delta\theta]$，边界和整体高度因子需固定，环面另加绕数约束。

大$K$小$t$时$\widetilde v=-2t\cos\Delta\theta+O(t^2)$，XY双刚度约$2t$、处于高温无序，对应原高度代价大、平相。小$K$时该Fourier权可能负，不可把整个形式$-\log$未经检验当全温正概率XY模型；对偶图和拓扑流需更完整处理。此限制不影响大$K$题问。

一维可直接独立伸长$n_i=-1,0,1$，规范化权$1,t,t$，$\operatorname{Var}n_i=2t/(1+2t)>0$。固定$h_0$后$\avg{(h_j-h_0)^2}=j\,2t/(1+2t)$，任意有限$K$长链都粗糙，只有$K\to\infty$完全平；周期闭链有限尺寸改为桥方差$\propto j(1-j/N)$，不产生有限温转变。
\end{solution}

\originalchapter{来源、全课程覆盖与证明范围}
\source{MIT OpenCourseWare公开课历、讲义、作业与复习包；逐件URL、SHA-256、字节数与页数在独立源码清单。冻结范围不因题目困难而删减。}
\originalsection{课历到中文编章的映射}
8.333的L1--L4热力学、L5--L6概率见第\ref{ch:thermo}章及作业1、2；L7--L8相空间与BBGKY见第\ref{ch:bbgky}章；L9H定理见第\ref{ch:boltzmann}章；L10--L11守恒与流体见第\ref{ch:hydro}章；L12--L14经典系综见第\ref{ch:ensemble}章；L15--L19累积、团簇、virial、变分与临界见第\ref{ch:clusters}章，并由既有作业5接口连接。L20分子振转、声子与黑体、L21--L23量子态/全同粒子系综由第\ref{ch:ensemble}章、相关作业和既有转子接口展开；L24Fermi、L25Bose、L26超流见第\ref{ch:qfluid}章及量子气体作业。该完整主线不是对原件每页的逐字翻译。

8.334的L1--L3集体场、LG与Goldstone见第\ref{ch:landau}章；L4--L5关联、Gaussian与Ginzburg见第\ref{ch:gaussian}章；L6--L8齐次、RG与Gaussian标度见第\ref{ch:scaling}章；L9--L12微扰及$\epsilon$展开见第\ref{ch:wilson}章；L13--L14实空间RG见第\ref{ch:realrg}章；L15--L17高低温图、对偶、phantom loops见第\ref{ch:duality}章；L18--L19二维Ising自由能与临界见第\ref{ch:ising2}章；L20--L22非线性$\sigma$、XY缺陷与Coulomb流见第\ref{ch:bkt}章；L23弹性与熔化见第\ref{ch:melting}章；L24--L25Brownian、平衡/守恒场动力学见第\ref{ch:dynamics}章；L26尺度不变、开放系统、生长见第\ref{ch:growth}章。第\ref{ch:random}章的无序与聚合物补充用于大纲/图展开的连接，不虚称另有未公开的26讲之后讲义。
\originalsection{题目计数与年份边界}
两门课12套作业按原卷顶层编号分别56与29；复习包8.333为9、7、10及历史期末21，8.334为11、9、12，合计164个\emph{原题位置}。本卷新写137个习题环境承载138个原位置（两维/三维Dirac题合编一环境），10个原位置完整回指卷内同题，16个原位置回指既有热统中文成果（作业5的12题及后续4次重复引用）。这些计数不能说成164道不同新增题。

8.333前三复习包原题页为Fall2007；期末包标2007复习、收1998、1999、2000、2003--2006历史考试。8.334三包为Spring2014复习题。2013/2014当期实际考试题面与学生最终项目未取得，不补造；课历列考试不等于该考试文件已公开。作业5的既有中文题面与解、固定PDF哈希和原编号全部核对，正常构建不依赖外卷文件。

课程讲义、作业和复习/官方解共78份PDF（52份讲义、12套作业、14份复习/官方解）；另有20份已取得的文字视频参考PDF，合计98份、1226个源物理页。各课有26条视频链接，但仅下载了各10份文字参考，不声称全52段视频逐分钟转录，也不把文字视频参考计入题集。源文件页数包含题面与答案重叠，不能当独立新增中文页。来源文件本地核查，公开只保链接、页数、字节数和哈希，不重分发第三方例外图片。
\originalsection{数学与物理复核的限度}
题面全部明确小问保留，包括选做；8.333作业6代表表He3能量系数原件漏印，未用它反造数据，所求四小问照完整条件解答。侧壁压力的能零、自由链模数、Lorentz气热流的peculiar速度、Coulomb点源符号、正交矩阵谱、相变凸性与各向异性复合算符等确定可核查问题单列修正。

RG一环、MK和块近似的指数是对应近似的结果；非解析流或runaway不单独证明全相图。一般二维Ising磁化定理、完整非平衡KPZ固定点与淬火无序非微扰理论未在本册从零证明，引用课程结果或明确补充假设。复习题中小$\epsilon$外推、Bose热容左右极限及定义、锁定刚度归一化均保留条件。有限矩阵、穷举和短Monte Carlo用于验证给定计算，不能作为无限体系相变证明。
\begin{thebibliography}{9}
\bibitem{333} Mehran Kardar，MIT OpenCourseWare，8.333 Statistical Mechanics I: Statistical Mechanics of Particles，Fall2013，讲义来自较早未指定年份版本。\url{https://ocw.mit.edu/courses/8-333-statistical-mechanics-i-statistical-mechanics-of-particles-fall-2013/}
\bibitem{334} Mehran Kardar，MIT OpenCourseWare，8.334 Statistical Mechanics II: Statistical Physics of Fields，Spring2014，讲义来自较早未指定年份版本。\url{https://ocw.mit.edu/courses/8-334-statistical-mechanics-ii-statistical-physics-of-fields-spring-2014/}
\bibitem{license} MIT OCW使用条款与逐件例外。\url{https://ocw.mit.edu/pages/privacy-and-terms-of-use/}；中文改编CC BY-NC-SA4.0。\url{https://creativecommons.org/licenses/by-nc-sa/4.0/}
\end{thebibliography}

\end{document}
